\input{ko_preamble}

% OK, start here.
%
\begin{document}

\title{변형 이론}


\maketitle

\phantomsection
\label{section-phantom}

\tableofcontents

\section{서론}
\label{section-introduction}

\noindent
이 장의 목표는 가군과 사상 등의 변형 이론을 비교적 쉽게 소개하는 것이다.
여기서는 소박한 코탄젠트 복합체를 사용하여 증명할 수 있는 결과들을 다룬다.
코탄젠트 복합체를 다루는 다음 장에서는 이 결과들을 조금 더 확장한다.
더 깊이 공부하려는 독자는 이 주제에 관한 Illusie의 논저
\cite{cotangent}를 참고할 수 있다.





\section{환의 변형과 소박한 코탄젠트 복합체}
\label{section-deformations}

\noindent
이 절에서는 소박한 코탄젠트 복합체를 사용하여 약간의 변형 이론을 전개한다.
먼저 핵이 제곱이 영인 아이디얼 $I$인 전사 환 준동형 $A' \to A$를 잡는다.
또한 환 준동형 $A \to B$, $B$-가군 $N$, 그리고 $A$-가군 준동형
$c : I \to N$이 주어졌다고 하자. 다음 도표의 물음표 자리를 채우는 대상을
찾을 수 있는지가 이 절의 첫 번째 문제이다.
\begin{equation}
\label{equation-to-solve}
\vcenter{
\xymatrix{
0 \ar[r] & N \ar[r] & {?} \ar[r] & B \ar[r] & 0 \\
0 \ar[r] & I \ar[u]^c \ar[r] & A' \ar[u] \ar[r] & A \ar[u] \ar[r] & 0
}
}
\end{equation}
또 그러한 해가 존재할 때 얼마나 유일한지도 묻는다. 더 정확히 말하면,
핵이 제곱이 영인 아이디얼이며 $N$과 동일시되고, $A' \to B'$가 주어진
준동형 $c$를 유도하는 $A'$-대수의 전사 준동형 $B' \to B$를 찾는다.
이러한 $B'$를 (\ref{equation-to-solve})의 {\it 해}라고 부른다.

\begin{lemma}
\label{lemma-huge-diagram}
다음과 같은 가환 도표가 주어졌다고 하자.
$$
\xymatrix{
& 0 \ar[r] & N_2 \ar[r] & B'_2 \ar[r] & B_2 \ar[r] & 0 \\
& 0 \ar[r]|\hole & I_2 \ar[u]_{c_2} \ar[r] &
A'_2 \ar[u] \ar[r]|\hole & A_2 \ar[u] \ar[r] & 0 \\
0 \ar[r] & N_1 \ar[ruu] \ar[r] & B'_1 \ar[r] & B_1 \ar[ruu] \ar[r] & 0 \\
0 \ar[r] & I_1 \ar[ruu]|\hole \ar[u]^{c_1} \ar[r] &
A'_1 \ar[ruu]|\hole \ar[u] \ar[r] & A_1 \ar[ruu]|\hole \ar[u] \ar[r] & 0
}
$$
앞쪽과 뒤쪽이 (\ref{equation-to-solve})의 해이면 다음이 성립한다.
\begin{enumerate}
\item 다음 군에는 표준적인 원소가 존재한다.
$\Ext^1_{B_1}(\NL_{B_1/A_1}, N_2)$
이 원소가 영이라는 것은 도표를 가환시키는 환 준동형
$B'_1 \to B'_2$가 존재하기 위한 필요충분조건이다.
\item 도표를 가환시키는 준동형 $B'_1 \to B'_2$가 존재하면,
그러한 모든 준동형의 집합은
$\Hom_{B_1}(\Omega_{B_1/A_1}, N_2)$ 아래의 주동차 공간이다.
\end{enumerate}
\end{lemma}

\begin{proof}
$E=B_1$을 집합으로 보자. 핵이 $J$인 전사 준동형
$A_1[E] \to B_1$을 생각하자. 이는 다음 식으로 소박한 코탄젠트
복합체를 정의할 때 사용하는 전사 준동형이다.
$$
\NL_{B_1/A_1} = (J/J^2 \to \Omega_{A_1[E]/A_1} \otimes_{A_1[E]} B_1)
$$
가환대수학, 절 \ref{algebra-section-netherlander}를 보라.
$\Omega_{A_1[E]/A_1} \otimes B_1$은 자유 $B_1$-가군이므로
$$
\Ext^1_{B_1}(\NL_{B_1/A_1}, N_2) =
\frac{\Hom_{B_1}(J/J^2, N_2)}
{\Hom_{B_1}(\Omega_{A_1[E]/A_1} \otimes B_1, N_2)}
$$
이다. 오른쪽 가군 안에 장애 원소를 구성하자.
$J'=\Ker(A'_1[E]\to B_1)$로 놓는다. 핵이 $I_1A'_1[E]$인
전사 준동형 $J'\to J$가 있음에 유의하자. 각 $e\in E$에 대응하는
변수를 $x_e\in A_1[E]$라 쓰자. $B_1$에서 $x_e$의 상을 들어 올린
$y_e\in B'_1$과, $B_2$에서 $x_e$의 상을 들어 올린
$z_e\in B'_2$를 택한다. 이 선택들은 다음 $A'_1$-대수 준동형들을 정한다.
$$
A'_1[E] \to B'_1 \quad\text{그리고}\quad A'_1[E] \to B'_2
$$
첫 번째 준동형은 $J'\to N_1$, $f'\mapsto f'(y_e)$를 주고,
두 번째 준동형은 $J'\to N_2$, $f'\mapsto f'(z_e)$를 준다.
계산하면 이 두 준동형은 $(J')^2$을 영으로 보낸다. 도표에서
$c_1$과 $c_2$가 들어 있는 왼쪽 정사각형이 가환하므로, $N_2$로 가는
준동형으로 보았을 때 둘은 $I_1A'_1[E]$ 위에서 일치한다.
$B'_1$은 $J'\to A'_1[E]$와 $J'\to N_1$의 밀어내기임에 유의하자.
따라서 $J'\to N_1\to N_2$와 $J'\to N_2$가 일치하면 도표를
가환시키는 준동형 $B'_1\to B'_2$를 얻는다. 이제 다음 준동형의
동치류를 장애 원소로 삼는다.
$$
J/J^2 \to N_2,\quad f \mapsto f'(z_e) - \nu(f'(y_e))
$$
여기서 $\nu:N_1\to N_2$는 주어진 준동형이고 $f'\in J'$는 $f$의
한 올림이다. 위의 논의로부터 이는 잘 정의된다. 어떤
$\delta_{i,e}\in N_i$에 대하여 $z_e$를 $z_e+\delta_{2,e}$로,
$y_e$를 $y_e+\delta_{1,e}$로 바꿀 자유가 있다. 그러면 위 준동형은
다음과 같이 바뀐다.
$$
f \mapsto f'(z_e + \delta_{2, e}) - \nu(f'(y_e + \delta_{1, e})) =
f'(z_e) - \nu(f'(z_e)) +
\sum (\delta_{2, e} - \nu(\delta_{1, e}))\frac{\partial f}{\partial x_e}
$$
이는 정확히 $J/J^2\to N_2$를 합성 준동형
$J/J^2\to\Omega_{A_1[E]/A_1}\otimes B_1\to N_2$만큼 바꾸는 것이다.
여기서 두 번째 준동형은 $\text{d}x_e$를
$\delta_{2,e}-\nu(\delta_{1,e})$로 보낸다. 따라서 장애 원소는
잘 정의되며, 그 값이 영인 것과 올림이 존재하는 것은 동치이다.

\medskip\noindent
(2)는 다음 관찰에서 따른다. 도표를 가환시키는 두 준동형
$\varphi,\psi:B'_1\to B'_2$가 주어지면 $\varphi-\psi$는
$A_1$-미분 $D:B_1\to N_2$를 거쳐 분해된다.
\begin{align*}
D(fg) & = \varphi(f'g') - \psi(f'g') \\
& =
\varphi(f')\varphi(g') - \psi(f')\psi(g') \\
& =
(\varphi(f') - \psi(f'))\varphi(g') + \psi(f')(\varphi(g') - \psi(g')) \\
& =
gD(f) + fD(g)
\end{align*}
따라서 $D$는 유일한 $B_1$-선형 준동형
$\Omega_{B_1/A_1}\to N_2$에 대응한다. 역으로 그러한 선형 준동형에서
미분 $D$를 얻는다. 도표를 가환시키는 환 준동형
$\psi:B'_1\to B'_2$가 하나 주어지면 $\psi+D$ 역시 도표를
가환시키는 또 다른 환 준동형이다.
\end{proof}

\begin{lemma}
\label{lemma-choices}
(\ref{equation-to-solve})의 해가 존재하면, 해의 동형류들의 집합은
$\Ext^1_B(\NL_{B/A}, N)$ 아래의 주동차 공간이다.
\end{lemma}

\begin{proof}
(\ref{equation-to-solve})의 두 해 $B'_1$, $B'_2$가 주어지면
보조정리 \ref{lemma-huge-diagram}에 의해 준동형 $B'_1\to B'_2$의
존재에 대한 장애 원소
$o(B'_1,B'_2)\in\Ext^1_B(\NL_{B/A},N)$을 얻는다. 이 원소는
분명히 동형의 존재에 대한 장애이므로 동형류들을 구별한다. 따라서 증명을
끝내기 위해서는 해 $B'$와 원소 $\xi\in\Ext^1_B(\NL_{B/A},N)$가
주어졌을 때
$o(B',B'_\xi)=\xi$를 만족하는 두 번째 해 $B'_\xi$를 찾을 수 있음을
보이면 충분하다.

\medskip\noindent
$E=B$를 집합으로 보자. 핵이 $J$인 전사 준동형 $A[E]\to B$를
생각하자. 이는 다음 식으로 소박한 코탄젠트 복합체를 정의할 때 쓰인다.
$$
\NL_{B/A} = (J/J^2 \to \Omega_{A[E]/A} \otimes_{A[E]} B)
$$
가환대수학, 절 \ref{algebra-section-netherlander}를 보라.
$\Omega_{A[E]/A}\otimes B$는 자유 $B$-가군이므로
$$
\Ext^1_B(\NL_{B/A}, N) =
\frac{\Hom_B(J/J^2, N)}
{\Hom_B(\Omega_{A[E]/A} \otimes B, N)}
$$
따라서 $\xi$를 준동형 $\delta:J/J^2\to N$의 동치류로 나타낼 수 있다.

\medskip\noindent
각 $e\in E$에 대응하는 변수를 $x_e\in A[E]$라 쓰자. $B$에서
$x_e$의 상을 들어 올린 $y_e\in B'$를 택한다. 이 선택들은
$A'$-대수 준동형 $\varphi:A'[E]\to B'$를 정한다.
$J'=\Ker(A'[E]\to B)$로 놓는다. $\varphi$는 준동형
$\varphi|_{J'}:J'\to N$을 유도하며, 다음 도표와 같이 $B'$는
밀어내기이다.
$$
\xymatrix{
0 \ar[r] & N \ar[r] & B' \ar[r] & B \ar[r] & 0 \\
0 \ar[r] & J' \ar[u]^{\varphi|_{J'}} \ar[r] & A'[E] \ar[u] \ar[r] &
B \ar[u]_{=} \ar[r] & 0
}
$$
$\psi:J'\to N$을 $\varphi|_{J'}$와 다음 합성 준동형의 합으로 둔다.
$$
J' \to J'/(J')^2 \to J/J^2 \xrightarrow{\delta} N.
$$
그러면 $\psi$에 따른 밀어내기는 위와 같은 도표를 이루는 또 하나의
환 확대 $B'_\xi$이다. 계산하면 원하는 대로
$o(B',B'_\xi)=\xi$임을 알 수 있다.
\end{proof}

\begin{lemma}
\label{lemma-extensions-of-algebras}
$A$를 환, $B$를 $A$-대수, $N$을 $B$-가군이라 하자. 다음과 같은
$A$-대수 확대들의 동형류 집합을 생각하자.
$$
0 \to N \to B' \to B \to 0
$$
여기서 $N$은 제곱이 영인 아이디얼이다. 이 집합은
$\Ext^1_B(\NL_{B/A},N)$과 표준적으로 전단사 대응한다.
\end{lemma}

\begin{proof}
이를 보이기 위해 (\ref{equation-to-solve})이 다음 도표로 주어지는 경우에
앞의 결과들을 적용한다.
$$
\xymatrix{
0 \ar[r] & N \ar[r] & {?} \ar[r] & B \ar[r] & 0 \\
0 \ar[r] & 0 \ar[u] \ar[r] & A \ar[u] \ar[r]^{\text{id}} & A \ar[u] \ar[r] & 0
}
$$
그러면 이 보조정리는 보조정리 \ref{lemma-choices}와 해 $N\oplus B$가
존재한다는 사실에서 따른다. 전단사 대응의 직접적인 구성은 다음 주를 보라.
\end{proof}

\begin{remark}
\label{remark-extensions-of-algebras}
$A\to B$와 $N$은 보조정리 \ref{lemma-extensions-of-algebras}에서와
같다고 하자. $\alpha:P\to B$를 $A$ 위에서의 $B$의 표시라 하자.
가환대수학, 절 \ref{algebra-section-netherlander}를 보라.
$J=\Ker(\alpha)$로 놓으면, $\alpha$에 딸린 소박한 코탄젠트 복합체
$\NL(\alpha)$는 복합체 $J/J^2\to\Omega_{P/A}\otimes_P B$이다. 또한
$$
\Ext^1_B(\NL(\alpha), N) =
\Coker\left(\Hom_B(\Omega_{P/A} \otimes_P B, N) \to \Hom_B(J/J^2, N)\right)
$$
이다. 이는 $\Omega_{P/A}$가 자유 가군이기 때문이다. 보조정리에 나오는
확대 $0\to N\to B'\to B\to0$을 생각하자. $P$는 $A$ 위의 다항식
대수이므로 $\alpha$를 $A$-대수 준동형 $\alpha':P'\to B'$로
들어 올릴 수 있다. $N$은 $B'$ 안에서 제곱이 영이므로
$\alpha'|_J:J\to N$은 $J\to J/J^2\to N$으로 분해된다.
보조정리의 대응은 이 확대를 위에 표시한 여핵 안의 준동형
$J/J^2\to N$의 동치류로 보낸다.
\end{remark}

\begin{lemma}
\label{lemma-extensions-of-algebras-functorial}
환 준동형 $A\to B\to C$, $B$-가군 $M$, $C$-가군 $N$,
$B$-선형 준동형 $c:M\to N$, 그리고 핵의 제곱이 영인 다음
$A$-대수 확대들이 주어졌다고 하자.
\begin{enumerate}
\item[(a)] $\xi\in\Ext^1_B(\NL_{B/A},M)$에 대응하는
$0\to M\to B'\to B\to0$,
\item[(b)] $\zeta\in\Ext^1_C(\NL_{C/A},N)$에 대응하는
$0\to N\to C'\to C\to0$.
\end{enumerate}
보조정리 \ref{lemma-extensions-of-algebras}를 보라. 이때 $B\to C$와
$c$에 양립하는 $A$-대수 준동형 $B'\to C'$가 존재하는 것과
$\xi$, $\zeta$가 $\Ext^1_B(\NL_{B/A},N)$의 같은 원소로 가는 것은
동치이다.
\end{lemma}

\begin{proof}
이 주장은 다음 준동형들이 있으므로 의미가 있다.
$$
\Ext^1_B(\NL_{B/A}, M) \to \Ext^1_B(\NL_{B/A}, N)
$$
첫 번째 것은 준동형 $M\to N$을 사용한다. 또한
$$
\Ext^1_C(\NL_{C/A}, N) \to \Ext^1_B(\NL_{C/A}, N) \to \Ext^1_B(\NL_{B/A}, N)
$$
도 있다. 여기서 첫 번째 화살표는 제한 준동형 $D(C)\to D(B)$를,
두 번째 화살표는 복합체의 표준 준동형
$\NL_{B/A}\to\NL_{C/A}$를 사용한다. 보조정리 \ref{lemma-huge-diagram}을
다음 도표에 적용하면 본 보조정리의 주장을 얻을 수 있다.
$$
\xymatrix{
& 0 \ar[r] & N \ar[r] & C' \ar[r] & C \ar[r] & 0 \\
& 0 \ar[r]|\hole & 0 \ar[u] \ar[r] &
A \ar[u] \ar[r]|\hole & A \ar[u] \ar[r] & 0 \\
0 \ar[r] & M \ar[ruu] \ar[r] & B' \ar[r] & B \ar[ruu] \ar[r] & 0 \\
0 \ar[r] & 0 \ar[ruu]|\hole \ar[u] \ar[r] &
A \ar[ruu]|\hole \ar[u] \ar[r] & A \ar[ruu]|\hole \ar[u] \ar[r] & 0
}
$$
여기에 보조정리 \ref{lemma-extensions-of-algebras}와
\ref{lemma-huge-diagram}의 증명에 나오는 구성들의 양립성을 사용한다.
그 양립성의 정확한 진술과 증명은 생략한다. 직접적인 논증은 다음 주를 보라.
\end{proof}

\begin{remark}
\label{remark-extensions-of-algebras-functorial}
$A\to B\to C$, $M$, $N$, $c:M\to N$,
$0\to M\to B'\to B\to0$, $\xi\in\Ext^1_B(\NL_{B/A},M)$,
$0\to N\to C'\to C\to0$, $\zeta\in\Ext^1_C(\NL_{C/A},N)$은
보조정리 \ref{lemma-extensions-of-algebras-functorial}에서와 같다고 하자.
$c:M\to N$에 따른 밀어내기로 다음 확대를 구성할 수 있다.
$$
\xymatrix{
0 \ar[r] & N \ar[r] & B'_1 \ar[r] & B \ar[r] & 0 \\
0 \ar[r] & M \ar[u]^c \ar[r] & B' \ar[u] \ar[r] & B \ar[u] \ar[r] & 0
}
$$
여기서 $M$은 반대각선으로 매장되며 $B'_1=(N\times B')/M$이다.
$B\to C$에 따른 당김으로 다음 확대를 구성할 수 있다.
$$
\xymatrix{
0 \ar[r] & N \ar[r] & C' \ar[r] & C \ar[r] & 0 \\
0 \ar[r] & N \ar[u] \ar[r] & B'_2 \ar[u] \ar[r] & B \ar[u] \ar[r] & 0
}
$$
여기서 $B'_2=C'\times_C B$는 환의 섬유곱이다. 간단한 도표 추적에
따르면 $B\to C$와 $c$에 양립하는 $A$-대수 준동형 $B'\to C'$가
존재하는 것과, $B$를 $N$으로 확대한 $A$-대수 확대들로서 $B'_1$과
$B'_2$가 동형인 것은 동치이다. 따라서 보조정리
\ref{lemma-extensions-of-algebras-functorial}을 보이려면, 보조정리
\ref{lemma-extensions-of-algebras}의 전단사 대응 아래에서 $B'_1$이
다음 준동형에 의한 $\xi$의 상에 대응하고
$\Ext^1_B(\NL_{B/A}, M) \to \Ext^1_B(\NL_{B/A}, N)$
또 $B'_2$가 다음 준동형에 의한 $\zeta$의 상에 대응함을 보이면 충분하다.
$\Ext^1_C(\NL_{C/A}, N) \to \Ext^1_B(\NL_{B/A}, N)$.
첫 번째 주장은 주 \ref{remark-extensions-of-algebras}에서 동치류를
구성한 방법으로부터 곧바로 따른다. 두 번째 주장을 위해 다음과 같은
$A$-대수의 가환 도표를 택하자.
$$
\xymatrix{
Q \ar[r]_\beta & C \\
P \ar[u]^\varphi \ar[r]^\alpha & B \ar[u]
}
$$
여기서 $\alpha$는 $A$ 위에서의 $B$의 표시이고 $\beta$는 $A$ 위에서의
$C$의 표시이다. 주 \ref{remark-extensions-of-algebras}와 거기에 적힌
참고 문헌을 보라. $J=\Ker(\alpha)$, $K=\Ker(\beta)$로 놓는다.
$\varphi$는 복합체의 준동형 $\NL(\alpha)\to\NL(\beta)$, 특히
$\bar\varphi:J/J^2\to K/K^2$를 유도한다. $\beta$의 올림인
$A$-대수 준동형 $\beta':Q\to C'$를 택한다. 그러면
$\alpha' = (\beta' \circ \varphi, \alpha) : P \to B'_2 = C' \times_C B$
는 $\alpha$의 올림이다. 이 선택 아래에서 $\beta'$가 유도하는 준동형
$K/K^2\to N$과 $\bar\varphi:J/J^2\to K/K^2$의 합성은
$\alpha'$를 $J/J^2$에 제한한 준동형이다. 주
\ref{remark-extensions-of-algebras}에서 동치류를 구성한 방법을 풀어 쓰면,
$B'_2$가 다음 준동형에 의한 $\zeta$의 상에 대응함을 알 수 있다.
$\Ext^1_C(\NL_{C/A}, N) \to \Ext^1_B(\NL_{B/A}, N)$.
\end{remark}

\begin{lemma}
\label{lemma-parametrize-solutions}
$0\to I\to A'\to A\to0$, $A\to B$, $c:I\to N$은
(\ref{equation-to-solve})에서와 같다고 하자. 보조정리
\ref{lemma-extensions-of-algebras}에 의해 $A$를 $I$로 확대한 확대 $A'$에
대응하는 원소를 $\xi\in\Ext^1_A(\NL_{A/A'},I)$라 쓰자. 해의 동형류
집합은 다음 준동형의 $\xi$의 상 위의 섬유와 표준적으로 전단사 대응한다.
$$
\Ext^1_B(\NL_{B/A'}, N) \to \Ext^1_A(\NL_{A/A'}, N)
$$
\end{lemma}

\begin{proof}
$A'\to B$와 $B$-가군 $N$에 보조정리
\ref{lemma-extensions-of-algebras}를 적용하면,
$\Ext^1_B(\NL_{B/A'},N)$의 원소 $\zeta$들이 $A'$-대수 확대
$0\to N\to B'\to B\to0$을 매개화함을 알 수 있다. $A'\to A\to B$와
$c:I\to N$에 보조정리
\ref{lemma-extensions-of-algebras-functorial}을 적용하면, $c$와
$A\to B$에 양립하는 $A'$-대수 준동형 $A'\to B'$가 존재하는 것과
$\zeta$가 $\xi$로 가는 것은 동치이다. 이는 물론 $B'$가
(\ref{equation-to-solve})의 해라는 것과 같은 말이다.
\end{proof}

\begin{remark}
\label{remark-parametrize-solutions}
보조정리 \ref{lemma-parametrize-solutions}의 상황에서는
$$
\Ext^1_A(\NL_{A/A'}, N) =
\Ext^1_B(\NL_{A/A'} \otimes_A^\mathbf{L} B, N) =
\Ext^1_B(\NL_{A/A'} \otimes_A B, N)
$$
임에 유의하자. 첫 번째 등식은 대수학 심화, 보조정리
\ref{more-algebra-lemma-tensor-hom-adjoint}에서, 두 번째 등식은
대수학 심화, 보조정리 \ref{more-algebra-lemma-tensor-NL}에서 따른다.
다음과 같은 복합체의 준동형들이 있다.
$$
\NL_{A/A'} \otimes_A B \to \NL_{B/A'} \to \NL_{B/A}
$$
이는 완전 삼각형에 매우 가깝다. 가환대수학, 보조정리
\ref{algebra-lemma-exact-sequence-NL}을 보라. 만약 이것이 완전
삼각형이라면 $\Ext^2_B(\NL_{B/A},N)$에서 $\xi$의 상이
(\ref{equation-to-solve})의 해가 존재하는 데 대한 장애 원소라고
결론 내릴 수 있을 것이다.
\end{remark}

\noindent
환 준동형 $A\to B$가 국소 완전 교차이면 해가 존재한다. 이는 일종의
올림 결과이다. 신토믹 환 준동형에 대해서는 환 준동형의 평활화,
명제 \ref{smoothing-proposition-lift-smooth}에서 이미 상당히 강한
올림 결과를 증명했음에 유의하자.

\begin{lemma}
\label{lemma-existence-lci}
$A\to B$가 국소 완전 교차 환 준동형이면
(\ref{equation-to-solve})의 해가 존재한다.
\end{lemma}

\begin{proof}[첫 번째 증명]
$B=A[x_1,\ldots,x_n]/J$로 쓰자. 대수학 심화, 정의
\ref{more-algebra-definition-local-complete-intersection}
에 의해 아이디얼 $J$는 Koszul-정칙이다. 따라서 $J$는 $H_1$-정칙이고
준정칙이다. 대수학 심화, 절 \ref{more-algebra-section-ideals}를 보라.
$J'\subset A'[x_1,\ldots,x_n]$를 $J$의 역상이라 하고,
$I[x_1,\ldots,x_n]$을
$A'[x_1,\ldots,x_n]\to A[x_1,\ldots,x_n]$의 핵이라 쓰자.
대수학 심화, 보조정리
\ref{more-algebra-lemma-conormal-sequence-H1-regular-ideal}에 의해
$I[x_1, \ldots, x_n] \cap (J')^2 = J'I[x_1, \ldots, x_n] =
JI[x_1, \ldots, x_n]$이다. 따라서 다음 짧은 완전열을 얻는다.
$$
0 \to I \otimes_A B \to J'/(J')^2 \to J/J^2 \to 0
$$
또 $J/J^2$은 사영 가군이므로(대수학 심화, 보조정리
\ref{more-algebra-lemma-quasi-regular-ideal-finite-projective})
이 완전열의 분할을 택할 수 있다.
$$
J'/(J')^2 = I \otimes_A B \oplus J/J^2 
$$
$(J')^2\subset J''\subset J'$를 위 분해의 두 번째 직합 인자로 가는
원소들의 집합이라 하자. 그러면
$$
0 \to I \otimes_A B \to A'[x_1, \ldots, x_n]/J'' \to B \to 0
$$
는 $N=I\otimes_A B$인 경우의 (\ref{equation-to-solve})의 해이다.
일반적인 경우는 주어진 준동형 $I\otimes_A B\to N$에 따른
밀어내기를 취하여 얻는다.
\end{proof}

\begin{proof}[두 번째 증명]
이 증명을 읽기 전에 주 \ref{remark-parametrize-solutions}를 읽어 두자.
대수학 심화, 보조정리 \ref{more-algebra-lemma-transitive-lci-at-end}에 의해
준동형들 $\NL_{A/A'}\otimes_A B\to\NL_{B/A'}\to\NL_{B/A}$는
$D(B)$에서 실제로 완전 삼각형을 이룬다. 따라서
$\Ext^2_B(\NL_{B/A},N)$이 영임을 보이면 충분하다. 대수학 심화,
보조정리 \ref{more-algebra-lemma-lci-NL}에 의해 복합체 $\NL_{B/A}$는
Tor-진폭이 $[-1,0]$에 놓이는 완전 복합체이다. 예를 들어 대수학 심화,
보조정리 \ref{more-algebra-lemma-splitting-unique}의 (1)에 의해 이 사실은
위 $\Ext^2$가 영임을 뜻한다.
\end{proof}







\section{환 달린 공간의 두꺼워짐}
\label{section-thickenings-spaces}

\noindent
이 절과 이어지는 몇 절에서는 다음 개념들을 사용한다.
\begin{enumerate}
\item 환 달린 공간 $(X',\mathcal{O}_{X'})$ 위의 아이디얼층
$\mathcal{I}\subset\mathcal{O}_{X'}$의 모든 국소 단면이 국소적으로
멱영이면 $\mathcal{I}$를 {\it 국소 멱영}이라고 한다. 가환대수학,
항목 \ref{algebra-item-ideal-locally-nilpotent}와 비교하라.
\item 환 달린 공간의 {\it 두꺼워짐}은 다음 조건들을 만족하는
환 달린 공간의 사상
$i : (X, \mathcal{O}_X) \to (X', \mathcal{O}_{X'})$
이다.
\begin{enumerate}
\item $i$는 위상 동형 $X\to X'$를 유도한다.
\item 준동형 $i^\sharp:\mathcal{O}_{X'}\to i_*\mathcal{O}_X$는
전사이다.
\item $i^\sharp$의 핵은 국소 멱영 아이디얼층이다.
\end{enumerate}
\item 환 달린 공간의 {\it 1차 두꺼워짐}은
$i : (X, \mathcal{O}_X) \to (X', \mathcal{O}_{X'})$
중 $\Ker(i^\sharp)$의 제곱이 영인 두꺼워짐이다.
\item {\it 두꺼워짐의 사상}, {\it 주어진 환 달린 공간 위의
두꺼워짐의 사상} 등은 자명한 방식으로 정의한다.
\end{enumerate}
$i:(X,\mathcal{O}_X)\to(X',\mathcal{O}_{X'})$가 환 달린 공간의
두꺼워짐이면 두 바탕 위상 공간을 동일시하고 $\mathcal{O}_X$,
$\mathcal{O}_{X'}$, $\mathcal{I}=\Ker(i^\sharp)$를 모두 $X=X'$ 위의
층으로 생각한다. 그러면 $\mathcal{O}_{X'}$-가군의 짧은 완전열
$$
0 \to \mathcal{I} \to \mathcal{O}_{X'} \to \mathcal{O}_X \to 0
$$
을 얻는다. 가군층, 보조정리 \ref{modules-lemma-i-star-equivalence}에 의해
$\mathcal{O}_X$-가군의 범주는 $\mathcal{I}$에 의해 소멸되는
$\mathcal{O}_{X'}$-가군의 범주와 동치이다. 특히 $i$가 1차 두꺼워짐이면
$\mathcal{I}$는 $\mathcal{O}_X$-가군이다.

\begin{situation}
\label{situation-morphism-thickenings}
두꺼워짐의 사상 $(f,f')$은 다음과 같은 환 달린 공간의 가환 도표로
주어진다.
\begin{equation}
\label{equation-morphism-thickenings}
\vcenter{
\xymatrix{
(X, \mathcal{O}_X) \ar[r]_i \ar[d]_f & (X', \mathcal{O}_{X'}) \ar[d]^{f'} \\
(S, \mathcal{O}_S) \ar[r]^t & (S', \mathcal{O}_{S'})
}
}
\end{equation}
여기서 가로 화살표들은 두꺼워짐이다. 이 상황에서는
$\mathcal{I} = \Ker(i^\sharp) \subset \mathcal{O}_{X'}$ 및
$\mathcal{J} = \Ker(t^\sharp) \subset \mathcal{O}_{S'}$로 놓는다.
바탕 위상 공간에서 $f=f'$이므로 위상적 역상 함자 $f^{-1}$과
$(f')^{-1}$을 동일시한다. 특히
$(f')^\sharp:f^{-1}\mathcal{O}_{S'}\to\mathcal{O}_{X'}$은 준동형
$f^{-1}\mathcal{J}\to\mathcal{I}$를 유도하고, 따라서 다음
$\mathcal{O}_{X'}$-가군 준동형을 유도한다.
$$
(f')^*\mathcal{J} \longrightarrow \mathcal{I}
$$
$i$와 $t$가 1차 두꺼워짐이면 $(f')^*\mathcal{J}=f^*\mathcal{J}$이고,
위 준동형은 $f^*\mathcal{J}\to\mathcal{I}$가 된다.
\end{situation}

\begin{definition}
\label{definition-strict-morphism-thickenings}
상황 \ref{situation-morphism-thickenings}에서 준동형
$(f')^*\mathcal{J}\longrightarrow\mathcal{I}$가 전사이면 $(f,f')$을
{\it 두꺼워짐의 엄격 사상}이라고 한다.
\end{definition}

\noindent
특히 다음 보조정리는 스킴의 두꺼워짐 사이의 사상
$(f,f'):(X\subset X')\to(S\subset S')$이 엄격인 것과
$X=S\times_{S'}X'$인 것이 동치임을 보인다.

\begin{lemma}
\label{lemma-strict-morphism-thickenings}
상황 \ref{situation-morphism-thickenings}에서 사상 $(f,f')$이 두꺼워짐의
엄격 사상인 것과 (\ref{equation-morphism-thickenings})이 환 달린 공간의
범주에서 데카르트 도표인 것은 동치이다.
\end{lemma}

\begin{proof}
생략한다.
\end{proof}




\section{환 달린 공간의 1차 두꺼워짐 위의 가군}
\label{section-modules-thickenings}

\noindent
이 절에서는 가군의 변형 이론에 필요한 몇 가지 준비 사항을 다룬다.
$i:(X,\mathcal{O}_X)\to(X',\mathcal{O}_{X'})$를 환 달린 공간의
1차 두꺼워짐이라 하자. 절 \ref{section-thickenings-spaces}에서 도입한
표기를 자유롭게 사용하며, 특히 바탕 위상 공간들을 동일시한다.
이 절에서는 다음과 같은 $\mathcal{O}_{X'}$-가군의 짧은 완전열을 다룬다.
\begin{equation}
\label{equation-extension}
0 \to \mathcal{K} \to \mathcal{F}' \to \mathcal{F} \to 0
\end{equation}
여기서 $\mathcal{F}$와 $\mathcal{K}$는 $\mathcal{O}_X$-가군이고
$\mathcal{F}'$은 $\mathcal{O}_{X'}$-가군이다. 이 상황에서는 표준적인
$\mathcal{O}_X$-가군 준동형
$$
c_{\mathcal{F}'} :
\mathcal{I} \otimes_{\mathcal{O}_X} \mathcal{F}
\longrightarrow
\mathcal{K}
$$
이 있다. 여기서 $\mathcal{I}=\Ker(i^\sharp)$이다. 구체적으로,
$\mathcal{I}$의 국소 단면 $f$와 $\mathcal{F}$의 국소 단면 $s$에 대해
$s$를 들어 올리는 $\mathcal{F}'$의 국소 단면을 $s'$라 하고
$c_{\mathcal{F}'}(f\otimes s)=fs'$로 둔다.

\begin{lemma}
\label{lemma-inf-map}
$i:(X,\mathcal{O}_X)\to(X',\mathcal{O}_{X'})$를 환 달린 공간의
1차 두꺼워짐이라 하자. 다음 확대들과
$$
0 \to \mathcal{K} \to \mathcal{F}' \to \mathcal{F} \to 0
\quad\text{그리고}\quad
0 \to \mathcal{L} \to \mathcal{G}' \to \mathcal{G} \to 0
$$
(\ref{equation-extension})에서와 같은 준동형
$\varphi:\mathcal{F}\to\mathcal{G}$ 및
$\psi:\mathcal{K}\to\mathcal{L}$이 주어졌다고 하자.
\begin{enumerate}
\item $\varphi$와 $\psi$에 양립하는 $\mathcal{O}_{X'}$-가군 준동형
$\varphi':\mathcal{F}'\to\mathcal{G}'$이 존재하면 다음 도표는 가환한다.
$$
\xymatrix{
\mathcal{I} \otimes_{\mathcal{O}_X} \mathcal{F}
\ar[r]_-{c_{\mathcal{F}'}} \ar[d]_{1 \otimes \varphi} &
\mathcal{K} \ar[d]^\psi \\
\mathcal{I} \otimes_{\mathcal{O}_X} \mathcal{G}
\ar[r]^-{c_{\mathcal{G}'}} &
\mathcal{L}
}
$$
\item $\varphi$와 $\psi$에 양립하는 $\mathcal{O}_{X'}$-가군 준동형
$\varphi':\mathcal{F}'\to\mathcal{G}'$들의 집합은 공집합이 아니면
$\Hom_{\mathcal{O}_X}(\mathcal{F},\mathcal{L})$ 아래의 주동차 공간이다.
\end{enumerate}
\end{lemma}

\begin{proof}
(1)은 준동형들의 정의에서 곧바로 따른다. (2)를 보자.
$\varphi'$과 $\varphi''$가 $\varphi$, $\psi$에 양립하는 두 준동형
$\mathcal{F}'\to\mathcal{G}'$이면 $\varphi'-\varphi''$는 다음과 같이
분해된다.
$$
\mathcal{F}' \to \mathcal{F} \to \mathcal{L} \to \mathcal{G}'
$$
가운데 준동형은 가군층, 보조정리
\ref{modules-lemma-i-star-equivalence}에 의해
$\Hom_{\mathcal{O}_X}(\mathcal{F},\mathcal{L})$의 유일한 원소에서 온다.
역으로 이 군의 원소 $\alpha$가 주어지면 위 합성에서 가운데에
$\alpha$를 넣어 얻은 준동형을 $\varphi'$에 더할 수 있다.
몇 가지 세부 사항은 생략한다.
\end{proof}

\begin{lemma}
\label{lemma-inf-obs-map}
$i:(X,\mathcal{O}_X)\to(X',\mathcal{O}_{X'})$를 환 달린 공간의
1차 두꺼워짐이라 하자. 다음 확대들과
$$
0 \to \mathcal{K} \to \mathcal{F}' \to \mathcal{F} \to 0
\quad\text{그리고}\quad
0 \to \mathcal{L} \to \mathcal{G}' \to \mathcal{G} \to 0
$$
(\ref{equation-extension})에서와 같은 준동형
$\varphi:\mathcal{F}\to\mathcal{G}$ 및
$\psi:\mathcal{K}\to\mathcal{L}$이 주어졌다고 하자. 다음 도표가
가환한다고 가정한다.
$$
\xymatrix{
\mathcal{I} \otimes_{\mathcal{O}_X} \mathcal{F}
\ar[r]_-{c_{\mathcal{F}'}} \ar[d]_{1 \otimes \varphi} &
\mathcal{K} \ar[d]^\psi \\
\mathcal{I} \otimes_{\mathcal{O}_X} \mathcal{G}
\ar[r]^-{c_{\mathcal{G}'}} &
\mathcal{L}
}
$$
그러면 다음 원소가 존재한다.
$$
o(\varphi, \psi) \in
\Ext^1_{\mathcal{O}_X}(\mathcal{F}, \mathcal{L})
$$
이 원소가 영이라는 것은 $\varphi$와 $\psi$에 양립하는 준동형
$\varphi':\mathcal{F}'\to\mathcal{G}'$이 존재하기 위한 필요충분조건이다.
\end{lemma}

\begin{proof}
다음 확대를 명시적으로 구성할 수 있다.
$$
0 \to \mathcal{L} \to \mathcal{H} \to \mathcal{F} \to 0
$$
$\mathcal{H}$는 다음 복합체의 가운데 코호몰로지로 잡는다.
$$
\mathcal{K}
\xrightarrow{1, - \psi}
\mathcal{F}' \oplus \mathcal{G}' \xrightarrow{\varphi, 1}
\mathcal{G}
$$
표기는 자명하다. 보조정리의 도표가 가환한다는 가정을 사용하여 국소
단면으로 계산하면 $\mathcal{H}$가 $\mathcal{I}$에 의해 소멸됨을 알 수
있다. 따라서 $\mathcal{H}$는 다음 군의 동치류를 정한다.
$$
\Ext^1_{\mathcal{O}_X}(\mathcal{F}, \mathcal{L})
\subset
\Ext^1_{\mathcal{O}_{X'}}(\mathcal{F}, \mathcal{L})
$$
끝으로 $\mathcal{H}$의 동치류는 $\psi$에 따른 확대 $\mathcal{F}'$의
밀어내기와 $\varphi$에 따른 확대 $\mathcal{G}'$의 당김의 차이다.
계산은 생략한다. 따라서 $\mathcal{H}$의 동치류가 영인 것과 다음과 같은
가환 도표가 존재하는 것은 동치이다.
$$
\xymatrix{
0 \ar[r] &
\mathcal{K} \ar[r] \ar[d]_{\psi} &
\mathcal{F}' \ar[r] \ar[d]_{\varphi'} &
\mathcal{F} \ar[r] \ar[d]_\varphi & 0\\
0 \ar[r] &
\mathcal{L} \ar[r] &
\mathcal{G}' \ar[r] &
\mathcal{G} \ar[r] & 0
}
$$
이는 원하는 결론이다.
\end{proof}

\begin{lemma}
\label{lemma-inf-ext}
$i:(X,\mathcal{O}_X)\to(X',\mathcal{O}_{X'})$를 환 달린 공간의
1차 두꺼워짐이라 하자. $\mathcal{O}_X$-가군 $\mathcal{F}$,
$\mathcal{K}$와 $\mathcal{O}_X$-선형 준동형
$c : \mathcal{I} \otimes_{\mathcal{O}_X} \mathcal{F} \to \mathcal{K}$.
이 주어졌다고 하자. $c_{\mathcal{F}'}=c$인 완전열
(\ref{equation-extension})이 하나라도 존재하면, 그러한 확대들의 동형류
집합은 $\Ext^1_{\mathcal{O}_X}(\mathcal{F},\mathcal{K})$ 아래의
주동차 공간이다.
\end{lemma}

\begin{proof}
다음 확대들이 주어졌다고 하자.
$$
0 \to \mathcal{K} \to \mathcal{F}'_1 \to \mathcal{F} \to 0
\quad\text{그리고}\quad
0 \to \mathcal{K} \to \mathcal{F}'_2 \to \mathcal{F} \to 0
$$
$c_{\mathcal{F}'_1}=c_{\mathcal{F}'_2}=c$라 하자. 확대군에서 두 확대의
차를 취하면(호몰로지 대수학, 절 \ref{homology-section-extensions} 참조)
다음 확대를 얻는다.
$$
0 \to \mathcal{K} \to \mathcal{E} \to \mathcal{F} \to 0
$$
여기서 $\mathcal{E}$는 $\mathcal{I}$에 의해 소멸된다. 국소 계산은
생략한다. 따라서 이 완전열은 $\mathcal{O}_X$-가군의 확대이다.
가군층, 보조정리 \ref{modules-lemma-i-star-equivalence}를 보라.
역으로 이러한 확대 $\mathcal{E}$가 주어지면
$c_{\mathcal{F}'}$를 바꾸지 않고 $\mathcal{E}$를
$\mathcal{O}_{X'}$-확대 $\mathcal{F}'$에 더할 수 있다. 몇 가지
세부 사항은 생략한다.
\end{proof}

\begin{lemma}
\label{lemma-inf-obs-ext}
$i:(X,\mathcal{O}_X)\to(X',\mathcal{O}_{X'})$를 환 달린 공간의
1차 두꺼워짐이라 하자. $\mathcal{O}_X$-가군 $\mathcal{F}$,
$\mathcal{K}$와 $\mathcal{O}_X$-선형 준동형
$c : \mathcal{I} \otimes_{\mathcal{O}_X} \mathcal{F} \to \mathcal{K}$.
이 주어졌다고 하자. 그러면 다음 원소가 존재한다.
$$
o(\mathcal{F}, \mathcal{K}, c) \in
\Ext^2_{\mathcal{O}_X}(\mathcal{F}, \mathcal{K})
$$
이 원소가 영이라는 것은 $c_{\mathcal{F}'}=c$인 완전열
(\ref{equation-extension})이 존재하기 위한 필요충분조건이다.
\end{lemma}

\begin{proof}
먼저 $\mathcal{K}$가 단사 $\mathcal{O}_X$-가군이면
$c_{\mathcal{F}'}=c$인 완전열 (\ref{equation-extension})이 실제로
존재함을 보이자. 평탄 $\mathcal{O}_{X'}$-가군 $\mathcal{H}'$과
전사 준동형
$\mathcal{H}' \to \mathcal{F}$
을 택한다. 가군층, 보조정리 \ref{modules-lemma-module-quotient-flat}를
보라. 그 핵을 $\mathcal{J}\subset\mathcal{H}'$라 하자.
$\mathcal{H}'$은 평탄하므로
$$
\mathcal{I} \otimes_{\mathcal{O}_{X'}} \mathcal{H}' =
\mathcal{I}\mathcal{H}'
\subset \mathcal{J} \subset \mathcal{H}'
$$
이다. 다음 준동형은 $\mathcal{I}\mathcal{J}$를 영으로 보낸다.
$$
\mathcal{I}\mathcal{H}' =
\mathcal{I} \otimes_{\mathcal{O}_{X'}} \mathcal{H}'
\longrightarrow
\mathcal{I} \otimes_{\mathcal{O}_{X'}} \mathcal{F} =
\mathcal{I} \otimes_{\mathcal{O}_X} \mathcal{F}
$$
실제로 $f$가 $\mathcal{I}$의 국소 단면이고 $s$가 $\mathcal{H}$의
국소 단면이면 $fs$는 $f\otimes\overline{s}$로 간다. 여기서
$\overline{s}$는 $\mathcal{F}$에서 $s$의 상이다. 따라서 다음과 같은
$\mathcal{O}_X$-가군의 도표를 얻는다.
$$
\xymatrix{
\mathcal{I}\mathcal{H}'/\mathcal{I}\mathcal{J}
\ar@{^{(}->}[r] \ar[d] &
\mathcal{J}/\mathcal{I}\mathcal{J} \ar@{..>}[d]_\gamma \\
\mathcal{I} \otimes_{\mathcal{O}_X} \mathcal{F} \ar[r]^-c &
\mathcal{K}
}
$$
$\mathcal{K}$가 단사 $\mathcal{O}_X$-가군이면 점선 화살표를 얻는다.
$\gamma$와 $\mathcal{J}\to\mathcal{J}/\mathcal{I}\mathcal{J}$의 합성을
$\gamma':\mathcal{J}\to\mathcal{K}$라 쓰자. 국소 계산에 따르면 다음
밀어내기는
$$
\xymatrix{
0 \ar[r] &
\mathcal{J} \ar[r] \ar[d]_{\gamma'} &
\mathcal{H}' \ar[r] \ar[d] &
\mathcal{F} \ar[r] \ar@{=}[d] &
0 \\
0 \ar[r] &
\mathcal{K} \ar[r] &
\mathcal{F}' \ar[r] &
\mathcal{F} \ar[r] &
0
}
$$
보조정리에서 제기한 문제의 해이다.

\medskip\noindent
이제 일반적인 경우를 다루자. $\mathcal{K}'$가 단사
$\mathcal{O}_X$-가군이 되도록 매장 $\mathcal{K}\subset\mathcal{K}'$을
택한다. 그 몫을 $\mathcal{Q}$라 하면 완전열
$$
0 \to \mathcal{K} \to \mathcal{K}' \to \mathcal{Q} \to 0
$$
을 얻는다. 합성 준동형을
$c':\mathcal{I}\otimes_{\mathcal{O}_X}\mathcal{F}\to\mathcal{K}'$라
쓰자. 앞 문단에 의해 다음과 같은 완전열이 존재한다.
$$
0 \to \mathcal{K}' \to \mathcal{E}' \to \mathcal{F} \to 0
$$
이는 (\ref{equation-extension})에서와 같고 $c_{\mathcal{E}'}=c'$이다.
$c'$와 $\mathcal{K}'\to\mathcal{Q}$의 합성은 영이므로,
$\mathcal{K}'\to\mathcal{Q}$에 따른 $\mathcal{E}'$의 밀어내기는
다음 확대이다.
$$
0 \to \mathcal{Q} \to \mathcal{D}' \to \mathcal{F} \to 0
$$
이는 (\ref{equation-extension})에서와 같고 $c_{\mathcal{D}'}=0$이다.
이는 정확히 $\mathcal{D}'$가 $\mathcal{I}$에 의해 소멸됨을 뜻한다.
다시 말해 $\mathcal{D}'$는 $\mathcal{O}_X$-가군의 확대이며 다음 원소를
정한다.
$$
o(\mathcal{F}, \mathcal{K}, c) \in
\Ext^1_{\mathcal{O}_X}(\mathcal{F}, \mathcal{Q}) =
\Ext^2_{\mathcal{O}_X}(\mathcal{F}, \mathcal{K})
$$
등식은 위 완전열에 딸린 긴 완전 코호몰로지 열과 단사 가군
$\mathcal{K}'$로 가는 고차 Ext 군들의 소멸에서 따른다.
$o(\mathcal{F},\mathcal{K},c)=0$이면 분할
$s:\mathcal{F}\to\mathcal{D}'$를 택하고
$$
\mathcal{F}' = \Ker(\mathcal{E}' \to \mathcal{D}'/s(\mathcal{F}))
$$
로 놓을 수 있다. 그러면 완전한 두 행을 가진 다음 도표를 얻으며,
$$
\xymatrix{
0 \ar[r] &
\mathcal{K} \ar[r] \ar[d] &
\mathcal{F}' \ar[r] \ar[d] &
\mathcal{F} \ar[r] \ar@{=}[d] &
0 \\
0 \ar[r] &
\mathcal{K}' \ar[r] &
\mathcal{E}' \ar[r] &
\mathcal{F} \ar[r] & 0
}
$$
이는 $c_{\mathcal{F}'}=c$임을 보인다. 역으로 $\mathcal{F}'$가 존재하면
보조정리 \ref{lemma-inf-ext}와 단사 가군 $\mathcal{K}'$로 가는 고차
Ext 군들의 소멸에 의해, $\mathcal{K}\to\mathcal{K}'$에 따른
$\mathcal{F}'$의 밀어내기는 $\mathcal{E}'$와 동형이다. 이로부터 위와
같은 도표를 얻고, $\mathcal{D}'$가 확대로서 분할됨을 알 수 있다.
즉 동치류 $o(\mathcal{F},\mathcal{K},c)$는 영이다.
\end{proof}

\begin{remark}
\label{remark-trivial-thickening}
$(X,\mathcal{O}_X)$를 환 달린 공간이라 하자. 1차 두꺼워짐
$i:(X,\mathcal{O}_X)\to(X',\mathcal{O}_{X'})$에 대하여 그
왼쪽 역인 환 달린 공간의 사상
$\pi:(X',\mathcal{O}_{X'})\to(X,\mathcal{O}_X)$가 존재하면 $i$를
{\it 자명하다}고 한다. 그러한 사상 $\pi$의 선택을 이 1차 두꺼워짐의
{\it 자명화}라고 한다. $\pi$가 주어지면 다음 분할을 얻는다.
\begin{equation}
\label{equation-splitting}
\mathcal{O}_{X'} = \mathcal{O}_X \oplus \mathcal{I}
\end{equation}
$\pi^\sharp$로 전사 준동형 $\mathcal{O}_{X'}\to\mathcal{O}_X$를
분할하여 얻은 $X$ 위의 대수층의 분할이다. 역으로 이러한 분할은 사상
$\pi$를 정한다. $(X,\mathcal{O}_X)$의 자명화된 1차 두꺼워짐들의 범주는
$\mathcal{O}_X$-가군의 범주와 동치이다.
\end{remark}

\begin{remark}
\label{remark-trivial-extension}
$i:(X,\mathcal{O}_X)\to(X',\mathcal{O}_{X'})$를 환 달린 공간의
자명한 1차 두꺼워짐이라 하고
$\pi:(X',\mathcal{O}_{X'})\to(X,\mathcal{O}_X)$를 그 자명화라 하자.
두 $\mathcal{O}_X$-가군과 준동형
$c : \mathcal{I} \otimes_{\mathcal{O}_X} \mathcal{F} \to \mathcal{K}$
으로 이루어진 임의의 삼중항 $(\mathcal{F},\mathcal{K},c)$에 대하여
다음과 같이 놓을 수 있다.
$$
\mathcal{F}'_{c, triv} = \mathcal{F} \oplus \mathcal{K}
$$
$\pi$에 딸린 분할 (\ref{equation-splitting})과 준동형 $c$를 사용하여
$\mathcal{O}_{X'}$-가군 구조를 정의하면 확대
(\ref{equation-extension})을 얻는다. $\mathcal{F}'_{c,triv}$를 $c$와
자명화 $\pi$에 대응하는, $\mathcal{F}$를 $\mathcal{K}$로 확대한
{\it 자명한 확대}라고 부른다. (\ref{equation-extension})과 같은 임의의
확대 $\mathcal{F}'$에 대하여
$\pi^\sharp:\mathcal{O}_X\to\mathcal{O}_{X'}$를 사용하면
$\mathcal{F}'$를 $\mathcal{O}_X$-가군의 확대로 볼 수 있고, 따라서
$\Ext^1_{\mathcal{O}_X}(\mathcal{F},\mathcal{K})$의 동치류
$\xi_{\mathcal{F}'}$를 얻는다. 보조정리 \ref{lemma-inf-ext}에 의해
$\mathcal{F}' \mapsto \xi_{\mathcal{F}'}$
는 전단사 대응
$$
\left\{
\begin{matrix}
\text{확대의 동형류}\\
\mathcal{F}'\text{는 (\ref{equation-extension})과 같고 }c = c_{\mathcal{F}'}
\end{matrix}
\right\}
\longrightarrow
\Ext^1_{\mathcal{O}_X}(\mathcal{F}, \mathcal{K})
$$
를 유도한다. 또한 자명한 확대 $\mathcal{F}'_{c,triv}$는 영 동치류로 간다.
\end{remark}

\begin{remark}
\label{remark-extension-functorial}
$(X,\mathcal{O}_X)$를 환 달린 공간이라 하자.
$(X,\mathcal{O}_X)\to(X'_i,\mathcal{O}_{X'_i})$, $i=1,2$를
아이디얼층이 $\mathcal{I}_i$인 1차 두꺼워짐이라 하자.
$h:(X'_1,\mathcal{O}_{X'_1})\to(X'_2,\mathcal{O}_{X'_2})$를
$(X,\mathcal{O}_X)$의 1차 두꺼워짐 사이의 사상이라 하자. 도표는
$$
\xymatrix{
& (X, \mathcal{O}_X) \ar[ld] \ar[rd] & \\
(X'_1, \mathcal{O}_{X'_1}) \ar[rr]^h & & 
(X'_2, \mathcal{O}_{X'_2})
}
$$
와 같다. 특히 $h^\sharp:\mathcal{O}_{X'_2}\to\mathcal{O}_{X'_1}$은
$\mathcal{O}_X$-가군 준동형 $\mathcal{I}_2\to\mathcal{I}_1$을
유도한다. $\mathcal{F}$를 $\mathcal{O}_X$-가군이라 하자.
각 $i=1,2$에 대하여 $(\mathcal{K}_i,c_i)$는
$\mathcal{O}_X$-가군 $\mathcal{K}_i$와 준동형
$c_i : \mathcal{I}_i \otimes_{\mathcal{O}_X} \mathcal{F} \to
\mathcal{K}_i$로 이루어졌다고 하자. 또한 다음 도표를 가환시키는
$\mathcal{O}_X$-가군 준동형 $\mathcal{K}_2\to\mathcal{K}_1$이
주어졌다고 하자.
$$
\xymatrix{
\mathcal{I}_2 \otimes_{\mathcal{O}_X} \mathcal{F}
\ar[r]_-{c_2} \ar[d] &
\mathcal{K}_2 \ar[d] \\
\mathcal{I}_1 \otimes_{\mathcal{O}_X} \mathcal{F}
\ar[r]^-{c_1} &
\mathcal{K}_1
}
$$
그러면 다음과 같은 표준적인 함자성이 있다.
$$
\left\{
\begin{matrix}
\mathcal{F}'_2\text{는 (\ref{equation-extension})과 같고 }\\
c_2 = c_{\mathcal{F}'_2}\text{이며 }\mathcal{K} = \mathcal{K}_2
\end{matrix}
\right\}
\longrightarrow
\left\{
\begin{matrix}
\mathcal{F}'_1\text{는 (\ref{equation-extension})과 같고 }\\
c_1 = c_{\mathcal{F}'_1}\text{이며 }\mathcal{K} = \mathcal{K}_1
\end{matrix}
\right\}
$$
구체적으로 모든 층 $\mathcal{O}_X$, $\mathcal{O}_{X'_i}$,
$\mathcal{F}$, $\mathcal{K}_i$ 등을 $X$ 위의 층으로 본다.
$\mathcal{F}'_2$가 주어지면 $\mathcal{F}'_1$을 다음 확대 도표에
들어가는 밀어내기로 정의한다.
$$
\xymatrix{
0 \ar[r] &
\mathcal{K}_2 \ar[r] \ar[d] &
\mathcal{F}'_2 \ar[r] \ar[d] &
\mathcal{F} \ar@{=}[d] \ar[r] & 0 \\
0 \ar[r] &
\mathcal{K}_1 \ar[r] &
\mathcal{F}'_1 \ar[r] &
\mathcal{F} \ar[r] & 0
}
$$
이 밀어내기 위의 $\mathcal{O}_{X'_1}$-가군 구조의 구성은 생략한다.
이 구성에는 $c_1$, $c_2$가 들어가는 도표의 가환성이 사용된다.
\end{remark}

\begin{remark}
\label{remark-trivial-extension-functorial}
$(X,\mathcal{O}_X)$, $(X,\mathcal{O}_X)\to(X'_i,\mathcal{O}_{X'_i})$,
$\mathcal{I}_i$, $h:(X'_1,\mathcal{O}_{X'_1})\to
(X'_2,\mathcal{O}_{X'_2})$는 주 \ref{remark-extension-functorial}에서와
같다고 하자. $\pi_1=h\circ\pi_2$를 만족하는 자명화
$\pi_i:X'_i\to X$가 주어졌다고 하자. 다시 말해 $h$가
$(X,\mathcal{O}_X)$의 자명화된 1차 두꺼워짐 사이의 사상이라고
가정한다. 각 $i=1,2$에 대하여 $(\mathcal{K}_i,c_i)$는
$\mathcal{O}_X$-가군 $\mathcal{K}_i$와 준동형
$c_i : \mathcal{I}_i \otimes_{\mathcal{O}_X} \mathcal{F} \to
\mathcal{K}_i$로 이루어졌다고 하자. 또한 다음 도표를 가환시키는
$\mathcal{O}_X$-가군 준동형 $\mathcal{K}_2\to\mathcal{K}_1$이
주어졌다고 하자.
$$
\xymatrix{
\mathcal{I}_2 \otimes_{\mathcal{O}_X} \mathcal{F}
\ar[r]_-{c_2} \ar[d] &
\mathcal{K}_2 \ar[d] \\
\mathcal{I}_1 \otimes_{\mathcal{O}_X} \mathcal{F}
\ar[r]^-{c_1} &
\mathcal{K}_1
}
$$
이 상황에서 주 \ref{remark-trivial-extension}의 구성은 다음 가환
도표를 유도한다.
$$
\xymatrix{
\{\mathcal{F}'_2\text{는 (\ref{equation-extension})과 같고 }
c_2 = c_{\mathcal{F}'_2}\text{이며 }\mathcal{K} = \mathcal{K}_2\}
\ar[d] \ar[rr] & &
\Ext^1_{\mathcal{O}_X}(\mathcal{F}, \mathcal{K}_2) \ar[d] \\
\{\mathcal{F}'_1\text{는 (\ref{equation-extension})과 같고 }
c_1 = c_{\mathcal{F}'_1}\text{이며 }\mathcal{K} = \mathcal{K}_1\}
\ar[rr] & &
\Ext^1_{\mathcal{O}_X}(\mathcal{F}, \mathcal{K}_1)
}
$$
여기서 오른쪽 세로 준동형은 $\Ext$의 함자성과 준동형
$\mathcal{K}_2\to\mathcal{K}_1$로 주어지고, 왼쪽 세로 준동형은
주 \ref{remark-extension-functorial}에서 구성한 것이다.
\end{remark}

\begin{remark}
\label{remark-short-exact-sequence-thickenings}
$(X,\mathcal{O}_X)$를 환 달린 공간이라 하자. $(X,\mathcal{O}_X)$의
1차 두꺼워짐 사이의 사상열
$$
(X'_1, \mathcal{O}_{X'_1}) \to
(X'_2, \mathcal{O}_{X'_2}) \to
(X'_3, \mathcal{O}_{X'_3})
$$
에서 아이디얼층 $\mathcal{I}_i$ 사이의 대응하는 준동형들이
$\mathcal{O}_X$-가군의 복합체
$\mathcal{I}_3\to\mathcal{I}_2\to\mathcal{I}_1$, 즉 합성이 영인 열을
이루면, 주어진 두꺼워짐의 열을 {\it 복합체}라고 한다. 이 경우 합성
$(X'_1,\mathcal{O}_{X'_1})\to(X_3',\mathcal{O}_{X'_3})$은
$(X,\mathcal{O}_X)\to(X'_3,\mathcal{O}_{X'_3})$을 거쳐 분해된다.
즉 $(X,\mathcal{O}_X)$의 1차 두꺼워짐
$(X'_1,\mathcal{O}_{X'_1})$은 자명하며 표준 자명화
$\pi:(X'_1,\mathcal{O}_{X'_1})\to(X,\mathcal{O}_X)$를 갖는다.

\medskip\noindent
$(X,\mathcal{O}_X)$의 1차 두꺼워짐 사이의 사상열
$$
(X'_1, \mathcal{O}_{X'_1}) \to
(X'_2, \mathcal{O}_{X'_2}) \to
(X'_3, \mathcal{O}_{X'_3})
$$
에서 아이디얼층 사이의 대응하는 준동형들이 $\mathcal{O}_X$-가군의
짧은 완전열
$$
0 \to \mathcal{I}_3 \to \mathcal{I}_2 \to \mathcal{I}_1 \to 0
$$
을 이루면, 주어진 두꺼워짐의 열을 {\it 짧은 완전열}이라고 한다.
\end{remark}

\begin{remark}
\label{remark-complex-thickenings-and-ses-modules}
$(X,\mathcal{O}_X)$를 환 달린 공간이라 하고 $\mathcal{F}$를
$\mathcal{O}_X$-가군이라 하자. 다음 열을
$$
(X'_1, \mathcal{O}_{X'_1}) \to
(X'_2, \mathcal{O}_{X'_2}) \to
(X'_3, \mathcal{O}_{X'_3})
$$
$(X,\mathcal{O}_X)$의 1차 두꺼워짐의 복합체라 하자. 주
\ref{remark-short-exact-sequence-thickenings}를 보라. 각 $i=1,2,3$에
대하여 $(\mathcal{K}_i,c_i)$는 $\mathcal{O}_X$-가군
$\mathcal{K}_i$와 준동형
$c_i : \mathcal{I}_i \otimes_{\mathcal{O}_X} \mathcal{F} \to
\mathcal{K}_i$로 이루어졌다고 하자. 다음과 같은
$\mathcal{O}_X$-가군의 짧은 완전열이 주어졌다고 하자.
$$
0 \to \mathcal{K}_3 \to \mathcal{K}_2 \to \mathcal{K}_1 \to 0
$$
또 다음 두 도표가 가환한다고 가정한다.
$$
\vcenter{
\xymatrix{
\mathcal{I}_2 \otimes_{\mathcal{O}_X} \mathcal{F}
\ar[r]_-{c_2} \ar[d] &
\mathcal{K}_2 \ar[d] \\
\mathcal{I}_1 \otimes_{\mathcal{O}_X} \mathcal{F}
\ar[r]^-{c_1} &
\mathcal{K}_1
}
}
\quad\text{그리고}\quad
\vcenter{
\xymatrix{
\mathcal{I}_3 \otimes_{\mathcal{O}_X} \mathcal{F}
\ar[r]_-{c_3} \ar[d] &
\mathcal{K}_3 \ar[d] \\
\mathcal{I}_2 \otimes_{\mathcal{O}_X} \mathcal{F}
\ar[r]^-{c_2} &
\mathcal{K}_2
}
}
$$
끝으로 $\mathcal{O}_{X'_2}$-가군의 확대
$$
0 \to \mathcal{K}_2 \to \mathcal{F}'_2 \to \mathcal{F} \to 0
$$
이 주어졌다고 하자. 이는 $\mathcal{K}=\mathcal{K}_2$인
(\ref{equation-extension})과 같고 $c_{\mathcal{F}'_2}=c_2$이다.
주 \ref{remark-extension-functorial}의 함자성을 적용하여 $X'_1$ 위의
확대 $\mathcal{F}'_1$을 얻을 수 있다. 이 특별한 경우의
$\mathcal{F}'_1$은 아래에서 설명한다. 주
\ref{remark-trivial-extension}에 따라, 주
\ref{remark-short-exact-sequence-thickenings}의 표준 자명화
$\pi : (X'_1, \mathcal{O}_{X'_1}) \to (X, \mathcal{O}_X)$
를 사용하면
$\xi_{\mathcal{F}'_1} \in
\Ext^1_{\mathcal{O}_X}(\mathcal{F}, \mathcal{K}_1)$을 얻는다.
마지막으로 다음 장애 원소가 있다.
$$
o(\mathcal{F}, \mathcal{K}_3, c_3) \in
\Ext^2_{\mathcal{O}_X}(\mathcal{F}, \mathcal{K}_3)
$$
보조정리 \ref{lemma-inf-obs-ext}를 보라. 이 상황에서 짧은 완전열
$0\to\mathcal{K}_3\to\mathcal{K}_2\to\mathcal{K}_1\to0$에서 오는
표준 준동형
$$
\partial :
\Ext^1_{\mathcal{O}_X}(\mathcal{F}, \mathcal{K}_1)
\longrightarrow
\Ext^2_{\mathcal{O}_X}(\mathcal{F}, \mathcal{K}_3)
$$
이 $\xi_{\mathcal{F}'_1}$을 장애 동치류
$o(\mathcal{F},\mathcal{K}_3,c_3)$로 보낸다고 {\bf 주장}한다.

\medskip\noindent
이 주장을 증명하기 위해 $\mathcal{K}$가 단사 $\mathcal{O}_X$-가군이
되도록 매장 $j:\mathcal{K}_3\to\mathcal{K}$를 택한다. $j$를 준동형
$j':\mathcal{K}_2\to\mathcal{K}$로 들어 올릴 수 있다.
$\mathcal{E}'_2=j'_*\mathcal{F}'_2$를 $j'$에 따른
$\mathcal{F}'_2$의 밀어내기로 놓는다. 그러면
$c_{\mathcal{E}'_2}=j'\circ c_2$이다. 도표로 쓰면 다음과 같다.
$$
\xymatrix{
0 \ar[r] &
\mathcal{K}_2 \ar[r] \ar[d]_{j'} &
\mathcal{F}'_2 \ar[r] \ar[d] &
\mathcal{F} \ar[r] \ar[d] & 0 \\
0 \ar[r] &
\mathcal{K} \ar[r] &
\mathcal{E}'_2 \ar[r] &
\mathcal{F} \ar[r] & 0
}
$$
$\mathcal{E}'_3=\mathcal{E}'_2$로 놓되,
$\mathcal{O}_{X'_3}\to\mathcal{O}_{X'_2}$를 통해
$\mathcal{O}_{X'_3}$-가군으로 본다. 그러면
$c_{\mathcal{E}'_3}=j\circ c_3$이다. 보조정리
\ref{lemma-inf-obs-ext}의 증명은 $o(\mathcal{F},\mathcal{K}_3,c_3)$를
다음 $\mathcal{O}_X$-가군 확대의 동치류의 경계로 구성한다.
$$
0 \to
\mathcal{K}/\mathcal{K}_3 \to
\mathcal{E}'_3/\mathcal{K}_3 \to
\mathcal{F} \to 0
$$
한편 $\mathcal{F}'_1=\mathcal{F}'_2/\mathcal{K}_3$임에 유의하자.
따라서 $\xi_{\mathcal{F}'_1}$은 다음 확대의 동치류이다.
$$
0 \to \mathcal{K}_2/\mathcal{K}_3 \to \mathcal{F}'_2/\mathcal{K}_3
\to \mathcal{F} \to 0
$$
여기서는 표준 자명화
$\pi:(X'_1,\mathcal{O}_{X'_1})\to(X,\mathcal{O}_X)$의
$\pi^\sharp$를 사용하여 이를 $\mathcal{O}_X$-가군의 열로 본다.
이제 다음 가환 도표가 있으므로 주장이 따른다.
$$
\xymatrix{
0 \ar[r] &
\mathcal{K}_2/\mathcal{K}_3 \ar[r] \ar[d] &
\mathcal{F}'_2/\mathcal{K}_3 \ar[r] \ar[d] &
\mathcal{F} \ar[r] \ar[d] & 0 \\
0 \ar[r] &
\mathcal{K}/\mathcal{K}_3 \ar[r] &
\mathcal{E}'_3/\mathcal{K}_3 \ar[r] &
\mathcal{F} \ar[r] & 0
}
$$
이 도표는 위에서 주어진 $\mathcal{O}_X$-가군 구조들에 대해
$\mathcal{O}_X$-선형이다.
\end{remark}








\section{환 달린 공간 위 가군의 무한소 변형}
\label{section-deformation-modules}

\noindent
$i : (X, \mathcal{O}_X) \to (X', \mathcal{O}_{X'})$를 환 달린 공간의
1차 두꺼워짐이라 하자. 절 \ref{section-thickenings-spaces}에서 도입한
표기를 자유롭게 사용한다. $\mathcal{F}'$를 $\mathcal{O}_{X'}$-가군이라
하고 $\mathcal{F} = i^*\mathcal{F}'$라 놓자. 이 상황에서는 다음
짧은 완전열이 있다.
$$
0 \to \mathcal{I}\mathcal{F}' \to \mathcal{F}' \to \mathcal{F} \to 0
$$
이는 $\mathcal{O}_{X'}$-가군의 열이다. $\mathcal{I}^2 = 0$이므로
$\mathcal{I}\mathcal{F}'$ 위의 $\mathcal{O}_{X'}$-가군 구조는 유일한
$\mathcal{O}_X$-가군 구조에서 온다. 따라서 위의 열은
(\ref{equation-extension})과 같은 확대이다. 특히
$\mathcal{F}' = \mathcal{O}_{X'}$이면
$i^*\mathcal{O}_{X'} = \mathcal{O}_X$이고
$\mathcal{I}\mathcal{O}_{X'} = \mathcal{I}$이므로 다음 구조층의 열을
되찾는다.
$$
0 \to \mathcal{I} \to \mathcal{O}_{X'} \to \mathcal{O}_X \to 0
$$

\begin{lemma}
\label{lemma-inf-map-special}
$i : (X, \mathcal{O}_X) \to (X', \mathcal{O}_{X'})$를 환 달린 공간의
1차 두꺼워짐이라 하자. $\mathcal{F}'$, $\mathcal{G}'$를
$\mathcal{O}_{X'}$-가군이라 하고
$\mathcal{F} = i^*\mathcal{F}'$ 및 $\mathcal{G} = i^*\mathcal{G}'$라
놓자. $\varphi : \mathcal{F} \to \mathcal{G}$를
$\mathcal{O}_X$-선형 준동형이라 하자. $\varphi$를
$\mathcal{O}_{X'}$-선형 준동형
$\varphi' : \mathcal{F}' \to \mathcal{G}'$로 올리는 방법들의 집합은
공집합이 아니면 다음 군 아래의 주동차 공간이다.
$\Hom_{\mathcal{O}_X}(\mathcal{F}, \mathcal{I}\mathcal{G}')$.
\end{lemma}

\begin{proof}
이는 보조정리 \ref{lemma-inf-map}의 특수한 경우이지만 직접 증명도
제시한다. 다음 가군의 짧은 완전열들이 있다.
$$
0 \to \mathcal{I} \to \mathcal{O}_{X'} \to \mathcal{O}_X \to 0
\quad\text{그리고}\quad
0 \to \mathcal{I}\mathcal{G}' \to \mathcal{G}' \to \mathcal{G} \to 0
$$
$\mathcal{F}'$에 대해서도 마찬가지이다. $\mathcal{I}$의 제곱이
영이므로 $\mathcal{I}$와 $\mathcal{I}\mathcal{G}'$ 위의
$\mathcal{O}_{X'}$-가군 구조는 유일한 $\mathcal{O}_X$-가군 구조에서
온다. 따라서
$$
\Hom_{\mathcal{O}_{X'}}(\mathcal{F}', \mathcal{I}\mathcal{G}') =
\Hom_{\mathcal{O}_X}(\mathcal{F}, \mathcal{I}\mathcal{G}')
\quad\text{그리고}\quad
\Hom_{\mathcal{O}_{X'}}(\mathcal{F}', \mathcal{G}) =
\Hom_{\mathcal{O}_X}(\mathcal{F}, \mathcal{G})
$$
이고, 이제 보조정리는 다음 완전열에서 따른다.
$$
0 \to \Hom_{\mathcal{O}_{X'}}(\mathcal{F}', \mathcal{I}\mathcal{G}') \to
\Hom_{\mathcal{O}_{X'}}(\mathcal{F}', \mathcal{G}') \to
\Hom_{\mathcal{O}_{X'}}(\mathcal{F}', \mathcal{G})
$$
호몰로지, 보조정리 \ref{homology-lemma-check-exactness}를 보라.
\end{proof}

\begin{lemma}
\label{lemma-deform-module}
$(f, f')$를 상황 \ref{situation-morphism-thickenings}에서와 같은
환 달린 공간의 1차 두꺼워짐 사이의 사상이라 하자.
$\mathcal{F}'$를 $\mathcal{O}_{X'}$-가군이라 하고
$\mathcal{F} = i^*\mathcal{F}'$라 놓자. $\mathcal{F}$가 $S$ 위에서
평탄하고 $(f, f')$가 두꺼워짐의 엄격 사상이라고 가정하자
(정의 \ref{definition-strict-morphism-thickenings}). 그러면 다음
조건들이 동치이다.
\begin{enumerate}
\item $\mathcal{F}'$는 $S'$ 위에서 평탄하다.
\item 표준 준동형
$f^*\mathcal{J} \otimes_{\mathcal{O}_X} \mathcal{F} \to
\mathcal{I}\mathcal{F}'$
은 동형사상이다.
\end{enumerate}
더욱이 이 경우 다음 준동형들은 동형사상이다.
$$
f^*\mathcal{J} \otimes_{\mathcal{O}_X} \mathcal{F} \to
\mathcal{I} \otimes_{\mathcal{O}_X} \mathcal{F} \to
\mathcal{I}\mathcal{F}'
$$
\end{lemma}

\begin{proof}
$(f, f')$가 두꺼워짐의 엄격 사상이므로 준동형
$f^*\mathcal{J} \to \mathcal{I}$은 전사이다. 따라서 마지막 명제는
(2)에서 따른다.

\medskip\noindent
(1)과 (2)의 동치를 증명하자. 이 조건들은 줄기에서 확인해도 된다.
$x \in X \subset X'$를 상이 $s = f(x) \in S \subset S'$인 점이라 하자.
$A' = \mathcal{O}_{S', s}$, $B' = \mathcal{O}_{X', x}$,
$A = \mathcal{O}_{S, s}$, 그리고 $B = \mathcal{O}_{X, x}$라 놓자.
어떤 제곱이 영인 아이디얼들에 대하여 $A = A'/J$이고 $B = B'/I$이다.
$(f, f')$가 두꺼워짐의 엄격 사상이므로 $I = JB'$이다.
$M' = \mathcal{F}'_x$ 및 $M = \mathcal{F}_x$라 놓자. 그러면
$M'$는 $B'$-가군이고 $M$은 $B$-가군이다.
$\mathcal{F} = i^*\mathcal{F}'$이므로 전사 준동형 $M' \to M$의 핵은
$IM' = JM'$이다. 따라서 다음 짧은 완전열이 있다.
$$
0 \to JM' \to M' \to M \to 0
$$
층, 보조정리 \ref{sheaves-lemma-stalk-pullback-modules}와
가군, 보조정리 \ref{modules-lemma-stalk-tensor-product}를 사용하여
당김과 텐서곱의 줄기들을 식별하면, 보조정리에 나오는 표준 준동형의
$x$에서의 줄기는 다음 준동형임을 알 수 있다.
$$
(J \otimes_A B) \otimes_B M = J \otimes_A M = J \otimes_{A'} M'
\longrightarrow JM'
$$
$\mathcal{F}$가 $S$ 위에서 평탄하다는 가정은 $M$이 평탄한
$A$-가군이라는 뜻이다.

\medskip\noindent
(1)을 가정하자. 대수학, 보조정리
\ref{algebra-lemma-characterize-flat}에 의하면 평탄성으로부터
$\text{Tor}_1^{A'}(M', A) = 0$을 얻는다. 대수학, 주
\ref{algebra-remark-Tor-ring-mod-ideal}에 의하면 이는
$J \otimes_{A'} M' \to M'$가 단사라는 뜻이다. 따라서
$J \otimes_A M \to JM'$은 동형사상이다.

\medskip\noindent
(2)를 가정하자. 그러면 $J \otimes_{A'} M' \to M'$는 단사이다.
따라서 대수학, 주 \ref{algebra-remark-Tor-ring-mod-ideal}에 의해
$\text{Tor}_1^{A'}(M', A) = 0$이다. 그러므로 대수학, 보조정리
\ref{algebra-lemma-what-does-it-mean}에 의해 $M'$는 $A'$ 위에서
평탄하다.
\end{proof}

\begin{lemma}
\label{lemma-inf-map-rel}
$(f, f')$를 상황 \ref{situation-morphism-thickenings}에서와 같은
1차 두꺼워짐 사이의 사상이라 하자. $\mathcal{F}'$, $\mathcal{G}'$를
$\mathcal{O}_{X'}$-가군이라 하고 $\mathcal{F} = i^*\mathcal{F}'$ 및
$\mathcal{G} = i^*\mathcal{G}'$라 놓자.
$\varphi : \mathcal{F} \to \mathcal{G}$를 $\mathcal{O}_X$-선형
준동형이라 하자. $\mathcal{G}'$가 $S'$ 위에서 평탄하고
$(f, f')$가 두꺼워짐의 엄격 사상이라고 가정하자. $\varphi$를
$\mathcal{O}_{X'}$-선형 준동형
$\varphi' : \mathcal{F}' \to \mathcal{G}'$로 올리는 방법들의 집합은
공집합이 아니면 다음 군 아래의 주동차 공간이다.
$$
\Hom_{\mathcal{O}_X}(\mathcal{F},
\mathcal{G} \otimes_{\mathcal{O}_X} f^*\mathcal{J})
$$
\end{lemma}

\begin{proof}
보조정리 \ref{lemma-inf-map-special}과
\ref{lemma-deform-module}을 결합하면 된다.
\end{proof}

\begin{lemma}
\label{lemma-inf-obs-map-special}
$i : (X, \mathcal{O}_X) \to (X', \mathcal{O}_{X'})$를 환 달린 공간의
1차 두꺼워짐이라 하자. $\mathcal{F}'$, $\mathcal{G}'$를
$\mathcal{O}_{X'}$-가군이라 하고 $\mathcal{F} = i^*\mathcal{F}'$ 및
$\mathcal{G} = i^*\mathcal{G}'$라 놓자.
$\varphi : \mathcal{F} \to \mathcal{G}$를 $\mathcal{O}_X$-선형
준동형이라 하자. 다음 원소가 존재한다.
$$
o(\varphi) \in
\Ext^1_{\mathcal{O}_X}(Li^*\mathcal{F}',
\mathcal{I}\mathcal{G}')
$$
이 원소가 영인 것은 $\varphi$를 $\mathcal{O}_{X'}$-선형 준동형
$\varphi' : \mathcal{F}' \to \mathcal{G}'$로 올릴 수 있기 위한
필요충분조건이다.
\end{lemma}

\begin{proof}
보조정리 \ref{lemma-inf-map-special}의 증명으로부터, 다음 준동형에
따른 $\varphi$의 경계가 영인 것이 올림의 존재를 위한
필요충분조건임을 알 수 있다.
$$
\Hom_{\mathcal{O}_X}(\mathcal{F}, \mathcal{G}) =
\Hom_{\mathcal{O}_{X'}}(\mathcal{F}', \mathcal{G}) \longrightarrow
\Ext^1_{\mathcal{O}_{X'}}(\mathcal{F}', \mathcal{I}\mathcal{G}')
$$
한편
$$
\Ext^1_{\mathcal{O}_{X'}}(\mathcal{F}', \mathcal{I}\mathcal{G}') =
\Ext^1_{\mathcal{O}_X}(Li^*\mathcal{F}', \mathcal{I}\mathcal{G}')
$$
는 유도 범주에서의 $i_* = Ri_*$와 $Li^*$의 수반성에서 따른다
(코호몰로지, 보조정리 \ref{cohomology-lemma-adjoint}).
\end{proof}

\begin{lemma}
\label{lemma-inf-obs-map-rel}
$(f, f')$를 상황 \ref{situation-morphism-thickenings}에서와 같은
1차 두꺼워짐 사이의 사상이라 하자. $\mathcal{F}'$, $\mathcal{G}'$를
$\mathcal{O}_{X'}$-가군이라 하고 $\mathcal{F} = i^*\mathcal{F}'$ 및
$\mathcal{G} = i^*\mathcal{G}'$라 놓자.
$\varphi : \mathcal{F} \to \mathcal{G}$를 $\mathcal{O}_X$-선형
준동형이라 하자. $\mathcal{F}'$와 $\mathcal{G}'$가 $S'$ 위에서
평탄하고 $(f, f')$가 두꺼워짐의 엄격 사상이라고 가정하자.
다음 원소가 존재한다.
$$
o(\varphi) \in  \Ext^1_{\mathcal{O}_X}(\mathcal{F},
\mathcal{G} \otimes_{\mathcal{O}_X} f^*\mathcal{J})
$$
이 원소가 영인 것은 $\varphi$를 $\mathcal{O}_{X'}$-선형 준동형
$\varphi' : \mathcal{F}' \to \mathcal{G}'$로 올릴 수 있기 위한
필요충분조건이다.
\end{lemma}

\begin{proof}[첫째 증명]
보조정리의 가정 아래 다음 등식이 성립한다고 주장하면, 결과는
보조정리 \ref{lemma-inf-obs-map-special}에서 따른다.
$$
\Ext^1_{\mathcal{O}_X}(Li^*\mathcal{F}',
\mathcal{I}\mathcal{G}') =
\Ext^1_{\mathcal{O}_X}(\mathcal{F},
\mathcal{G} \otimes_{\mathcal{O}_X} f^*\mathcal{J})
$$
실제로 보조정리 \ref{lemma-deform-module}에 의해
$\mathcal{I}\mathcal{G}' =
\mathcal{G} \otimes_{\mathcal{O}_X} f^*\mathcal{J}$이다. 한편
다음에 유의하자.
$$
H^{-1}(Li^*\mathcal{F}') =
\text{Tor}_1^{\mathcal{O}_{X'}}(\mathcal{F}', \mathcal{O}_X)
$$
(국소 계산은 생략한다). 다음 짧은 완전열을 사용하면
$$
0 \to \mathcal{I} \to \mathcal{O}_{X'} \to \mathcal{O}_X \to 0
$$
이 $\text{Tor}_1$은 준동형
$\mathcal{I} \otimes_{\mathcal{O}_X} \mathcal{F} \to \mathcal{I}\mathcal{F}'$
의 핵으로 계산됨을 알 수 있다. 이 핵은 보조정리
\ref{lemma-deform-module}의 마지막 명제에 의해 영이다. 따라서
$\tau_{\geq -1}Li^*\mathcal{F}' = \mathcal{F}$이다. 한편
$$
\Ext^1_{\mathcal{O}_X}(Li^*\mathcal{F}',
\mathcal{I}\mathcal{G}') =
\Ext^1_{\mathcal{O}_X}(\tau_{\geq -1}Li^*\mathcal{F}',
\mathcal{I}\mathcal{G}')
$$
는 유도 범주, 보조정리
\ref{derived-lemma-negative-vanishing}의 쌍대형에서 따른다.
\end{proof}

\begin{proof}[둘째 증명]
보조정리 \ref{lemma-inf-obs-map}을 다음과 같이 적용할 수 있다.
보조정리 \ref{lemma-deform-module}에 의해
$\mathcal{K} = \mathcal{I} \otimes_{\mathcal{O}_X} \mathcal{F}$이고
$\mathcal{L} = \mathcal{I} \otimes_{\mathcal{O}_X} \mathcal{G}$이며,
$c_{\mathcal{F}'} = 1 \otimes 1$이고 $c_{\mathcal{G}'} = 1 \otimes 1$이다.
$\psi = 1 \otimes \varphi$로 택하면 보조정리의 도표가 가환한다.
따라서 $o(\varphi) = o(\varphi, 1 \otimes \varphi)$로 택하면 된다.
\end{proof}

\begin{lemma}
\label{lemma-inf-ext-rel}
$(f, f')$를 상황 \ref{situation-morphism-thickenings}에서와 같은
1차 두꺼워짐 사이의 사상이라 하자. $\mathcal{F}$를
$\mathcal{O}_X$-가군이라 하자. $(f, f')$가 두꺼워짐의 엄격 사상이고
$\mathcal{F}$가 $S$ 위에서 평탄하다고 가정하자.
$S'$ 위에서 평탄한 $\mathcal{O}_{X'}$-가군 $\mathcal{F}'$와 동형사상
$\alpha : i^*\mathcal{F}' \to \mathcal{F}$로 이루어진 쌍
$(\mathcal{F}', \alpha)$가 하나라도 존재하면, 그러한 쌍들의 동형류
집합은 다음 군 아래의 주동차 공간이다.
$\Ext^1_{\mathcal{O}_X}(
\mathcal{F}, \mathcal{I} \otimes_{\mathcal{O}_X} \mathcal{F})$.
\end{lemma}

\begin{proof}
그러한 가군이 하나 존재한다고 가정하면, 보조정리
\ref{lemma-deform-module}에 의해 표준 준동형
$$
f^*\mathcal{J} \otimes_{\mathcal{O}_X} \mathcal{F} \to
\mathcal{I} \otimes_{\mathcal{O}_X} \mathcal{F}
$$
은 동형사상이다. $\mathcal{K} =
\mathcal{I} \otimes_{\mathcal{O}_X} \mathcal{F}$ 및 $c = 1$로 놓고
보조정리 \ref{lemma-inf-ext}를 적용하자. 보조정리
\ref{lemma-deform-module}에 의해 대응하는 확대 $\mathcal{F}'$는 모두
$S'$ 위에서 평탄하다.
\end{proof}

\begin{lemma}
\label{lemma-inf-obs-ext-rel}
$(f, f')$를 상황 \ref{situation-morphism-thickenings}에서와 같은
1차 두꺼워짐 사이의 사상이라 하자. $\mathcal{F}$를
$\mathcal{O}_X$-가군이라 하자. $(f, f')$가 두꺼워짐의 엄격 사상이고
$\mathcal{F}$가 $S$ 위에서 평탄하다고 가정하자. $S'$ 위에서 평탄하며
$i^*\mathcal{F}' \cong \mathcal{F}$를 만족하는
$\mathcal{O}_{X'}$-가군 $\mathcal{F}'$가 존재하기 위한 필요충분조건은
다음과 같다.
\begin{enumerate}
\item 표준 준동형 $
f^*\mathcal{J} \otimes_{\mathcal{O}_X} \mathcal{F} \to
\mathcal{I} \otimes_{\mathcal{O}_X} \mathcal{F}$
은 동형사상이다.
\item 보조정리 \ref{lemma-inf-obs-ext}의 동치류
$o(\mathcal{F}, \mathcal{I} \otimes_{\mathcal{O}_X} \mathcal{F}, 1)
\in \Ext^2_{\mathcal{O}_X}(
\mathcal{F}, \mathcal{I} \otimes_{\mathcal{O}_X} \mathcal{F})$
는 영이다.
\end{enumerate}
\end{lemma}

\begin{proof}
이는 보조정리 \ref{lemma-deform-module}에 주어진 $S'$ 위에서 평탄한
$\mathcal{O}_{X'}$-가군의 특징짓기와 보조정리
\ref{lemma-inf-obs-ext}에서 바로 따른다.
\end{proof}






\section{환 달린 공간의 평탄한 두꺼워짐 위 평탄한 가군에 대한 응용}
\label{section-flat}

\noindent
다음 환 달린 공간의 가환 도표를 생각하자.
$$
\xymatrix{
(X, \mathcal{O}_X) \ar[r]_i \ar[d]_f & (X', \mathcal{O}_{X'}) \ar[d]^{f'} \\
(S, \mathcal{O}_S) \ar[r]^t & (S', \mathcal{O}_{S'})
}
$$
가로 화살표들은 상황 \ref{situation-morphism-thickenings}에서와 같은
1차 두꺼워짐이다. $\mathcal{I} = \Ker(i^\sharp) \subset \mathcal{O}_{X'}$
및 $\mathcal{J} = \Ker(t^\sharp) \subset \mathcal{O}_{S'}$라 놓자.
$\mathcal{F}$를 $\mathcal{O}_X$-가군이라 하고 다음을 가정하자.
\begin{enumerate}
\item $(f, f')$는 두꺼워짐의 엄격 사상이다.
\item $f'$는 평탄하다.
\item $\mathcal{F}$는 $S$ 위에서 평탄하다.
\end{enumerate}
(1) $+$ (2)로부터 $\mathcal{I} = f^*\mathcal{J}$가 따름에 유의하자
(보조정리 \ref{lemma-deform-module}을 $\mathcal{O}_{X'}$에 적용한다).
이 가정들 아래에서는 앞 절의 이론이 특히 간결해진다. 지금까지 얻은
결과를 다음 보조정리로 요약한다.

\begin{lemma}
\label{lemma-flat}
위 상황에서 다음이 성립한다.
\begin{enumerate}
\item $S'$ 위에서 평탄하고 $i^*\mathcal{F}' \cong \mathcal{F}$를
만족하는 $\mathcal{O}_{X'}$-가군 $\mathcal{F}'$가 존재하기 위한
필요충분조건은 보조정리 \ref{lemma-inf-obs-ext}의 동치류
$o(\mathcal{F}, f^*\mathcal{J} \otimes_{\mathcal{O}_X} \mathcal{F}, 1)
\in \Ext^2_{\mathcal{O}_X}(
\mathcal{F}, f^*\mathcal{J} \otimes_{\mathcal{O}_X} \mathcal{F})$
가 영인 것이다.
\item 그러한 가군이 존재하면 올림들의 동형류 집합은 다음 군 아래의
주동차 공간이다.
$\Ext^1_{\mathcal{O}_X}(
\mathcal{F}, f^*\mathcal{J} \otimes_{\mathcal{O}_X} \mathcal{F})$.
\item 올림 $\mathcal{F}'$가 주어졌을 때, 당기면
$\text{id}_\mathcal{F}$가 되는 $\mathcal{F}'$의 자기동형사상들의
집합은 다음 군과 표준적으로 동형이다.
$\Ext^0_{\mathcal{O}_X}(
\mathcal{F}, f^*\mathcal{J} \otimes_{\mathcal{O}_X} \mathcal{F})$.
\end{enumerate}
\end{lemma}

\begin{proof}
위에서 $\mathcal{I} = f^*\mathcal{J}$임을 보았으므로 (1)은 보조정리
\ref{lemma-inf-obs-ext-rel}에서 따른다. (2)는 보조정리
\ref{lemma-inf-ext-rel}에서, (3)은 보조정리
\ref{lemma-inf-map-rel}에서 따른다.
\end{proof}

\begin{situation}
\label{situation-ses-flat-thickenings}
$f : (X, \mathcal{O}_X) \to (S, \mathcal{O}_S)$를 환 달린 공간의
사상이라 하자. 다음 가환 도표를 생각하자.
$$
\xymatrix{
(X'_1, \mathcal{O}'_1) \ar[r]_h \ar[d]_{f'_1} &
(X'_2, \mathcal{O}'_2) \ar[r] \ar[d]_{f'_2} &
(X'_3, \mathcal{O}'_3) \ar[d]_{f'_3} \\
(S'_1, \mathcal{O}_{S'_1}) \ar[r] &
(S'_2, \mathcal{O}_{S'_2}) \ar[r] &
(S'_3, \mathcal{O}_{S'_3})
}
$$
여기서 (a) 윗행은 $X$의 1차 두꺼워짐들의 짧은 완전열이고,
(b) 아랫행은 $S$의 1차 두꺼워짐들의 짧은 완전열이며, (c) 각
$f'_i$의 제한은 $f$이고, (d) 각 쌍 $(f, f_i')$는 두꺼워짐의 엄격
사상이며, (e) 각 $f'_i$는 평탄하다. 끝으로 $\mathcal{F}'_2$를
$S'_2$ 위에서 평탄한 $\mathcal{O}'_2$-가군이라 하고
$\mathcal{F} = \mathcal{F}'_2|_X$라 놓자.
$\pi : X'_1 \to X$를 표준 분할이라 하자
(주 \ref{remark-short-exact-sequence-thickenings}).
\end{situation}

\begin{lemma}
\label{lemma-verify-iv}
상황 \ref{situation-ses-flat-thickenings}에서 가군
$\pi^*\mathcal{F}$와 $h^*\mathcal{F}'_2$는 $S'_1$ 위에서 평탄하고
$X$ 위로 제한하면 $\mathcal{F}$가 되는 $\mathcal{O}'_1$-가군이다.
이 둘의 차이(보조정리 \ref{lemma-flat})는
$\Ext^1_{\mathcal{O}_X}(
\mathcal{F}, f^*\mathcal{J}_1 \otimes_{\mathcal{O}_X} \mathcal{F})$
의 원소 $\theta$이고, 그 경계는
$\Ext^2_{\mathcal{O}_X}(
\mathcal{F}, f^*\mathcal{J}_3 \otimes_{\mathcal{O}_X} \mathcal{F})$
에서 $\mathcal{F}$를 $S'_3$ 위에서 평탄한 $\mathcal{O}'_3$-가군으로
올리는 데 대한 장애 원소(보조정리 \ref{lemma-flat})와 같다.
\end{lemma}

\begin{proof}
$\pi^*\mathcal{F}$와 $h^*\mathcal{F}'_2$는 모두 $X$ 위로 제한하면
$\mathcal{F}$가 되고, $\pi^*\mathcal{F} \to \mathcal{F}$와
$h^*\mathcal{F}'_2 \to \mathcal{F}$의 핵은
$f^*\mathcal{J}_1 \otimes_{\mathcal{O}_X} \mathcal{F}$로 주어진다.
따라서 보조정리 \ref{lemma-deform-module}에 의해 평탄하다.
다음 가군의 열이 짧은 완전열이므로 경계를 취할 수 있다.
$$
0 \to f^*\mathcal{J}_3 \otimes_{\mathcal{O}_X} \mathcal{F} \to
f^*\mathcal{J}_2 \otimes_{\mathcal{O}_X} \mathcal{F} \to
f^*\mathcal{J}_1 \otimes_{\mathcal{O}_X} \mathcal{F} \to 0
$$
이는 상황 \ref{situation-ses-flat-thickenings}의 가정들과
$\mathcal{F}$가 $S$ 위에서 평탄하다는 사실에서 따른다. 장애 동치류에
대한 명제는 주 \ref{remark-complex-thickenings-and-ses-modules}의
결과를 이 특수한 상황에 그대로 옮긴 것이다.
\end{proof}






\section{환 달린 공간의 변형과 소박한 코탄젠트 복합체}
\label{section-deformations-ringed-spaces}

\noindent
이 절에서는 소박한 코탄젠트 복합체를 사용하여 변형 이론의 일부를
다룬다. 환 달린 공간의 1차 두꺼워짐
$t : (S, \mathcal{O}_S) \to (S', \mathcal{O}_{S'})$에서 시작하자.
$\mathcal{J} = \Ker(t^\sharp)$라 쓰고 $S$와 $S'$의 바탕 위상공간을
동일시한다. 또한 환 달린 공간의 사상
$f : (X, \mathcal{O}_X) \to (S, \mathcal{O}_S)$, $\mathcal{O}_X$-가군
$\mathcal{G}$, 그리고 가군층의 $f$-사상
$c : \mathcal{J} \to \mathcal{G}$가 주어졌다고 가정하자
(층, 정의 \ref{sheaves-definition-f-map} 및 절
\ref{sheaves-section-ringed-spaces-functoriality-modules}).
이 절에서는 다음 도표의 물음표를 채울 수 있는지를 묻는다.
\begin{equation}
\label{equation-to-solve-ringed-spaces}
\vcenter{
\xymatrix{
0 \ar[r] & \mathcal{G} \ar[r] & {?} \ar[r] & \mathcal{O}_X \ar[r] & 0 \\
0 \ar[r] & \mathcal{J} \ar[u]^c \ar[r] & \mathcal{O}_{S'} \ar[u] \ar[r] &
\mathcal{O}_S \ar[u] \ar[r] & 0
}
}
\end{equation}
(세로 화살표들은 $f$-사상이다.) 또한 해가 존재할 때 얼마나
유일한지도 묻는다. 더 정확히는, 1차 두꺼워짐
$i : (X, \mathcal{O}_X) \to (X', \mathcal{O}_{X'})$과
(\ref{equation-morphism-thickenings})에서와 같은 두꺼워짐의 사상
$(f, f')$를 찾는다. 여기서 $\Ker(i^\sharp)$를 $\mathcal{G}$와
동일시하며 $(f')^\sharp$가 주어진 사상 $c$를 유도해야 한다.
이때 $X'$를 (\ref{equation-to-solve-ringed-spaces})의 {\it 해}라고
부른다.

\begin{lemma}
\label{lemma-huge-diagram-ringed-spaces}
다음 환 달린 공간의 사상들의 가환 도표가 주어졌다고 가정하자.
\begin{equation}
\label{equation-huge-1}
\vcenter{
\xymatrix{
& (X_2, \mathcal{O}_{X_2}) \ar[r]_{i_2} \ar[d]_{f_2} \ar[ddl]_g &
(X'_2, \mathcal{O}_{X'_2}) \ar[d]^{f'_2} \\
& (S_2, \mathcal{O}_{S_2}) \ar[r]^{t_2} \ar[ddl]|\hole &
(S'_2, \mathcal{O}_{S'_2}) \ar[ddl] \\
(X_1, \mathcal{O}_{X_1}) \ar[r]_{i_1} \ar[d]_{f_1} &
(X'_1, \mathcal{O}_{X'_1}) \ar[d]^{f'_1} \\
(S_1, \mathcal{O}_{S_1}) \ar[r]^{t_1} &
(S'_1, \mathcal{O}_{S'_1})
}
}
\end{equation}
가로 화살표들은 1차 두꺼워짐이다. $\mathcal{G}_j = \Ker(i_j^\sharp)$라
놓고, 다음 가환 도표를 유도하는 가군의 $g$-사상
$\nu : \mathcal{G}_1 \to \mathcal{G}_2$가 주어졌다고 가정하자.
\begin{equation}
\label{equation-huge-2}
\vcenter{
\xymatrix{
& 0 \ar[r] & \mathcal{G}_2 \ar[r] &
\mathcal{O}_{X'_2} \ar[r] &
\mathcal{O}_{X_2} \ar[r] & 0 \\
& 0 \ar[r]|\hole &
\mathcal{J}_2 \ar[u]_{c_2} \ar[r] &
\mathcal{O}_{S'_2} \ar[u] \ar[r]|\hole &
\mathcal{O}_{S_2} \ar[u] \ar[r] & 0 \\
0 \ar[r] & \mathcal{G}_1 \ar[ruu] \ar[r] &
\mathcal{O}_{X'_1} \ar[r] &
\mathcal{O}_{X_1} \ar[ruu] \ar[r] & 0 \\
0 \ar[r] & \mathcal{J}_1 \ar[ruu]|\hole \ar[u]^{c_1} \ar[r] &
\mathcal{O}_{S'_1} \ar[ruu]|\hole \ar[u] \ar[r] &
\mathcal{O}_{S_1} \ar[ruu]|\hole \ar[u] \ar[r] & 0
}
}
\end{equation}
이 도표의 앞면과 뒷면은 (\ref{equation-to-solve-ringed-spaces})의
해이다.
\begin{enumerate}
\item 다음 군에는 표준 원소가 존재한다.
$\Ext^1_{\mathcal{O}_{X_2}}(Lg^*\NL_{X_1/S_1}, \mathcal{G}_2)$
이 원소가 영인 것은 $\nu$와 양립하며
(\ref{equation-huge-1})에 들어맞는 환 달린 공간의 사상
$X'_2 \to X'_1$이 존재하기 위한 필요충분조건이다.
\item $\nu$와 양립하며 (\ref{equation-huge-1})에 들어맞는 사상
$X'_2 \to X'_1$이 존재하면, 그러한 모든 사상의 집합은 다음 군
아래의 주동차 공간이다.
$$
\Hom_{\mathcal{O}_{X_1}}(\Omega_{X_1/S_1}, g_*\mathcal{G}_2) =
\Hom_{\mathcal{O}_{X_2}}(g^*\Omega_{X_1/S_1}, \mathcal{G}_2) =
\Ext^0_{\mathcal{O}_{X_2}}(Lg^*\NL_{X_1/S_1}, \mathcal{G}_2).
$$
\end{enumerate}
\end{lemma}

\begin{proof}
소박한 코탄젠트 복합체 $\NL_{X_1/S_1}$은 가군, 정의
\ref{modules-definition-cotangent-complex-morphism-ringed-topoi}에서
정의한다. 보조정리의 마지막 명제에 나오는 등식들은 $g^*$가
$g_*$의 왼쪽 수반이라는 사실,
$H^0(\NL_{X_1/S_1}) = \Omega_{X_1/S_1}$이라는 사실(소박한 코탄젠트
복합체의 구성에 의한다), 그리고 $Lg^*$가 $g^*$의 왼쪽 유도
함자라는 사실에서 따른다. 따라서 증명의 나머지에서는 다음 군들을
사용한다.
$\Ext^k_{\mathcal{O}_{X_2}}(Lg^*\NL_{X_1/S_1}, \mathcal{G}_2)$,
$k = 0, 1$. 먼저 보조정리에 나오는 모든 환 달린 공간의 바탕
위상공간이 같은 경우로 환원할 수 있음을 보이자.

\medskip\noindent
이를 위해 $g^{-1}\NL_{X_1/S_1}$이 환층의 준동형
$g^{-1}f_1^{-1}\mathcal{O}_{S_1} \to g^{-1}\mathcal{O}_{X_1}$의
소박한 코탄젠트 복합체와 같음에 유의하자. 가군, 보조정리
\ref{modules-lemma-pullback-NL}를 보라. 더욱이
$\NL_{X_1/S_1}$의 차수 $0$인 항은 평탄한 $\mathcal{O}_{X_1}$-가군이다.
따라서 표준 준동형
$$
Lg^*\NL_{X_1/S_1}
\longrightarrow
g^{-1}\NL_{X_1/S_1} \otimes_{g^{-1}\mathcal{O}_{X_1}} \mathcal{O}_{X_2}
$$
은 차수 $0$과 $-1$의 코호몰로지층에서 동형사상을 유도한다.
따라서 보조정리의 Ext 군들을 다음 군들로 바꾸어도 된다.
$$
\Ext^k_{g^{-1}\mathcal{O}_{X_1}}(g^{-1}\NL_{X_1/S_1}, \mathcal{G}_2) =
\Ext^k_{g^{-1}\mathcal{O}_{X_1}}(
\NL_{g^{-1}\mathcal{O}_{X_1}/g^{-1}f_1^{-1}\mathcal{O}_{S_1}}, \mathcal{G}_2)
$$
(\ref{equation-huge-1})에 들어맞고 $\nu$와 양립하는 환 달린 공간의
사상 $X'_2 \to X'_1$들의 집합은 $f^\sharp$ 및 $\nu$와 양립하는
$g^{-1}f_1^{-1}\mathcal{O}_{S'_1}$-대수 준동형
$g^{-1}\mathcal{O}_{X'_1} \to \mathcal{O}_{X'_2}$들의 집합과
일대일 대응한다. 따라서 $X$ 위의 층들로 이루어진 도표
(\ref{equation-huge-2})가 주어졌고, 그 도표에 들어맞는 환층의 준동형
$\mathcal{O}_{X'_1} \to \mathcal{O}_{X'_2}$를 찾는 경우를 가정해도
된다.

\medskip\noindent
보조정리의 증명의 나머지에서는 모든 바탕 위상공간이 같다고
가정한다. 즉 공간 $X$ 위의 층들로 이루어진 도표
(\ref{equation-huge-2})가 주어졌고, 그 도표에 들어맞는 환층의 준동형
$\mathcal{O}_{X'_1} \to \mathcal{O}_{X'_2}$를 찾는다. Ext 군으로는
$\Ext^k_{\mathcal{O}_{X_1}}(
\NL_{\mathcal{O}_{X_1}/\mathcal{O}_{S_1}}, \mathcal{G}_2)$, $k = 0, 1$.

\medskip\noindent
1단계. 장애 동치류의 구성. 다음 집합층을 생각하자.
$$
\mathcal{E} = \mathcal{O}_{X'_1} \times_{\mathcal{O}_{X_2}} \mathcal{O}_{X'_2}
$$
여기에는 전사 사상 $\alpha : \mathcal{E} \to \mathcal{O}_{X_1}$이
따라오므로 $\NL_{\mathcal{O}_{X_1}/\mathcal{O}_{S_1}}$ 대신
$\NL(\alpha)$를 사용할 수 있다. 가군, 보조정리
\ref{modules-lemma-NL-up-to-qis}를 보라. 다음과 같이 놓자.
$$
\mathcal{I}' =
\Ker(\mathcal{O}_{S'_1}[\mathcal{E}] \to \mathcal{O}_{X_1})
\quad\text{그리고}\quad
\mathcal{I} =
\Ker(\mathcal{O}_{S_1}[\mathcal{E}] \to \mathcal{O}_{X_1})
$$
핵이 $\mathcal{J}_1\mathcal{O}_{S'_1}[\mathcal{E}]$인 전사 준동형
$\mathcal{I}' \to \mathcal{I}$이 있다. 두 개의
$\mathcal{O}_{S'_1}$-대수 준동형
$$
a : \mathcal{O}_{S'_1}[\mathcal{E}] \to \mathcal{O}_{X'_1}
\quad\text{그리고}\quad
b : \mathcal{O}_{S'_1}[\mathcal{E}] \to \mathcal{O}_{X'_2}
$$
를 얻는다. 이들은 각각 준동형
$a|_{\mathcal{I}'} : \mathcal{I}' \to \mathcal{G}_1$과
$b|_{\mathcal{I}'} : \mathcal{I}' \to \mathcal{G}_2$를 유도한다.
$a$와 $b$는 모두 $(\mathcal{I}')^2$을 영으로 보낸다. 더욱이
(\ref{equation-huge-2})의 왼쪽 정사각형이 가환하므로
$\mathcal{G}_2$로 가는 사상으로서 $a$와 $b$는
$\mathcal{J}_1\mathcal{O}_{S'_1}[\mathcal{E}]$ 위에서 일치한다.
따라서 차
$b|_{\mathcal{I}'} - \nu \circ a|_{\mathcal{I}'}$
는 잘 정의된 $\mathcal{O}_{X_1}$-선형 준동형
$$
\xi : \mathcal{I}/\mathcal{I}^2 \longrightarrow \mathcal{G}_2
$$
를 유도한다. 이 준동형은 $\mathcal{I}$의 국소절단 $f$의 동치류를
$\nu(a(f')) - b(f')$로 보내는데, 여기서 $f'$는 $f$를
$\mathcal{I}'$의 국소절단으로 올린 것이다. 그 상을
$[\xi] \in \Ext^1_{\mathcal{O}_{X_1}}(\NL(\alpha), \mathcal{G}_2)$
라 하자(아래를 보라).

\medskip\noindent
2단계. $[\xi]$의 소멸은 필요하다. 다음과 같이 쓰자.
$\Omega = \Omega_{\mathcal{O}_{S_1}[\mathcal{E}]/\mathcal{O}_{S_1}}
\otimes_{\mathcal{O}_{S_1}[\mathcal{E}]} \mathcal{O}_{X_1}$.
$\NL(\alpha) = (\mathcal{I}/\mathcal{I}^2 \to \Omega)$가 다음 완전
삼각형에 들어맞음에 유의하자.
$$
\Omega[0] \to
\NL(\alpha) \to
\mathcal{I}/\mathcal{I}^2[1] \to
\Omega[1]
$$
따라서 어떤 준동형 $\Omega \to \mathcal{G}_2$에 대하여 $\xi$가 합성
$\mathcal{I}/\mathcal{I}^2 \to \Omega \to \mathcal{G}_2$일 때, 그리고
그때에만 $[\xi]$가 영이다. (\ref{equation-huge-2})에 들어맞는 환층의
준동형 $\varphi : \mathcal{O}_{X'_1} \to \mathcal{O}_{X'_2}$가
존재한다고 하자. 이 경우 준동형
$\mathcal{O}_{S'_1}[\mathcal{E}] \to \mathcal{G}_2$,
$f' \mapsto b(f') - \varphi(a(f'))$를 생각하자. 계산하면 이 준동형은
$\mathcal{J}_1\mathcal{O}_{S'_1}[\mathcal{E}]$을 영으로 보내며 미분
$\mathcal{O}_{S_1}[\mathcal{E}] \to \mathcal{G}_2$를 유도함을 알 수
있다. 그 결과 얻는 선형 준동형 $\Omega \to \mathcal{G}_2$는 이 경우
$[\xi] = 0$임을 보인다.

\medskip\noindent
3단계. $[\xi]$의 소멸은 충분하다.
$\theta : \Omega \to \mathcal{G}_2$를 $\mathcal{O}_{X_1}$-선형
준동형이라 하자. 이때 $\xi$는
$\theta \circ (\mathcal{I}/\mathcal{I}^2 \to \Omega)$와 같다고
가정한다. 계산하면
$$
b + \theta \circ d : \mathcal{O}_{S'_1}[\mathcal{E}] \to \mathcal{O}_{X'_2}
$$
의 $\mathcal{I}'$ 위 제한은
$\nu \circ a : \mathcal{I}' \to \mathcal{G}_2$와 일치한다.
$\mathcal{O}_{X'_1}$은 $\mathcal{I}' \to
\mathcal{O}_{S'_1}[\mathcal{E}]$과 $\mathcal{I}' \to \mathcal{G}_1$의
밀어내기이므로, 준동형 $b + \theta \circ d$와 $a$는
(\ref{equation-huge-2})에 들어맞는 준동형
$\mathcal{O}_{X'_1} \to \mathcal{O}_{X'_2}$를 정한다.

\medskip\noindent
위의 특수한 경우에서 (2)의 증명. 생략한다. 도움말: 이는 보조정리
\ref{lemma-huge-diagram}의 (2)의 증명과 완전히 같다.
\end{proof}

\begin{lemma}
\label{lemma-NL-represent-ext-class}
$X$를 위상공간이라 하자. $\mathcal{A} \to \mathcal{B}$를 환층의
준동형이라 하고 $\mathcal{G}$를 $\mathcal{B}$-가군이라 하자.
$\xi \in \Ext^1_\mathcal{B}(\NL_{\mathcal{B}/\mathcal{A}}, \mathcal{G})$. 
라 하자. 집합층의 사상 $\alpha : \mathcal{E} \to \mathcal{B}$를
적절히 택하면, $\xi \in \Ext^1_\mathcal{B}(\NL(\alpha), \mathcal{G})$는
준동형 $\mathcal{I}/\mathcal{I}^2 \to \mathcal{G}$의 동치류로
표현된다(표기는 증명을 보라).
\end{lemma}

\begin{proof}
사상 $\alpha : \mathcal{E} \to \mathcal{B}$가 주어지고
$\mathcal{A}[\mathcal{E}] \to \mathcal{B}$가 핵이 $\mathcal{I}$인
전사 준동형이면, 복합체
$\NL(\alpha) = (\mathcal{I}/\mathcal{I}^2 \to 
\Omega_{\mathcal{A}[\mathcal{E}]/\mathcal{A}}
\otimes_{\mathcal{A}[\mathcal{E}]} \mathcal{B})$는
$\NL_{\mathcal{B}/\mathcal{A}}$와 표준적으로 동형임을 상기하자.
가군, 보조정리 \ref{modules-lemma-NL-up-to-qis}를 보라. 더욱이
$\Omega = \Omega_{\mathcal{A}[\mathcal{E}]/\mathcal{A}}
\otimes_{\mathcal{A}[\mathcal{E}]} \mathcal{B}$는 다음 준층에 연관된
층이다.
$U \mapsto \bigoplus_{e \in \mathcal{E}(U)} \mathcal{B}(U)$.
다시 말해 $\Omega$는 집합층 $\mathcal{E}$ 위의 자유
$\mathcal{B}$-가군이고, 특히 표준 사상 $\mathcal{E} \to \Omega$가
있다.

\medskip\noindent
이제 어떤 $\mathcal{E}$를 택하자(예를 들어 소박한 코탄젠트 복합체의
정의에서처럼 $\mathcal{E} = \mathcal{B}$로 택한다). $\xi$를
준동형 $\mathcal{I}/\mathcal{I}^2 \to \mathcal{G}$의 동치류로
표현하는 데 대한 장애 원소는
$\Ext^1_\mathcal{B}(\Omega, \mathcal{G})$에 속한다. 이 원소가
$\mathcal{B}$-가군의 확대
$0 \to \mathcal{G} \to \mathcal{H} \to \Omega \to 0$로 표현된다고
하자. 유도된 사상 $\alpha' : \mathcal{E}' \to \mathcal{B}$를 갖는
집합층 $\mathcal{E}' = \mathcal{E} \times_\Omega \mathcal{H}$를
생각하자. $\mathcal{I}' =
\Ker(\mathcal{A}[\mathcal{E}'] \to \mathcal{B})$ 및
$\Omega' = \Omega_{\mathcal{A}[\mathcal{E}']/\mathcal{A}}
\otimes_{\mathcal{A}[\mathcal{E}']} \mathcal{B}$라 놓자.
준동형 $\Omega' \to \Omega$에 따라 확대 $\mathcal{H}$를 당기면
분할된다. 실제로 $\Omega'$은 집합층 $\mathcal{E}'$ 위의 자유
$\mathcal{B}$-가군이고, 구성상 다음 가환 도표가 있기 때문이다.
$$
\xymatrix{
\mathcal{E}' \ar[r] \ar[d] & \mathcal{E} \ar[d] \\
\mathcal{H} \ar[r] & \Omega
}
$$
따라서 준동형 $\NL(\alpha') \to \NL(\alpha)$에 따른 $\xi$의 당김은
$\Ext^1_\mathcal{B}(\Omega', \mathcal{G})$에서 영으로 간다. 이로써
증명이 끝난다.
\end{proof}

\begin{lemma}
\label{lemma-choices-ringed-spaces}
(\ref{equation-to-solve-ringed-spaces})의 해가 존재하면, 해들의 동형류
집합은 $\Ext^1_{\mathcal{O}_X}(\NL_{X/S}, \mathcal{G})$ 아래의
주동차 공간이다.
\end{lemma}

\begin{proof}
(\ref{equation-to-solve-ringed-spaces})의 두 해 $X'_1$과 $X'_2$가
주어지면, 보조정리 \ref{lemma-huge-diagram-ringed-spaces}에 의해
사상 $X'_1 \to X'_2$의 존재에 대한 장애 원소
$o(X'_1, X'_2) \in \Ext^1_{\mathcal{O}_X}(\NL_{X/S}, \mathcal{G})$
를 곧바로 얻는다. 이 원소는 동형사상의 존재에 대한 장애 원소이므로
동형류들을 구별한다. 따라서 증명을 끝내려면 해 $X'$와 원소
$\xi \in \Ext^1_{\mathcal{O}_X}(\NL_{X/S}, \mathcal{G})$
가 주어졌을 때 $o(X', X'_\xi) = \xi$를 만족하는 두 번째 해
$X'_\xi$를 찾을 수 있음을 보이면 충분하다.

\medskip\noindent
동치류 $\xi$에 대하여 보조정리
\ref{lemma-NL-represent-ext-class}에서와 같이
$\alpha : \mathcal{E} \to \mathcal{O}_X$를 택하자. 핵이
$\mathcal{I}$인 전사 준동형
$f^{-1}\mathcal{O}_S[\mathcal{E}] \to \mathcal{O}_X$과 그에 대응하는
소박한 코탄젠트 복합체
$\NL(\alpha) = (\mathcal{I}/\mathcal{I}^2 \to
\Omega_{f^{-1}\mathcal{O}_S[\mathcal{E}]/f^{-1}\mathcal{O}_S}
\otimes_{f^{-1}\mathcal{O}_S[\mathcal{E}]} \mathcal{O}_X)$를 생각하자.
이 보조정리에 의해 $\xi$는 준동형
$\delta : \mathcal{I}/\mathcal{I}^2 \to \mathcal{G}$의 동치류이다.
$\mathcal{E}$를
$\mathcal{E} \times_{\mathcal{O}_X} \mathcal{O}_{X'}$로 바꾸면,
$\alpha$가 사상 $\alpha' : \mathcal{E} \to \mathcal{O}_{X'}$를 거쳐
분해된다고 가정할 수도 있다.

\medskip\noindent
이 선택들은 $f^{-1}\mathcal{O}_{S'}$-대수 준동형
$\varphi : f^{-1}\mathcal{O}_{S'}[\mathcal{E}] \to \mathcal{O}_{X'}$를
정한다. $\mathcal{I}' =
\Ker(f^{-1}\mathcal{O}_{S'}[\mathcal{E}] \to \mathcal{O}_X)$라 놓자.
$\varphi$는 준동형
$\varphi|_{\mathcal{I}'} : \mathcal{I}' \to \mathcal{G}$
를 유도하고 $\mathcal{O}_{X'}$은 다음 도표에서와 같은 밀어내기이다.
$$
\xymatrix{
0 \ar[r] & \mathcal{G} \ar[r] & \mathcal{O}_{X'} \ar[r] &
\mathcal{O}_X \ar[r] & 0 \\
0 \ar[r] & \mathcal{I}' \ar[u]^{\varphi|_{\mathcal{I}'}} \ar[r] &
f^{-1}\mathcal{O}_{S'}[\mathcal{E}] \ar[u] \ar[r] &
\mathcal{O}_X \ar[u]_{=} \ar[r] & 0
}
$$
$\psi : \mathcal{I}' \to \mathcal{G}$를
$\varphi|_{\mathcal{I}'}$와 다음 합성의 합이라 하자.
$$
\mathcal{I}' \to \mathcal{I}'/(\mathcal{I}')^2 \to
\mathcal{I}/\mathcal{I}^2 \xrightarrow{\delta} \mathcal{G}.
$$
그러면 $\psi$에 따른 밀어내기는 위와 같은 도표에 들어맞는 또 다른
환의 확대 $\mathcal{O}_{X'_\xi}$이다. 계산하면(생략한다)
$o(X', X'_\xi) = \xi$임을 알 수 있다.
\end{proof}

\begin{lemma}
\label{lemma-extensions-of-relative-ringed-spaces}
$f : (X, \mathcal{O}_X) \to (S, \mathcal{O}_S)$를 환 달린 공간의
사상이라 하고 $\mathcal{G}$를 $\mathcal{O}_X$-가군이라 하자.
다음 $f^{-1}\mathcal{O}_S$-대수의 확대들의 동형류 집합을 생각하자.
$$
0 \to \mathcal{G} \to \mathcal{O}_{X'} \to \mathcal{O}_X \to 0
$$
여기서 $\mathcal{G}$는 제곱이 영인 아이디얼이다.\footnote{다시 말해,
$\mathcal{O}_X$-가군의 동형사상
$\mathcal{G} \to \Ker(i^\sharp)$가 주어진 $S$ 위의 1차 두꺼워짐
$i : X \to X'$들의 동형류 집합이다.}
이 집합은 다음 군과 표준적으로 전단사이다.
$\Ext^1_{\mathcal{O}_X}(\NL_{X/S}, \mathcal{G})$.
\end{lemma}

\begin{proof}
이를 증명하기 위해, (\ref{equation-to-solve-ringed-spaces})이 다음
도표로 주어지는 경우에 앞의 결과들을 적용하자.
$$
\xymatrix{
0 \ar[r] &
\mathcal{G} \ar[r] &
{?} \ar[r] &
\mathcal{O}_X \ar[r] & 0 \\
0 \ar[r] &
0 \ar[u] \ar[r] &
\mathcal{O}_S \ar[u] \ar[r]^{\text{id}} &
\mathcal{O}_S \ar[u] \ar[r] & 0
}
$$
해 $\mathcal{G} \oplus \mathcal{O}_X$가 존재하므로, 보조정리는
보조정리 \ref{lemma-choices-ringed-spaces}에서 따른다. 전단사의 직접
구성은 아래의 주를 보라.
\end{proof}

\begin{remark}
\label{remark-extensions-of-relative-ringed-spaces}
$f : (X, \mathcal{O}_X) \to (S, \mathcal{O}_S)$와 $\mathcal{G}$를
보조정리 \ref{lemma-extensions-of-relative-ringed-spaces}에서와 같이
택하자. 보조정리에서와 같은 확대
$0 \to \mathcal{G} \to \mathcal{O}_{X'} \to \mathcal{O}_X \to 0$
를 생각하자. 집합층 $\mathcal{E}$와 다음 가환 도표를 택할 수 있다.
$$
\xymatrix{
\mathcal{E} \ar[d]_{\alpha'} \ar[rd]^\alpha \\
\mathcal{O}_{X'} \ar[r] & \mathcal{O}_X
}
$$
여기서 $f^{-1}\mathcal{O}_S[\mathcal{E}] \to \mathcal{O}_X$은 핵이
$\mathcal{J}$인 전사 준동형이다. 예를 들어
$\mathcal{O}_{X'}$로 전사하는 임의의 집합층을 택할 수 있다. 그러면
$$
\NL_{X/S} \cong \NL(\alpha) = 
\left(
\mathcal{J}/\mathcal{J}^2
\longrightarrow
\Omega_{f^{-1}\mathcal{O}_S[\mathcal{E}]/f^{-1}\mathcal{O}_S}
\otimes_{f^{-1}\mathcal{O}_S[\mathcal{E}]} \mathcal{O}_X\right)
$$
이다. 가군, 절 \ref{modules-section-netherlander}, 특히 보조정리
\ref{modules-lemma-NL-up-to-qis}를 보라. 물론 $\alpha'$는 준동형
$f^{-1}\mathcal{O}_S[\mathcal{E}] \to \mathcal{O}_{X'}$
를 정하고, 이는 다시 준동형
$$
\mathcal{J}/\mathcal{J}^2 \longrightarrow \mathcal{G}
$$
를 정하며, 이는 다시 다음 군의 원소를 정한다.\par\noindent
$\Ext^1_{\mathcal{O}_X}(\NL(\alpha), \mathcal{G}) =
\Ext^1_{\mathcal{O}_X}(\NL_{X/S}, \mathcal{G})$
보조정리의 전단사에서 이 원소는 $\mathcal{O}_{X'}$에 대응한다.
\end{remark}

\begin{lemma}
\label{lemma-extensions-of-relative-ringed-spaces-functorial}
$f : (X, \mathcal{O}_X) \to (S, \mathcal{O}_S)$와
$g : (Y, \mathcal{O}_Y) \to (X, \mathcal{O}_X)$를 환 달린 공간의
사상이라 하자. $\mathcal{F}$를 $\mathcal{O}_X$-가군,
$\mathcal{G}$를 $\mathcal{O}_Y$-가군이라 하고
$c : \mathcal{F} \to \mathcal{G}$를 $g$-사상이라 하자. 끝으로 다음을
생각하자.
\begin{enumerate}
\item[(a)] $0 \to \mathcal{F} \to \mathcal{O}_{X'} \to \mathcal{O}_X \to 0$
은 $\xi \in \Ext^1_{\mathcal{O}_X}(\NL_{X/S}, \mathcal{F})$에 대응하는
$f^{-1}\mathcal{O}_S$-대수의 확대이다.
\item[(b)] $0 \to \mathcal{G} \to \mathcal{O}_{Y'} \to \mathcal{O}_Y \to 0$
은 $\zeta \in \Ext^1_{\mathcal{O}_Y}(\NL_{Y/S}, \mathcal{G})$에 대응하는
$g^{-1}f^{-1}\mathcal{O}_S$-대수의 확대이다.
\end{enumerate}
보조정리 \ref{lemma-extensions-of-relative-ringed-spaces}를 보라.
$g$ 및 $c$와 양립하는 $S$-사상 $g' : Y' \to X'$가 존재할
필요충분조건은 $\xi$와 $\zeta$의 상이 다음 군에서 같은 원소인
것이다.
$\Ext^1_{\mathcal{O}_Y}(Lg^*\NL_{X/S}, \mathcal{G})$.
\end{lemma}

\begin{proof}
다음 준동형들이 있으므로 명제는 의미가 있다.
$$
\Ext^1_{\mathcal{O}_X}(\NL_{X/S}, \mathcal{F}) \to
\Ext^1_{\mathcal{O}_Y}(Lg^*\NL_{X/S}, Lg^*\mathcal{F}) \to
\Ext^1_{\mathcal{O}_Y}(Lg^*\NL_{X/S}, \mathcal{G})
$$
첫 준동형들은 사상
$Lg^*\mathcal{F} \to g^*\mathcal{F} \xrightarrow{c} \mathcal{G}$를
사용하고, 다음 준동형은
$$
\Ext^1_{\mathcal{O}_Y}(\NL_{Y/S}, \mathcal{G}) \to
\Ext^1_{\mathcal{O}_Y}(Lg^*\NL_{X/S}, \mathcal{G})
$$
사상 $Lg^*\NL_{X/S} \to \NL_{Y/S}$를 사용한다. 보조정리
\ref{lemma-huge-diagram-ringed-spaces}을 다음 도표에 적용하면
보조정리의 명제를 얻을 수 있다.
$$
\xymatrix{
& 0 \ar[r] &
\mathcal{G} \ar[r] &
\mathcal{O}_{Y'} \ar[r] &
\mathcal{O}_Y \ar[r] & 0 \\
& 0 \ar[r]|\hole & 0 \ar[u] \ar[r] &
\mathcal{O}_S \ar[u] \ar[r]|\hole &
\mathcal{O}_S \ar[u] \ar[r] & 0 \\
0 \ar[r] &
\mathcal{F} \ar[ruu] \ar[r] &
\mathcal{O}_{X'} \ar[r] &
\mathcal{O}_X \ar[ruu] \ar[r] & 0 \\
0 \ar[r] & 0 \ar[ruu]|\hole \ar[u] \ar[r] &
\mathcal{O}_S \ar[ruu]|\hole \ar[u] \ar[r] &
\mathcal{O}_S \ar[ruu]|\hole \ar[u] \ar[r] & 0
}
$$
여기에 보조정리 \ref{lemma-extensions-of-relative-ringed-spaces}와
\ref{lemma-huge-diagram-ringed-spaces}의 증명에 나오는 구성들 사이의
양립성을 함께 사용한다. 그 양립성의 명제와 증명은 생략한다. 직접
논증은 아래의 주를 보라.
\end{proof}

\begin{remark}
\label{remark-extensions-of-relative-ringed-spaces-functorial}
$f : (X, \mathcal{O}_X) \to (S, \mathcal{O}_S)$,
$g : (Y, \mathcal{O}_Y) \to (X, \mathcal{O}_X)$,
$\mathcal{F}$,
$\mathcal{G}$,
$c : \mathcal{F} \to \mathcal{G}$,
$0 \to \mathcal{F} \to \mathcal{O}_{X'} \to \mathcal{O}_X \to 0$,
$\xi \in \Ext^1_{\mathcal{O}_X}(\NL_{X/S}, \mathcal{F})$,
$0 \to \mathcal{G} \to \mathcal{O}_{Y'} \to \mathcal{O}_Y \to 0$,
그리고 $\zeta \in \Ext^1_{\mathcal{O}_Y}(\NL_{Y/S}, \mathcal{G})$를
보조정리 \ref{lemma-extensions-of-relative-ringed-spaces-functorial}에서와
같이 택하자. $c : g^{-1}\mathcal{F} \to \mathcal{G}$에 따른 밀어내기를
사용하여 다음 확대를 구성할 수 있다.
$$
\xymatrix{
0 \ar[r] &
\mathcal{G} \ar[r] &
\mathcal{O}'_1 \ar[r] &
g^{-1}\mathcal{O}_X \ar[r] & 0 \\
0 \ar[r] &
g^{-1}\mathcal{F} \ar[u]^c \ar[r] &
g^{-1}\mathcal{O}_{X'} \ar[u] \ar[r] &
g^{-1}\mathcal{O}_X \ar@{=}[u] \ar[r] & 0
}
$$
$g^\sharp : g^{-1}\mathcal{O}_X \to \mathcal{O}_Y$에 따른 당김을
사용하여 다음 확대를 구성할 수 있다.
$$
\xymatrix{
0 \ar[r] &
\mathcal{G} \ar[r] &
\mathcal{O}_{Y'} \ar[r] &
\mathcal{O}_Y \ar[r] & 0 \\
0 \ar[r] &
\mathcal{G} \ar@{=}[u] \ar[r] &
\mathcal{O}'_2 \ar[u] \ar[r] &
g^{-1}\mathcal{O}_X \ar[u] \ar[r] & 0
}
$$
도표 추적에 따르면, $g$ 및 $c$와 양립하는 $S$-사상 $Y' \to X'$가
존재할 필요충분조건은 $g^{-1}\mathcal{O}_X$의 $\mathcal{G}$에 의한
$g^{-1}f^{-1}\mathcal{O}_S$-대수 확대들로서 $\mathcal{O}'_1$과
$\mathcal{O}'_2$가 동형인 것이다. 보조정리
\ref{lemma-extensions-of-relative-ringed-spaces}에 의해 이 확대들은
다음 등식의 왼쪽 항으로 분류된다.
$$
\Ext^1_{g^{-1}\mathcal{O}_X}(
\NL_{g^{-1}\mathcal{O}_X/g^{-1}f^{-1}\mathcal{O}_S}, \mathcal{G}) =
\Ext^1_{\mathcal{O}_Y}(Lg^*\NL_{X/S}, \mathcal{G})
$$
여기서 등식은 텐서-Hom 수반성과 다음 등식들에서 따른다.
$$
\NL_{g^{-1}\mathcal{O}_X/g^{-1}f^{-1}\mathcal{O}_S} = g^{-1}\NL_{X/S}
\quad\text{그리고}\quad
Lg^*\NL_{X/S} =
g^{-1}\NL_{X/S} \otimes_{g^{-1}\mathcal{O}_X}^\mathbf{L} \mathcal{O}_Y
$$
첫째 등식은 가군, 보조정리 \ref{modules-lemma-pullback-NL}를 보라.
둘째 등식은 유도 당김의 정의에서 따른다. 따라서 보조정리
\ref{lemma-extensions-of-relative-ringed-spaces-functorial}을 보이려면
$\mathcal{O}'_1$이 $\xi$의 상에 대응하고 $\mathcal{O}'_2$가
$\zeta$의 상에 대응함을 보이면 충분하다. $\xi$와
$\mathcal{O}'_1$ 사이의 대응은 주
\ref{remark-extensions-of-relative-ringed-spaces}에서 동치류 $\xi$를
구성한 방법에서 바로 알 수 있다. $\zeta$와 $\mathcal{O}'_2$ 사이의
대응을 위해 먼저 다음 가환 도표를 택한다.
$$
\xymatrix{
\mathcal{E} \ar[d]_{\beta'} \ar[rd]^\beta \\
\mathcal{O}_{Y'} \ar[r] & \mathcal{O}_Y
}
$$
여기서 $g^{-1}f^{-1}\mathcal{O}_S[\mathcal{E}] \to \mathcal{O}_Y$는
핵이 $\mathcal{K}$인 전사 준동형이다. 다음으로 가환 도표
$$
\xymatrix{
\mathcal{E} \ar[d]_{\beta'} &
\mathcal{E}' \ar[l]^\varphi \ar[d]_{\alpha'} \ar[rd]^\alpha \\
\mathcal{O}_{Y'} &
\mathcal{O}'_2 \ar[l] \ar[r] &
g^{-1}\mathcal{O}_X
}
$$
를 택한다. 여기서
$g^{-1}f^{-1}\mathcal{O}_S[\mathcal{E}'] \to g^{-1}\mathcal{O}_X$는
핵이 $\mathcal{J}$인 전사 준동형이다. 예를 들어 집합층으로서
$\mathcal{E}' = \mathcal{E} \amalg \mathcal{O}'_2$로 택하면 된다.
사상 $\varphi$는 복합체의 사상 $\NL(\alpha) \to \NL(\beta)$를
유도하고(표기는 가군, 절 \ref{modules-section-netherlander}에서와
같다), 특히
$\bar\varphi : \mathcal{J}/\mathcal{J}^2 \to \mathcal{K}/\mathcal{K}^2$.
를 유도한다. 이때 $\NL(\alpha) \cong \NL_{Y/S}$이고
$\NL(\beta) \cong \NL_{g^{-1}\mathcal{O}_X/g^{-1}f^{-1}\mathcal{O}_S}$
이다. 복합체의 사상 $\NL(\alpha) \to \NL(\beta)$는 보조정리
\ref{lemma-extensions-of-relative-ringed-spaces-functorial}의 명제에서
사용한 사상 $Lg^*\NL_{X/S} \to \NL_{Y/S}$를 나타낸다(그 증명의 첫
부분을 보라). 이제 $\zeta$는 $\beta'$가 유도하는 준동형
$\mathcal{K}/\mathcal{K}^2 \to \mathcal{G}$의 동치류에 대응한다.
주 \ref{remark-extensions-of-relative-ringed-spaces}를 보라. 마찬가지로
확대 $\mathcal{O}'_2$는 $\alpha'$가 유도하는 준동형
$\mathcal{J}/\mathcal{J}^2 \to \mathcal{G}$에 대응한다. 위의 가환
도표에 의해 이 준동형은 $\beta'$가 유도하는 준동형
$\mathcal{K}/\mathcal{K}^2 \to \mathcal{G}$와 다음 준동형의 합성이다.
$\bar\varphi : \mathcal{J}/\mathcal{J}^2 \to \mathcal{K}/\mathcal{K}^2$.
이것이 필요한 양립성을 증명한다.
\end{remark}

\begin{lemma}
\label{lemma-parametrize-solutions-ringed-spaces}
$t : (S, \mathcal{O}_S) \to (S', \mathcal{O}_{S'})$,
$\mathcal{J} = \Ker(t^\sharp)$,
$f : (X, \mathcal{O}_X) \to (S, \mathcal{O}_S)$, $\mathcal{G}$,
그리고 $c : \mathcal{J} \to \mathcal{G}$를
(\ref{equation-to-solve-ringed-spaces})에서와 같이 택하자. 보조정리
\ref{lemma-extensions-of-relative-ringed-spaces}에 의해
$\mathcal{O}_S$의 $\mathcal{J}$에 의한 확대 $\mathcal{O}_{S'}$에
대응하는 원소를
$\xi \in \Ext^1_{\mathcal{O}_S}(\NL_{S/S'}, \mathcal{J})$라 쓰자.
해들의 동형류 집합은 다음 준동형에서 $\xi$의 상 위의 올과
표준적으로 전단사이다.
$$
\Ext^1_{\mathcal{O}_X}(\NL_{X/S'}, \mathcal{G}) \to
\Ext^1_{\mathcal{O}_X}(Lf^*\NL_{S/S'}, \mathcal{G})
$$
\end{lemma}

\begin{proof}
보조정리 \ref{lemma-extensions-of-relative-ringed-spaces}를
$X \to S'$와 $\mathcal{O}_X$-가군 $\mathcal{G}$에 적용하면, 다음
군의 원소 $\zeta$들이
$\Ext^1_{\mathcal{O}_X}(\NL_{X/S'}, \mathcal{G})$
다음 $f^{-1}\mathcal{O}_{S'}$-대수의 확대들을 매개화함을 알 수 있다.
$0 \to \mathcal{G} \to \mathcal{O}_{X'} \to \mathcal{O}_X \to 0$
보조정리
\ref{lemma-extensions-of-relative-ringed-spaces-functorial}을
$X \to S \to S'$와 $c : \mathcal{J} \to \mathcal{G}$에 적용하면,
$c$ 및 $f : X \to S$와 양립하는 $S'$-사상 $X' \to S'$가 존재할
필요충분조건은 $\zeta$가 $\xi$로 가는 것임을 알 수 있다. 이는
정확히 $\mathcal{O}_{X'}$이
(\ref{equation-to-solve-ringed-spaces})의 해라는 뜻이다.
\end{proof}

\begin{remark}
\label{remark-parametrize-solutions-ringed-spaces}
보조정리 \ref{lemma-parametrize-solutions-ringed-spaces}의 상황에서는
다음 복합체의 사상들이 있다.
$$
Lf^*\NL_{S'/S} \to \NL_{X/S'} \to \NL_{X/S}
$$
이 사상들은 완전 삼각형을 이루는 것에 가깝다. 가군, 보조정리
\ref{modules-lemma-exact-sequence-NL-ringed-topoi}를 보라. 만일 이것이
완전 삼각형이라면, $\Ext^2_{\mathcal{O}_X}(\NL_{X/S}, \mathcal{G})$에서
$\xi$의 상이 (\ref{equation-to-solve-ringed-spaces})의 해의 존재에
대한 장애 원소라고 결론 내릴 수 있을 것이다.
\end{remark}













\section{스킴의 변형}
\label{section-deformations-schemes}

\noindent
이 절에서는 절 \ref{section-deformations-ringed-spaces}의 결과들이
스킴의 변형에 대하여 무엇을 뜻하는지 명시한다.

\begin{lemma}
\label{lemma-deform}
$S \subset S'$를 스킴의 1차 두꺼워짐이라 하고 $f : X \to S$를
스킴의 평탄 사상이라 하자. 스킴의 평탄 사상 $f' : X' \to S'$와
$S$ 위의 동형사상 $a : X \to X' \times_{S'} S$가 존재하면 다음이
성립한다.
\begin{enumerate}
\item 쌍 $(f' : X' \to S', a)$들의 동형류 집합은
$\Ext^1_{\mathcal{O}_X}(\NL_{X/S}, f^*\mathcal{C}_{S/S'})$ 아래의
주동차 공간이다.
\item $X' \times_{S'} S$ 위에서 항등사상으로 환원되는 $S'$ 위의
자기동형사상 $\varphi : X' \to X'$들의 집합은
$\Ext^0_{\mathcal{O}_X}(\NL_{X/S}, f^*\mathcal{C}_{S/S'})$이다.
\end{enumerate}
\end{lemma}

\begin{proof}
먼저 사상의 심화, 절 \ref{more-morphisms-section-thickenings}에서 정의한
스킴의 두꺼워짐은 절 \ref{section-thickenings-spaces}의 의미에서
두꺼워짐인 스킴의 사상과 같은 것임을 관찰하자. $X$를 $X'$의
닫힌 부분스킴으로 볼 수 있으며, 이때
$(f, f') : (X \subset X') \to (S \subset S')$는 1차 두꺼워짐 사이의
사상이다. 그러면 사상의 심화, 보조정리
\ref{more-morphisms-lemma-deform}(또는 더 일반적인 보조정리
\ref{lemma-deform-module})에 의해 $X'$ 안의 $X$의 아이디얼층은
$f^*\mathcal{C}_{S/S'}$와 같다. 따라서 다음 가환 도표가 있다.
$$
\xymatrix{
0 \ar[r] & f^*\mathcal{C}_{S/S'} \ar[r] &
\mathcal{O}_{X'} \ar[r] &
\mathcal{O}_X \ar[r] & 0 \\
0 \ar[r] & \mathcal{C}_{S/S'} \ar[u] \ar[r] &
\mathcal{O}_{S'} \ar[u] \ar[r] &
\mathcal{O}_S \ar[u] \ar[r] & 0
}
$$
여기서 세로 화살표들은 $f$-사상이다.
(\ref{equation-to-solve-ringed-spaces})와 비교하라. 따라서 (1)은
보조정리 \ref{lemma-choices-ringed-spaces}에서, (2)는 보조정리
\ref{lemma-huge-diagram-ringed-spaces}의 (2)에서 따른다. 사상의 심화,
절 \ref{more-morphisms-section-netherlander}에서 스킴의 사상에 대해
정의한 $\NL_{X/S}$는 절 \ref{section-deformations-ringed-spaces}에서
사용한 $\NL_{X/S}$와 일치함에 유의하라.
\end{proof}










\section{환 달린 토포스의 두꺼워짐}
\label{section-thickenings-ringed-topoi}

\noindent
이 절은 환 달린 토포스에 대한 절 \ref{section-thickenings-spaces}의
유사형이다. 이어지는 몇 절에서는 다음 개념들을 사용한다.
\begin{enumerate}
\item 환 달린 토포스 $(\Sh(\mathcal{D}), \mathcal{O}')$ 위의 아이디얼층
$\mathcal{I} \subset \mathcal{O}'$에서 $\mathcal{I}$의 모든 국소절단이
국소적으로 멱영이면, 이 아이디얼층을 {\it 국소 멱영}이라고 한다.
\item 환 달린 토포스의 {\it 두꺼워짐}은 환 달린 토포스의 사상
$i : (\Sh(\mathcal{C}), \mathcal{O}) \to (\Sh(\mathcal{D}), \mathcal{O}')$
으로서 다음을 만족하는 것이다.
\begin{enumerate}
\item $i_*$는 동치 $\Sh(\mathcal{C}) \to \Sh(\mathcal{D})$이다.
\item 사상 $i^\sharp : \mathcal{O}' \to i_*\mathcal{O}$은 전사이다.
\item $i^\sharp$의 핵은 국소 멱영 아이디얼층이다.
\end{enumerate}
\item 환 달린 토포스의 {\it 1차 두꺼워짐}은 환 달린 토포스의 두꺼워짐
$i : (\Sh(\mathcal{C}), \mathcal{O}) \to (\Sh(\mathcal{D}), \mathcal{O}')$
으로서 $\Ker(i^\sharp)$의 제곱이 영인 것이다.
\item {\it 환 달린 토포스의 두꺼워짐 사이의 사상},
{\it 바탕 환 달린 토포스 위의 환 달린 토포스의 두꺼워짐 사이의 사상}
등의 정의는 명백하다.
\end{enumerate}
$i : (\Sh(\mathcal{C}), \mathcal{O}) \to (\Sh(\mathcal{D}), \mathcal{O}')$
가 환 달린 토포스의 두꺼워짐이면 바탕 토포스들을 동일시하고
$\mathcal{O}$, $\mathcal{O}'$, 그리고
$\mathcal{I} = \Ker(i^\sharp)$를 $\mathcal{C}$ 위의 층으로 본다.
다음 $\mathcal{O}'$-가군의 짧은 완전열을 얻는다.
$$
0 \to \mathcal{I} \to \mathcal{O}' \to \mathcal{O} \to 0
$$
사이트 위의 가군, 보조정리
\ref{sites-modules-lemma-i-star-equivalence}에 의해
$\mathcal{O}$-가군의 범주는 $\mathcal{I}$에 의해 소멸되는
$\mathcal{O}'$-가군의 범주와 동치이다. 특히 $i$가 1차 두꺼워짐이면
$\mathcal{I}$은 $\mathcal{O}$-가군이다.

\begin{situation}
\label{situation-morphism-thickenings-ringed-topoi}
환 달린 토포스의 두꺼워짐 사이의 사상 $(f, f')$는 다음 가환 도표로
주어진다.
\begin{equation}
\label{equation-morphism-thickenings-ringed-topoi}
\vcenter{
\xymatrix{
(\Sh(\mathcal{C}), \mathcal{O}) \ar[r]_i \ar[d]_f &
(\Sh(\mathcal{D}), \mathcal{O}') \ar[d]^{f'} \\
(\Sh(\mathcal{B}), \mathcal{O}_\mathcal{B}) \ar[r]^t &
(\Sh(\mathcal{B}'), \mathcal{O}_{\mathcal{B}'})
}
}
\end{equation}
가로 화살표들은 두꺼워짐이다. 이 상황에서
$\mathcal{I} = \Ker(i^\sharp) \subset \mathcal{O}'$ 및
$\mathcal{J} = \Ker(t^\sharp) \subset \mathcal{O}_{\mathcal{B}'}$라
놓는다. 바탕 토포스 위에서 $f = f'$이므로 당김 함자 $f^{-1}$과
$(f')^{-1}$을 동일시한다. 준동형
$(f')^\sharp : f^{-1}\mathcal{O}_{\mathcal{B}'} \to \mathcal{O}'$
는 특히 준동형 $f^{-1}\mathcal{J} \to \mathcal{I}$을 유도하고, 따라서
다음 $\mathcal{O}'$-가군의 준동형을 유도함에 유의하자.
$$
(f')^*\mathcal{J} \longrightarrow \mathcal{I}
$$
$i$와 $t$가 1차 두꺼워짐이면 $(f')^*\mathcal{J} = f^*\mathcal{J}$이고,
위 준동형은 $f^*\mathcal{J} \to \mathcal{I}$이 된다.
\end{situation}

\begin{definition}
\label{definition-strict-morphism-thickenings-ringed-topoi}
상황 \ref{situation-morphism-thickenings-ringed-topoi}에서 준동형
$(f')^*\mathcal{J} \longrightarrow \mathcal{I}$이 전사이면
$(f, f')$를 {\it 두꺼워짐의 엄격 사상}이라고 한다.
\end{definition}








\section{환 달린 토포스의 1차 두꺼워짐 위의 가군}
\label{section-modules-thickenings-ringed-topoi}

\noindent
이 절에서는 가군의 변형 이론을 위한 몇 가지 준비 사항을 다룬다.
$i : (\Sh(\mathcal{C}, \mathcal{O}) \to (\Sh(\mathcal{D}), \mathcal{O}')$
를 환 달린 토포스의 1차 두꺼워짐이라 하자. 절
\ref{section-thickenings-ringed-topoi}에서 도입한 표기를 자유롭게
사용하며, 특히 바탕 토포스들을 동일시한다. 이 절에서는 다음
짧은 완전열을 생각한다.
\begin{equation}
\label{equation-extension-ringed-topoi}
0 \to \mathcal{K} \to \mathcal{F}' \to \mathcal{F} \to 0
\end{equation}
이는 $\mathcal{O}'$-가군의 열이며, $\mathcal{F}$와 $\mathcal{K}$는
$\mathcal{O}$-가군이고 $\mathcal{F}'$는 $\mathcal{O}'$-가군이다.
이 상황에서는 표준 $\mathcal{O}$-가군 준동형
$$
c_{\mathcal{F}'} :
\mathcal{I} \otimes_\mathcal{O} \mathcal{F}
\longrightarrow
\mathcal{K}
$$
가 있다. 여기서 $\mathcal{I} = \Ker(i^\sharp)$이다. 구체적으로
$\mathcal{I}$의 국소절단 $f$와 $\mathcal{F}$의 국소절단 $s$가
주어졌을 때, $s$를 올리는 $\mathcal{F}'$의 국소절단을 $s'$라 하고
$c_{\mathcal{F}'}(f \otimes s) = fs'$로 놓는다.

\begin{lemma}
\label{lemma-inf-map-ringed-topoi}
$i : (\Sh(\mathcal{C}), \mathcal{O}) \to (\Sh(\mathcal{D}), \mathcal{O}')$
를 환 달린 토포스의 1차 두꺼워짐이라 하자. 다음 확대들과
$$
0 \to \mathcal{K} \to \mathcal{F}' \to \mathcal{F} \to 0
\quad\text{그리고}\quad
0 \to \mathcal{L} \to \mathcal{G}' \to \mathcal{G} \to 0
$$
(\ref{equation-extension-ringed-topoi})에서와 같은 준동형
$\varphi : \mathcal{F} \to \mathcal{G}$ 및
$\psi : \mathcal{K} \to \mathcal{L}$이 주어졌다고 가정하자.
\begin{enumerate}
\item $\varphi$ 및 $\psi$와 양립하는 $\mathcal{O}'$-가군 준동형
$\varphi' : \mathcal{F}' \to \mathcal{G}'$가 존재하면 다음 도표는
가환한다.
$$
\xymatrix{
\mathcal{I} \otimes_\mathcal{O} \mathcal{F}
\ar[r]_-{c_{\mathcal{F}'}} \ar[d]_{1 \otimes \varphi} &
\mathcal{K} \ar[d]^\psi \\
\mathcal{I} \otimes_\mathcal{O} \mathcal{G}
\ar[r]^-{c_{\mathcal{G}'}} &
\mathcal{L}
}
$$
\item $\varphi$ 및 $\psi$와 양립하는 $\mathcal{O}'$-가군 준동형
$\varphi' : \mathcal{F}' \to \mathcal{G}'$들의 집합은 공집합이
아니면 다음 군 아래의 주동차 공간이다.
$\Hom_\mathcal{O}(\mathcal{F}, \mathcal{L})$.
\end{enumerate}
\end{lemma}

\begin{proof}
(1)은 준동형들의 설명에서 바로 따른다. (2)를 위해 $\varphi'$와
$\varphi''$가 $\varphi$ 및 $\psi$와 양립하는 두 준동형
$\mathcal{F}' \to \mathcal{G}'$이라 하자. 그러면
$\varphi' - \varphi''$는 다음과 같이 분해된다.
$$
\mathcal{F}' \to \mathcal{F} \to \mathcal{L} \to \mathcal{G}'
$$
사이트 위의 가군, 보조정리
\ref{sites-modules-lemma-i-star-equivalence}에 의해 가운데 준동형은
$\Hom_\mathcal{O}(\mathcal{F}, \mathcal{L})$의 유일한 원소에서 온다.
역으로 이 군의 원소 $\alpha$가 주어지면 위의 합성에서 가운데
준동형을 $\alpha$로 택하여 $\varphi'$에 더할 수 있다. 세부 사항은
생략한다.
\end{proof}

\begin{lemma}
\label{lemma-inf-obs-map-ringed-topoi}
$i : (\Sh(\mathcal{C}), \mathcal{O}) \to (\Sh(\mathcal{D}), \mathcal{O}')$
를 환 달린 토포스의 1차 두꺼워짐이라 하자. 다음 확대들과
$$
0 \to \mathcal{K} \to \mathcal{F}' \to \mathcal{F} \to 0
\quad\text{그리고}\quad
0 \to \mathcal{L} \to \mathcal{G}' \to \mathcal{G} \to 0
$$
(\ref{equation-extension-ringed-topoi})에서와 같은 준동형
$\varphi : \mathcal{F} \to \mathcal{G}$ 및
$\psi : \mathcal{K} \to \mathcal{L}$이 주어졌다고 하자. 다음 도표가
가환한다고 가정한다.
$$
\xymatrix{
\mathcal{I} \otimes_\mathcal{O} \mathcal{F}
\ar[r]_-{c_{\mathcal{F}'}} \ar[d]_{1 \otimes \varphi} &
\mathcal{K} \ar[d]^\psi \\
\mathcal{I} \otimes_\mathcal{O} \mathcal{G}
\ar[r]^-{c_{\mathcal{G}'}} &
\mathcal{L}
}
$$
그러면 다음 원소가 존재한다.
$$
o(\varphi, \psi) \in
\Ext^1_\mathcal{O}(\mathcal{F}, \mathcal{L})
$$
이 원소가 영이라는 것은 $\varphi$와 $\psi$에 양립하는 준동형
$\varphi' : \mathcal{F}' \to \mathcal{G}'$가 존재하기 위한
필요충분조건이다.
\end{lemma}

\begin{proof}
다음 확대를 명시적으로 구성할 수 있다.
$$
0 \to \mathcal{L} \to \mathcal{H} \to \mathcal{F} \to 0
$$
$\mathcal{H}$는 다음 복합체의 가운데 코호몰로지로 잡는다.
$$
\mathcal{K}
\xrightarrow{1, - \psi}
\mathcal{F}' \oplus \mathcal{G}' \xrightarrow{\varphi, 1}
\mathcal{G}
$$
표기는 자명하다. 보조정리의 도표가 가환한다는 가정을 사용하여 국소
절단으로 계산하면 $\mathcal{H}$가 $\mathcal{I}$에 의해 소멸됨을 알 수
있다. 따라서 $\mathcal{H}$는 다음 군의 동치류를 정한다.
$$
\Ext^1_\mathcal{O}(\mathcal{F}, \mathcal{L})
\subset
\Ext^1_{\mathcal{O}'}(\mathcal{F}, \mathcal{L})
$$
끝으로 $\mathcal{H}$의 동치류는 $\psi$에 따른 확대 $\mathcal{F}'$의
밀어내기와 $\varphi$에 따른 확대 $\mathcal{G}'$의 당김의 차이다.
계산은 생략한다. 따라서 $\mathcal{H}$의 동치류가 영인 것과 다음과 같은
가환 도표가 존재하는 것은 동치이다.
$$
\xymatrix{
0 \ar[r] &
\mathcal{K} \ar[r] \ar[d]_{\psi} &
\mathcal{F}' \ar[r] \ar[d]_{\varphi'} &
\mathcal{F} \ar[r] \ar[d]_\varphi & 0\\
0 \ar[r] &
\mathcal{L} \ar[r] &
\mathcal{G}' \ar[r] &
\mathcal{G} \ar[r] & 0
}
$$
이는 원하는 결론이다.
\end{proof}

\begin{lemma}
\label{lemma-inf-ext-ringed-topoi}
$i : (\Sh(\mathcal{C}), \mathcal{O}) \to (\Sh(\mathcal{D}), \mathcal{O}')$
를 환 달린 토포스의 1차 두꺼워짐이라 하자. $\mathcal{O}$-가군
$\mathcal{F}$, $\mathcal{K}$와 $\mathcal{O}$-선형 준동형
$c : \mathcal{I} \otimes_\mathcal{O} \mathcal{F} \to \mathcal{K}$.
이 주어졌다고 하자. $c_{\mathcal{F}'} = c$인 완전열
(\ref{equation-extension-ringed-topoi})이 하나라도 존재하면, 그러한
확대들의 동형류 집합은
$\Ext^1_\mathcal{O}(\mathcal{F}, \mathcal{K})$ 아래의 주동차 공간이다.
\end{lemma}

\begin{proof}
다음 확대들이 주어졌다고 하자.
$$
0 \to \mathcal{K} \to \mathcal{F}'_1 \to \mathcal{F} \to 0
\quad\text{그리고}\quad
0 \to \mathcal{K} \to \mathcal{F}'_2 \to \mathcal{F} \to 0
$$
$c_{\mathcal{F}'_1} = c_{\mathcal{F}'_2} = c$라 하자. 확대군에서 두
확대의 차를 취하면(호몰로지 대수학, 절
\ref{homology-section-extensions} 참조) 다음 확대를 얻는다.
$$
0 \to \mathcal{K} \to \mathcal{E} \to \mathcal{F} \to 0
$$
여기서 $\mathcal{E}$는 $\mathcal{I}$에 의해 소멸된다. 국소 계산은
생략한다. 따라서 이 완전열은 $\mathcal{O}$-가군의 확대이다. 사이트
위의 가군, 보조정리 \ref{sites-modules-lemma-i-star-equivalence}를 보라.
역으로 이러한 확대 $\mathcal{E}$가 주어지면
$c_{\mathcal{F}'}$를 바꾸지 않고 $\mathcal{E}$를
$\mathcal{O}'$-확대 $\mathcal{F}'$에 더할 수 있다. 몇 가지 세부 사항은
생략한다.
\end{proof}

\begin{lemma}
\label{lemma-inf-obs-ext-ringed-topoi}
$i : (\Sh(\mathcal{C}), \mathcal{O}) \to (\Sh(\mathcal{D}), \mathcal{O}')$
를 환 달린 토포스의 1차 두꺼워짐이라 하자. $\mathcal{O}$-가군
$\mathcal{F}$, $\mathcal{K}$와 $\mathcal{O}$-선형 준동형
$c : \mathcal{I} \otimes_\mathcal{O} \mathcal{F} \to \mathcal{K}$.
이 주어졌다고 하자. 그러면 다음 원소가 존재한다.
$$
o(\mathcal{F}, \mathcal{K}, c) \in
\Ext^2_\mathcal{O}(\mathcal{F}, \mathcal{K})
$$
이 원소가 영이라는 것은 $c_{\mathcal{F}'} = c$인 완전열
(\ref{equation-extension-ringed-topoi})이 존재하기 위한 필요충분조건이다.
\end{lemma}

\begin{proof}
먼저 $\mathcal{K}$가 단사 $\mathcal{O}$-가군이면
$c_{\mathcal{F}'} = c$인 완전열
(\ref{equation-extension-ringed-topoi})이 실제로 존재함을 보이자. 평탄
$\mathcal{O}'$-가군 $\mathcal{H}'$과 전사 준동형
$\mathcal{H}' \to \mathcal{F}$
을 택한다. 사이트 위의 가군, 보조정리
\ref{sites-modules-lemma-module-quotient-flat}를 보라. 그 핵을
$\mathcal{J} \subset \mathcal{H}'$라 하자. $\mathcal{H}'$은 평탄하므로
$$
\mathcal{I} \otimes_{\mathcal{O}'} \mathcal{H}' =
\mathcal{I}\mathcal{H}'
\subset \mathcal{J} \subset \mathcal{H}'
$$
이다. 다음 준동형은 $\mathcal{I}\mathcal{J}$를 영으로 보낸다.
$$
\mathcal{I}\mathcal{H}' =
\mathcal{I} \otimes_{\mathcal{O}'} \mathcal{H}'
\longrightarrow
\mathcal{I} \otimes_{\mathcal{O}'} \mathcal{F} =
\mathcal{I} \otimes_\mathcal{O} \mathcal{F}
$$
실제로 $f$가 $\mathcal{I}$의 국소절단이고 $s$가 $\mathcal{H}$의
국소절단이면 $fs$는 $f \otimes \overline{s}$로 간다. 여기서
$\overline{s}$는 $\mathcal{F}$에서 $s$의 상이다. 따라서 다음과 같은
$\mathcal{O}$-가군의 도표를 얻는다.
$$
\xymatrix{
\mathcal{I}\mathcal{H}'/\mathcal{I}\mathcal{J}
\ar@{^{(}->}[r] \ar[d] &
\mathcal{J}/\mathcal{I}\mathcal{J} \ar@{..>}[d]_\gamma \\
\mathcal{I} \otimes_\mathcal{O} \mathcal{F} \ar[r]^-c &
\mathcal{K}
}
$$
$\mathcal{K}$가 단사 $\mathcal{O}$-가군이면 점선 화살표를 얻는다.
$\gamma$와 $\mathcal{J} \to \mathcal{J}/\mathcal{I}\mathcal{J}$의 합성을
$\gamma' : \mathcal{J} \to \mathcal{K}$라 쓰자. 국소 계산에 따르면
다음 밀어내기는
$$
\xymatrix{
0 \ar[r] &
\mathcal{J} \ar[r] \ar[d]_{\gamma'} &
\mathcal{H}' \ar[r] \ar[d] &
\mathcal{F} \ar[r] \ar@{=}[d] &
0 \\
0 \ar[r] &
\mathcal{K} \ar[r] &
\mathcal{F}' \ar[r] &
\mathcal{F} \ar[r] &
0
}
$$
보조정리에서 제기한 문제의 해이다.

\medskip\noindent
이제 일반적인 경우를 다루자. $\mathcal{K}'$가 단사
$\mathcal{O}$-가군이 되도록 매장 $\mathcal{K} \subset \mathcal{K}'$을
택한다. 그 몫을 $\mathcal{Q}$라 하면 완전열
$$
0 \to \mathcal{K} \to \mathcal{K}' \to \mathcal{Q} \to 0
$$
을 얻는다. 합성 준동형을
$c' : \mathcal{I} \otimes_\mathcal{O} \mathcal{F} \to \mathcal{K}'$라
쓰자. 앞 문단에 의해 다음과 같은 완전열이 존재한다.
$$
0 \to \mathcal{K}' \to \mathcal{E}' \to \mathcal{F} \to 0
$$
이는 (\ref{equation-extension-ringed-topoi})에서와 같고
$c_{\mathcal{E}'} = c'$이다. $c'$와 $\mathcal{K}' \to \mathcal{Q}$의
합성은 영이므로, $\mathcal{K}' \to \mathcal{Q}$에 따른
$\mathcal{E}'$의 밀어내기는 다음 확대이다.
$$
0 \to \mathcal{Q} \to \mathcal{D}' \to \mathcal{F} \to 0
$$
이는 (\ref{equation-extension-ringed-topoi})에서와 같고
$c_{\mathcal{D}'} = 0$이다. 이는 정확히 $\mathcal{D}'$가
$\mathcal{I}$에 의해 소멸됨을 뜻한다. 다시 말해 $\mathcal{D}'$는
$\mathcal{O}$-가군의 확대이며 다음 원소를 정한다.
$$
o(\mathcal{F}, \mathcal{K}, c) \in
\Ext^1_\mathcal{O}(\mathcal{F}, \mathcal{Q}) =
\Ext^2_\mathcal{O}(\mathcal{F}, \mathcal{K})
$$
등식은 위 완전열에 딸린 긴 완전 코호몰로지 열과 단사 가군
$\mathcal{K}'$로 가는 고차 Ext 군들의 소멸에서 따른다.
$o(\mathcal{F}, \mathcal{K}, c) = 0$이면 분할
$s : \mathcal{F} \to \mathcal{D}'$를 택하고
$$
\mathcal{F}' = \Ker(\mathcal{E}' \to \mathcal{D}'/s(\mathcal{F}))
$$
로 놓을 수 있다. 그러면 완전한 두 행을 가진 다음 도표를 얻으며,
$$
\xymatrix{
0 \ar[r] &
\mathcal{K} \ar[r] \ar[d] &
\mathcal{F}' \ar[r] \ar[d] &
\mathcal{F} \ar[r] \ar@{=}[d] &
0 \\
0 \ar[r] &
\mathcal{K}' \ar[r] &
\mathcal{E}' \ar[r] &
\mathcal{F} \ar[r] & 0
}
$$
이는 $c_{\mathcal{F}'} = c$임을 보인다. 역으로 $\mathcal{F}'$가
존재하면 보조정리 \ref{lemma-inf-ext-ringed-topoi}와 단사 가군
$\mathcal{K}'$로 가는 고차 Ext 군들의 소멸에 의해,
$\mathcal{K} \to \mathcal{K}'$에 따른 $\mathcal{F}'$의 밀어내기는
$\mathcal{E}'$와 동형이다. 이로부터 위와 같은 도표를 얻고,
$\mathcal{D}'$가 확대로서 분할됨을 알 수 있다. 즉 동치류
$o(\mathcal{F}, \mathcal{K}, c)$는 영이다.
\end{proof}

\begin{remark}
\label{remark-trivial-thickening-ringed-topoi}
$(\Sh(\mathcal{C}), \mathcal{O})$를 환 달린 토포스라 하자. 1차 두꺼워짐
$i : (\Sh(\mathcal{C}), \mathcal{O}) \to
(\Sh(\mathcal{D}), \mathcal{O}')$에 대하여 그 왼쪽 역인 환 달린 토포스의
사상
$\pi : (\Sh(\mathcal{D}), \mathcal{O}') \to (\Sh(\mathcal{C}), \mathcal{O})$
가 존재하면 $i$를 {\it 자명하다}고 한다. 그러한 사상 $\pi$의 선택을
이 1차 두꺼워짐의 {\it 자명화}라고 한다. $\pi$가 주어지면 다음 분할을
얻는다.
\begin{equation}
\label{equation-splitting-ringed-topoi}
\mathcal{O}' = \mathcal{O} \oplus \mathcal{I}
\end{equation}
$\pi^\sharp$로 전사 준동형 $\mathcal{O}' \to \mathcal{O}$를 분할하여
얻은 $\mathcal{C}$ 위의 대수층의 분할이다. 역으로 이러한 분할은 사상
$\pi$를 정한다. $(\Sh(\mathcal{C}), \mathcal{O})$의 자명화된 1차
두꺼워짐들의 범주는 $\mathcal{O}$-가군의 범주와 동치이다.
\end{remark}

\begin{remark}
\label{remark-trivial-extension-ringed-topoi}
$i : (\Sh(\mathcal{C}), \mathcal{O}) \to (\Sh(\mathcal{D}), \mathcal{O}')$
를 환 달린 토포스의 자명한 1차 두꺼워짐이라 하고
$\pi : (\Sh(\mathcal{D}), \mathcal{O}') \to
(\Sh(\mathcal{C}), \mathcal{O})$를 그 자명화라 하자. 두
$\mathcal{O}$-가군과 준동형
$c : \mathcal{I} \otimes_\mathcal{O} \mathcal{F} \to \mathcal{K}$
으로 이루어진 임의의 삼중항 $(\mathcal{F}, \mathcal{K}, c)$에 대하여
다음과 같이 놓을 수 있다.
$$
\mathcal{F}'_{c, triv} = \mathcal{F} \oplus \mathcal{K}
$$
$\pi$에 딸린 분할 (\ref{equation-splitting-ringed-topoi})과 준동형 $c$를 사용하여
$\mathcal{O}'$-가군 구조를 정의하면 확대
(\ref{equation-extension-ringed-topoi})을 얻는다.
$\mathcal{F}'_{c, triv}$를 $c$와 자명화 $\pi$에 대응하는,
$\mathcal{F}$를 $\mathcal{K}$로 확대한 {\it 자명한 확대}라고 부른다.
(\ref{equation-extension-ringed-topoi})과 같은 임의의 확대
$\mathcal{F}'$에 대하여 $\pi^\sharp : \mathcal{O} \to \mathcal{O}'$를
사용하면 $\mathcal{F}'$를 $\mathcal{O}$-가군의 확대로 볼 수 있고,
따라서 $\Ext^1_\mathcal{O}(\mathcal{F}, \mathcal{K})$의 동치류
$\xi_{\mathcal{F}'}$를 얻는다. 보조정리
\ref{lemma-inf-ext-ringed-topoi}에 의해
$\mathcal{F}' \mapsto \xi_{\mathcal{F}'}$
는 전단사 대응
$$
\left\{
\begin{matrix}
\text{확대의 동형류}\\
\mathcal{F}'\text{는 (\ref{equation-extension-ringed-topoi})과 같고 }
c = c_{\mathcal{F}'}
\end{matrix}
\right\}
\longrightarrow
\Ext^1_\mathcal{O}(\mathcal{F}, \mathcal{K})
$$
를 유도한다. 또한 자명한 확대 $\mathcal{F}'_{c, triv}$는 영 동치류로 간다.
\end{remark}

\begin{remark}
\label{remark-extension-functorial-ringed-topoi}
$(\Sh(\mathcal{C}), \mathcal{O})$를 환 달린 토포스라 하자.
$(\Sh(\mathcal{C}), \mathcal{O}) \to (\Sh(\mathcal{D}_i), \mathcal{O}'_i)$,
$i = 1, 2$를 아이디얼층이 $\mathcal{I}_i$인 1차 두꺼워짐이라 하자.
$h : (\Sh(\mathcal{D}_1), \mathcal{O}'_1) \to
(\Sh(\mathcal{D}_2), \mathcal{O}'_2)$
를 $(\Sh(\mathcal{C}), \mathcal{O})$의 1차 두꺼워짐 사이의 사상이라
하자. 도표는
$$
\xymatrix{
& (\Sh(\mathcal{C}), \mathcal{O}) \ar[ld] \ar[rd] & \\
(\Sh(\mathcal{D}_1), \mathcal{O}'_1) \ar[rr]^h & & 
(\Sh(\mathcal{D}_2), \mathcal{O}'_2)
}
$$
와 같다. 특히 $h^\sharp : \mathcal{O}'_2 \to \mathcal{O}'_1$은
$\mathcal{O}$-가군 준동형 $\mathcal{I}_2 \to \mathcal{I}_1$을
유도한다. $\mathcal{F}$를 $\mathcal{O}$-가군이라 하자. 각
$i = 1, 2$에 대하여 $(\mathcal{K}_i, c_i)$는 $\mathcal{O}$-가군
$\mathcal{K}_i$와 준동형
$c_i : \mathcal{I}_i \otimes_\mathcal{O} \mathcal{F} \to
\mathcal{K}_i$로 이루어졌다고 하자. 또한 다음 도표를 가환시키는
$\mathcal{O}$-가군 준동형 $\mathcal{K}_2 \to \mathcal{K}_1$이
주어졌다고 하자.
$$
\xymatrix{
\mathcal{I}_2 \otimes_\mathcal{O} \mathcal{F}
\ar[r]_-{c_2} \ar[d] &
\mathcal{K}_2 \ar[d] \\
\mathcal{I}_1 \otimes_\mathcal{O} \mathcal{F}
\ar[r]^-{c_1} &
\mathcal{K}_1
}
$$
그러면 다음과 같은 표준적인 함자성이 있다.
$$
\left\{
\begin{matrix}
\mathcal{F}'_2\text{는 (\ref{equation-extension-ringed-topoi})과 같고 }\\
c_2 = c_{\mathcal{F}'_2}\text{이며 }\mathcal{K} = \mathcal{K}_2
\end{matrix}
\right\}
\longrightarrow
\left\{
\begin{matrix}
\mathcal{F}'_1\text{는 (\ref{equation-extension-ringed-topoi})과 같고 }\\
c_1 = c_{\mathcal{F}'_1}\text{이며 }\mathcal{K} = \mathcal{K}_1
\end{matrix}
\right\}
$$
구체적으로 모든 층 $\mathcal{O}$, $\mathcal{O}'_i$,
$\mathcal{F}$, $\mathcal{K}_i$ 등을 $\mathcal{C}$ 위의 층으로 본다.
$\mathcal{F}'_2$가 주어지면 $\mathcal{F}'_1$을 다음 확대 도표에
들어가는 밀어내기로 정의한다.
$$
\xymatrix{
0 \ar[r] &
\mathcal{K}_2 \ar[r] \ar[d] &
\mathcal{F}'_2 \ar[r] \ar[d] &
\mathcal{F} \ar@{=}[d] \ar[r] & 0 \\
0 \ar[r] &
\mathcal{K}_1 \ar[r] &
\mathcal{F}'_1 \ar[r] &
\mathcal{F} \ar[r] & 0
}
$$
이 밀어내기 위의 $\mathcal{O}'_1$-가군 구조의 구성은 생략한다.
이 구성에는 $c_1$과 $c_2$가 들어가는 도표의 가환성이 사용된다.
\end{remark}

\begin{remark}
\label{remark-trivial-extension-functorial-ringed-topoi}
$(\Sh(\mathcal{C}), \mathcal{O})$,
$(\Sh(\mathcal{C}), \mathcal{O}) \to (\Sh(\mathcal{D}_i), \mathcal{O}'_i)$,
$\mathcal{I}_i$, $h : (\Sh(\mathcal{D}_1), \mathcal{O}'_1) \to
(\Sh(\mathcal{D}_2), \mathcal{O}'_2)$는 주
\ref{remark-extension-functorial-ringed-topoi}에서와 같다고 하자.
$\pi_1 = h \circ \pi_2$를 만족하는 자명화
$\pi_i : (\Sh(\mathcal{D}_i), \mathcal{O}'_i) \to
(\Sh(\mathcal{C}), \mathcal{O})$가 주어졌다고 하자. 다시 말해 $h$가
$(\Sh(\mathcal{C}), \mathcal{O})$의 자명화된 1차 두꺼워짐 사이의
사상이라고 가정한다. 각 $i = 1, 2$에 대하여
$(\mathcal{K}_i, c_i)$는 $\mathcal{O}$-가군 $\mathcal{K}_i$와 준동형
$c_i : \mathcal{I}_i \otimes_\mathcal{O} \mathcal{F} \to
\mathcal{K}_i$로 이루어졌다고 하자. 또한 다음 도표를 가환시키는
$\mathcal{O}$-가군 준동형 $\mathcal{K}_2 \to \mathcal{K}_1$이
주어졌다고 하자.
$$
\xymatrix{
\mathcal{I}_2 \otimes_\mathcal{O} \mathcal{F}
\ar[r]_-{c_2} \ar[d] &
\mathcal{K}_2 \ar[d] \\
\mathcal{I}_1 \otimes_\mathcal{O} \mathcal{F}
\ar[r]^-{c_1} &
\mathcal{K}_1
}
$$
이 상황에서 주 \ref{remark-trivial-extension-ringed-topoi}의 구성은
다음 가환 도표를 유도한다.
$$
\xymatrix{
\{\mathcal{F}'_2\text{는 (\ref{equation-extension-ringed-topoi})과 같고 }
c_2 = c_{\mathcal{F}'_2}\text{이며 }\mathcal{K} = \mathcal{K}_2\}
\ar[d] \ar[rr] & &
\Ext^1_\mathcal{O}(\mathcal{F}, \mathcal{K}_2) \ar[d] \\
\{\mathcal{F}'_1\text{는 (\ref{equation-extension-ringed-topoi})과 같고 }
c_1 = c_{\mathcal{F}'_1}\text{이며 }\mathcal{K} = \mathcal{K}_1\}
\ar[rr] & &
\Ext^1_\mathcal{O}(\mathcal{F}, \mathcal{K}_1)
}
$$
여기서 오른쪽 세로 준동형은 $\Ext$의 함자성과 준동형
$\mathcal{K}_2 \to \mathcal{K}_1$로 주어지고, 왼쪽 세로 준동형은
주 \ref{remark-extension-functorial-ringed-topoi}에서 구성한 것이다.
\end{remark}

\begin{remark}
\label{remark-obstruction-extension-functorial-ringed-topoi}
$(\Sh(\mathcal{C}), \mathcal{O})$,
$(\Sh(\mathcal{C}), \mathcal{O}) \to (\Sh(\mathcal{D}_i), \mathcal{O}'_i)$,
$\mathcal{I}_i$, $h : (\Sh(\mathcal{D}_1), \mathcal{O}'_1) \to
(\Sh(\mathcal{D}_2), \mathcal{O}'_2)$는 주
\ref{remark-extension-functorial-ringed-topoi}에서와 같다고 하자. 특히
$h^\sharp : \mathcal{O}'_2 \to \mathcal{O}'_1$은 $\mathcal{O}$-가군
준동형 $\mathcal{I}_2 \to \mathcal{I}_1$을 유도한다.
$\mathcal{F}$를 $\mathcal{O}$-가군이라 하자. 각 $i = 1, 2$에 대하여
$(\mathcal{K}_i, c_i)$는 $\mathcal{O}$-가군 $\mathcal{K}_i$와 준동형
$c_i : \mathcal{I}_i \otimes_\mathcal{O} \mathcal{F} \to
\mathcal{K}_i$로 이루어졌다고 하자. 또한 다음 도표를 가환시키는
$\mathcal{O}$-가군 준동형 $\mathcal{K}_2 \to \mathcal{K}_1$이
주어졌다고 하자.
$$
\xymatrix{
\mathcal{I}_2 \otimes_\mathcal{O} \mathcal{F}
\ar[r]_-{c_2} \ar[d] &
\mathcal{K}_2 \ar[d] \\
\mathcal{I}_1 \otimes_\mathcal{O} \mathcal{F}
\ar[r]^-{c_1} &
\mathcal{K}_1
}
$$
그러면 준동형
$$
\Ext^2_\mathcal{O}(\mathcal{F}, \mathcal{K}_2)
\longrightarrow
\Ext^2_\mathcal{O}(\mathcal{F}, \mathcal{K}_1)
$$
이 $o(\mathcal{F}, \mathcal{K}_2, c_2)$를
$o(\mathcal{F}, \mathcal{K}_1, c_1)$로 보낸다고 {\bf 주장}한다.

\medskip\noindent
이 주장을 증명하기 위해 매장
$j_2 : \mathcal{K}_2 \to \mathcal{K}_2'$
을 택하되 $\mathcal{K}_2'$이 단사 $\mathcal{O}$-가군이 되게 한다.
보조정리 \ref{lemma-inf-obs-ext-ringed-topoi}의 증명에서와 같이
다음 $\mathcal{O}_2$-가군의 확대를 택할 수 있다.
$$
0 \to \mathcal{K}_2' \to \mathcal{E}_2 \to \mathcal{F} \to 0
$$
$c_{\mathcal{E}_2} = j_2 \circ c_2$를 만족하도록 택한다.
보조정리 \ref{lemma-inf-obs-ext-ringed-topoi}의 증명은
$o(\mathcal{F}, \mathcal{K}_2, c_2)$
를 다음 $\mathcal{O}$-가군 완전열의 Yoneda 확대 동치류로 구성한다.
여기서 Yoneda 확대는 유도 범주, 절 \ref{derived-section-ext}의 의미이다.
$$
0 \to
\mathcal{K}_2 \to \mathcal{K}_2' \to
\mathcal{E}_2/\mathcal{K}_2 \to
\mathcal{F} \to 0
$$
$\mathcal{K}_1'$을 $\mathcal{K}_2 \to
\mathcal{K}_1 \oplus \mathcal{K}_2'$의 여핵이라 하자. 가환 정사각형을
이루는 단사 준동형 $j_1 : \mathcal{K}_1 \to \mathcal{K}_1'$과 준동형
$\mathcal{K}_2' \to \mathcal{K}_1'$이 있다. 다음 밀어내기를 취한다.
$$
\xymatrix{
0 \ar[r] &
\mathcal{K}_2' \ar[r] \ar[d] &
\mathcal{E}_2 \ar[r] \ar[d] &
\mathcal{F} \ar[r] \ar[d] & 0 \\
0 \ar[r] &
\mathcal{K}_1' \ar[r] &
\mathcal{E}_1 \ar[r] &
\mathcal{F} \ar[r] & 0
}
$$
$\mathcal{E}_1$ 위에는 표준 $\mathcal{O}_1$-가군 구조가 있고, 이 구조에
대해 $c_{\mathcal{E}_1} = j_1 \circ c_1$이다. 여기에는 위의 $c_1$과
$c_2$가 들어가는 도표의 가환성이 사용된다. 보조정리
\ref{lemma-inf-obs-ext-ringed-topoi}의 절차에 따르면
$o(\mathcal{F}, \mathcal{K}_1, c_1)$은 다음 $\mathcal{O}$-가군
완전열의 Yoneda 확대 동치류이다.
$$
0 \to
\mathcal{K}_1 \to
\mathcal{K}_1' \to
\mathcal{E}_1/\mathcal{K}_1 \to
\mathcal{F} \to 0
$$
다음 완전열 사이의 준동형들이 있으므로
$$
\xymatrix{
0 \ar[r] &
\mathcal{K}_2 \ar[d] \ar[r] &
\mathcal{K}_2' \ar[d] \ar[r] &
\mathcal{E}_2/\mathcal{K}_2 \ar[r] \ar[d] &
\mathcal{F} \ar[r] \ar@{=}[d] &
0 \\
0 \ar[r] &
\mathcal{K}_2 \ar[r] &
\mathcal{K}_2' \ar[r] &
\mathcal{E}_2/\mathcal{K}_2 \ar[r] &
\mathcal{F} \ar[r] &
0
}
$$
주장이 참임을 알 수 있다.
\end{remark}

\begin{remark}
\label{remark-short-exact-sequence-thickenings-ringed-topoi}
$(\Sh(\mathcal{C}), \mathcal{O})$를 환 달린 토포스라 하자.
$(\Sh(\mathcal{C}), \mathcal{O})$의 1차 두꺼워짐 사이의 사상열
$$
(\Sh(\mathcal{D}_1), \mathcal{O}'_1) \to
(\Sh(\mathcal{D}_2), \mathcal{O}'_2) \to
(\Sh(\mathcal{D}_3), \mathcal{O}'_3)
$$
에서 아이디얼층 $\mathcal{I}_i$ 사이의 대응하는 준동형들이
$\mathcal{O}$-가군의 복합체
$\mathcal{I}_3 \to \mathcal{I}_2 \to \mathcal{I}_1$, 즉 합성이 영인 열을
이루면, 주어진 두꺼워짐의 열을 {\it 복합체}라고 한다. 이 경우 합성
$(\Sh(\mathcal{D}_1), \mathcal{O}'_1) \to
(\Sh(\mathcal{D}_3), \mathcal{O}'_3)$은
$(\Sh(\mathcal{C}), \mathcal{O}) \to
(\Sh(\mathcal{D}_3), \mathcal{O}'_3)$을 거쳐 분해된다. 즉
$(\Sh(\mathcal{C}), \mathcal{O})$의 1차 두꺼워짐
$(\Sh(\mathcal{D}_1), \mathcal{O}'_1)$은 자명하며 표준 자명화
$\pi : (\Sh(\mathcal{D}_1), \mathcal{O}'_1) \to
(\Sh(\mathcal{C}), \mathcal{O})$.

\medskip\noindent
$(\Sh(\mathcal{C}), \mathcal{O})$의 1차 두꺼워짐 사이의 사상열
$$
(\Sh(\mathcal{D}_1), \mathcal{O}'_1) \to
(\Sh(\mathcal{D}_2), \mathcal{O}'_2) \to
(\Sh(\mathcal{D}_3), \mathcal{O}'_3)
$$
에서 아이디얼층 사이의 대응하는 준동형들이 $\mathcal{O}$-가군의
짧은 완전열
$$
0 \to \mathcal{I}_3 \to \mathcal{I}_2 \to \mathcal{I}_1 \to 0
$$
을 이루면, 주어진 두꺼워짐의 열을 {\it 짧은 완전열}이라고 한다.
\end{remark}

\begin{remark}
\label{remark-complex-thickenings-and-ses-modules-ringed-topoi}
$(\Sh(\mathcal{C}), \mathcal{O})$를 환 달린 토포스라 하고
$\mathcal{F}$를 $\mathcal{O}$-가군이라 하자. 다음 열을
$$
(\Sh(\mathcal{D}_1), \mathcal{O}'_1) \to
(\Sh(\mathcal{D}_2), \mathcal{O}'_2) \to
(\Sh(\mathcal{D}_3), \mathcal{O}'_3)
$$
$(\Sh(\mathcal{C}), \mathcal{O})$의 1차 두꺼워짐의 복합체라 하자. 주
\ref{remark-short-exact-sequence-thickenings-ringed-topoi}를 보라. 각
$i = 1, 2, 3$에 대하여 $(\mathcal{K}_i, c_i)$는 $\mathcal{O}$-가군
$\mathcal{K}_i$와 준동형
$c_i : \mathcal{I}_i \otimes_\mathcal{O} \mathcal{F} \to
\mathcal{K}_i$로 이루어졌다고 하자. 다음과 같은 $\mathcal{O}$-가군의
짧은 완전열이 주어졌다고 하자.
$$
0 \to \mathcal{K}_3 \to \mathcal{K}_2 \to \mathcal{K}_1 \to 0
$$
또 다음 두 도표가 가환한다고 가정한다.
$$
\vcenter{
\xymatrix{
\mathcal{I}_2 \otimes_\mathcal{O} \mathcal{F}
\ar[r]_-{c_2} \ar[d] &
\mathcal{K}_2 \ar[d] \\
\mathcal{I}_1 \otimes_\mathcal{O} \mathcal{F}
\ar[r]^-{c_1} &
\mathcal{K}_1
}
}
\quad\text{그리고}\quad
\vcenter{
\xymatrix{
\mathcal{I}_3 \otimes_\mathcal{O} \mathcal{F}
\ar[r]_-{c_3} \ar[d] &
\mathcal{K}_3 \ar[d] \\
\mathcal{I}_2 \otimes_\mathcal{O} \mathcal{F}
\ar[r]^-{c_2} &
\mathcal{K}_2
}
}
$$
끝으로 $\mathcal{O}'_2$-가군의 확대
$$
0 \to \mathcal{K}_2 \to \mathcal{F}'_2 \to \mathcal{F} \to 0
$$
이 주어졌다고 하자. 이는 $\mathcal{K} = \mathcal{K}_2$인
(\ref{equation-extension-ringed-topoi})과 같고
$c_{\mathcal{F}'_2} = c_2$이다. 주
\ref{remark-extension-functorial-ringed-topoi}의 함자성을 적용하여
$\mathcal{O}'_1$-가군의 확대 $\mathcal{F}'_1$을 얻을 수 있다. 이 특별한
경우의 $\mathcal{F}'_1$은 아래에서 설명한다. 주
\ref{remark-trivial-extension-ringed-topoi}에 따라, 주
\ref{remark-short-exact-sequence-thickenings-ringed-topoi}의 표준 자명화
$\pi : (\Sh(\mathcal{D}_1), \mathcal{O}'_1) \to
(\Sh(\mathcal{C}), \mathcal{O})$를 사용하면
$\xi_{\mathcal{F}'_1} \in
\Ext^1_\mathcal{O}(\mathcal{F}, \mathcal{K}_1)$을 얻는다. 마지막으로
다음 장애 원소가 있다.
$$
o(\mathcal{F}, \mathcal{K}_3, c_3) \in
\Ext^2_\mathcal{O}(\mathcal{F}, \mathcal{K}_3)
$$
보조정리 \ref{lemma-inf-obs-ext-ringed-topoi}를 보라. 이 상황에서 짧은
완전열
$0 \to \mathcal{K}_3 \to \mathcal{K}_2 \to \mathcal{K}_1 \to 0$에서
오는 표준 준동형
$$
\partial :
\Ext^1_\mathcal{O}(\mathcal{F}, \mathcal{K}_1)
\longrightarrow
\Ext^2_\mathcal{O}(\mathcal{F}, \mathcal{K}_3)
$$
이 $\xi_{\mathcal{F}'_1}$을 장애 동치류
$o(\mathcal{F}, \mathcal{K}_3, c_3)$로 보낸다고 {\bf 주장}한다.

\medskip\noindent
이 주장을 증명하기 위해 $\mathcal{K}$가 단사 $\mathcal{O}$-가군이
되도록 매장 $j : \mathcal{K}_3 \to \mathcal{K}$를 택한다. $j$를 준동형
$j' : \mathcal{K}_2 \to \mathcal{K}$로 들어 올릴 수 있다.
$\mathcal{E}'_2 = j'_*\mathcal{F}'_2$를 $j'$에 따른
$\mathcal{F}'_2$의 밀어내기로 놓는다. 그러면
$c_{\mathcal{E}'_2} = j' \circ c_2$이다. 도표로 쓰면 다음과 같다.
$$
\xymatrix{
0 \ar[r] &
\mathcal{K}_2 \ar[r] \ar[d]_{j'} &
\mathcal{F}'_2 \ar[r] \ar[d] &
\mathcal{F} \ar[r] \ar[d] & 0 \\
0 \ar[r] &
\mathcal{K} \ar[r] &
\mathcal{E}'_2 \ar[r] &
\mathcal{F} \ar[r] & 0
}
$$
$\mathcal{E}'_3 = \mathcal{E}'_2$로 놓되,
$\mathcal{O}'_3 \to \mathcal{O}'_2$를 통해 $\mathcal{O}'_3$-가군으로
본다. 그러면 $c_{\mathcal{E}'_3} = j \circ c_3$이다. 보조정리
\ref{lemma-inf-obs-ext-ringed-topoi}의 증명은
$o(\mathcal{F}, \mathcal{K}_3, c_3)$
를 다음 $\mathcal{O}$-가군 확대의 동치류의 경계로 구성한다.
$$
0 \to
\mathcal{K}/\mathcal{K}_3 \to
\mathcal{E}'_3/\mathcal{K}_3 \to
\mathcal{F} \to 0
$$
한편 $\mathcal{F}'_1 = \mathcal{F}'_2/\mathcal{K}_3$임에 유의하자.
따라서 $\xi_{\mathcal{F}'_1}$은 다음 확대의 동치류이다.
$$
0 \to \mathcal{K}_2/\mathcal{K}_3 \to \mathcal{F}'_2/\mathcal{K}_3
\to \mathcal{F} \to 0
$$
여기서는 표준 자명화
$\pi : (\Sh(\mathcal{D}_1), \mathcal{O}'_1) \to
(\Sh(\mathcal{C}), \mathcal{O})$의 $\pi^\sharp$를 사용하여 이를
$\mathcal{O}$-가군의 열로 본다. 이제 다음 가환 도표가 있으므로 주장이
따른다.
$$
\xymatrix{
0 \ar[r] &
\mathcal{K}_2/\mathcal{K}_3 \ar[r] \ar[d] &
\mathcal{F}'_2/\mathcal{K}_3 \ar[r] \ar[d] &
\mathcal{F} \ar[r] \ar[d] & 0 \\
0 \ar[r] &
\mathcal{K}/\mathcal{K}_3 \ar[r] &
\mathcal{E}'_3/\mathcal{K}_3 \ar[r] &
\mathcal{F} \ar[r] & 0
}
$$
이 도표는 위에서 주어진 $\mathcal{O}$-가군 구조들에 대해
$\mathcal{O}$-선형이다.
\end{remark}







\section{환 달린 토포스 위 가군의 무한소 변형}
\label{section-deformation-modules-ringed-topoi}

\noindent
$i : (\Sh(\mathcal{C}), \mathcal{O}) \to (\Sh(\mathcal{D}), \mathcal{O}')$
를 환 달린 토포스의 1차 두꺼워짐이라 하자. 절
\ref{section-thickenings-ringed-topoi}에서 도입한 표기를 자유롭게
사용한다. $\mathcal{F}'$를 $\mathcal{O}'$-가군이라 하고
$\mathcal{F} = i^*\mathcal{F}'$라 놓자. 이 상황에서는 다음 짧은
완전열이 있다.
$$
0 \to \mathcal{I}\mathcal{F}' \to \mathcal{F}' \to \mathcal{F} \to 0
$$
이는 $\mathcal{O}'$-가군의 열이다. $\mathcal{I}^2 = 0$이므로
$\mathcal{I}\mathcal{F}'$ 위의 $\mathcal{O}'$-가군 구조는 유일한
$\mathcal{O}$-가군 구조에서 온다. 따라서 위의 열은
(\ref{equation-extension-ringed-topoi})과 같은 확대이다. 특히
$\mathcal{F}' = \mathcal{O}'$이면 $i^*\mathcal{O}' = \mathcal{O}$이고
$\mathcal{I}\mathcal{O}' = \mathcal{I}$이므로 다음 구조층의 열을 되찾는다.
$$
0 \to \mathcal{I} \to \mathcal{O}' \to \mathcal{O} \to 0
$$

\begin{lemma}
\label{lemma-inf-map-special-ringed-topoi}
$i : (\Sh(\mathcal{C}), \mathcal{O}) \to (\Sh(\mathcal{D}), \mathcal{O}')$
를 환 달린 토포스의 1차 두꺼워짐이라 하자. $\mathcal{F}'$,
$\mathcal{G}'$를 $\mathcal{O}'$-가군이라 하고
$\mathcal{F} = i^*\mathcal{F}'$ 및 $\mathcal{G} = i^*\mathcal{G}'$라
놓자. $\varphi : \mathcal{F} \to \mathcal{G}$를 $\mathcal{O}$-선형
준동형이라 하자. $\varphi$를 $\mathcal{O}'$-선형 준동형
$\varphi' : \mathcal{F}' \to \mathcal{G}'$로 올리는 방법들의 집합은
공집합이 아니면 다음 군 아래의 주동차 공간이다.
$\Hom_\mathcal{O}(\mathcal{F}, \mathcal{I}\mathcal{G}')$.
\end{lemma}

\begin{proof}
이는 보조정리 \ref{lemma-inf-map-ringed-topoi}의 특수한 경우이지만 직접
증명도 제시한다. 다음 가군의 짧은 완전열들이 있다.
$$
0 \to \mathcal{I} \to \mathcal{O}' \to \mathcal{O} \to 0
\quad\text{그리고}\quad
0 \to \mathcal{I}\mathcal{G}' \to \mathcal{G}' \to \mathcal{G} \to 0
$$
$\mathcal{F}'$에 대해서도 마찬가지이다. $\mathcal{I}$의 제곱이 영이므로
$\mathcal{I}$와 $\mathcal{I}\mathcal{G}'$ 위의 $\mathcal{O}'$-가군
구조는 유일한 $\mathcal{O}$-가군 구조에서 온다. 따라서
$$
\Hom_{\mathcal{O}'}(\mathcal{F}', \mathcal{I}\mathcal{G}') =
\Hom_\mathcal{O}(\mathcal{F}, \mathcal{I}\mathcal{G}')
\quad\text{그리고}\quad
\Hom_{\mathcal{O}'}(\mathcal{F}', \mathcal{G}) =
\Hom_\mathcal{O}(\mathcal{F}, \mathcal{G})
$$
이고, 이제 보조정리는 다음 완전열에서 따른다.
$$
0 \to \Hom_{\mathcal{O}'}(\mathcal{F}', \mathcal{I}\mathcal{G}') \to
\Hom_{\mathcal{O}'}(\mathcal{F}', \mathcal{G}') \to
\Hom_{\mathcal{O}'}(\mathcal{F}', \mathcal{G})
$$
호몰로지 대수학, 보조정리 \ref{homology-lemma-check-exactness}를 보라.
\end{proof}

\begin{lemma}
\label{lemma-deform-module-ringed-topoi}
$(f, f')$를 상황 \ref{situation-morphism-thickenings-ringed-topoi}에서와
같은 환 달린 토포스의 1차 두꺼워짐 사이의 사상이라 하자.
$\mathcal{F}'$를 $\mathcal{O}'$-가군이라 하고
$\mathcal{F} = i^*\mathcal{F}'$라 놓자. $\mathcal{F}$가
$\mathcal{O}_\mathcal{B}$ 위에서 평탄하고 $(f, f')$가 두꺼워짐의 엄격
사상이라고 가정하자(정의
\ref{definition-strict-morphism-thickenings-ringed-topoi}). 그러면 다음
조건들이 동치이다.
\begin{enumerate}
\item $\mathcal{F}'$는 $\mathcal{O}_{\mathcal{B}'}$ 위에서 평탄하다.
\item 표준 준동형
$f^*\mathcal{J} \otimes_\mathcal{O} \mathcal{F} \to
\mathcal{I}\mathcal{F}'$
은 동형사상이다.
\end{enumerate}
더욱이 이 경우 다음 준동형들은 동형사상이다.
$$
f^*\mathcal{J} \otimes_\mathcal{O} \mathcal{F} \to
\mathcal{I} \otimes_\mathcal{O} \mathcal{F} \to
\mathcal{I}\mathcal{F}'
$$
는 모두 동형사상이다.
\end{lemma}

\begin{proof}
$(f, f')$가 두꺼워짐의 엄격 사상이므로 준동형
$f^*\mathcal{J} \to \mathcal{I}$은 전사이다. 따라서 마지막 명제는
(2)에서 따른다.

\medskip\noindent
(1)과 (2)의 동치를 증명하자. 정의상 $\mathcal{O}_\mathcal{B}$ 위의
평탄성은 $f^{-1}\mathcal{O}_\mathcal{B}$ 위의 평탄성을 뜻한다.
$f^{-1}\mathcal{O}_{\mathcal{B}'}$ 위의 평탄성도 마찬가지이다.
$(f, f')$의 엄격성과 $\mathcal{F} = i^*\mathcal{F}'$이라는 가정으로부터
$$
\mathcal{F} = \mathcal{F}'/(f^{-1}\mathcal{J})\mathcal{F}'
$$
라는 $\mathcal{C}$ 위의 층의 등식이 따른다. 또한
$f^*\mathcal{J} \otimes_\mathcal{O} \mathcal{F} =
f^{-1}\mathcal{J} \otimes_{f^{-1}\mathcal{O}_\mathcal{B}} \mathcal{F}$이다.
따라서 (1)과 (2)의 동치는 사이트 위의 가군, 보조정리
\ref{sites-modules-lemma-flat-over-thickening}에서 따른다.
\end{proof}

\begin{lemma}
\label{lemma-deform-fp-module-ringed-topoi}
$(f, f')$를 상황 \ref{situation-morphism-thickenings-ringed-topoi}에서와
같은 환 달린 토포스의 1차 두꺼워짐 사이의 사상이라 하자.
$\mathcal{F}'$를 $\mathcal{O}'$-가군이라 하고
$\mathcal{F} = i^*\mathcal{F}'$라 놓자. $\mathcal{F}'$가
$\mathcal{O}_{\mathcal{B}'}$ 위에서 평탄하고 $(f, f')$가 두꺼워짐의
엄격 사상이라고 가정하자. 그러면 다음 조건들이 동치이다.
\begin{enumerate}
\item $\mathcal{F}'$는 유한 표시 $\mathcal{O}'$-가군이다.
\item $\mathcal{F}$는 유한 표시 $\mathcal{O}$-가군이다.
\end{enumerate}
\end{lemma}

\begin{proof}
(1) $\Rightarrow$ (2)는 사이트 위의 가군, 보조정리
\ref{sites-modules-lemma-local-pullback}에서 따른다. 역으로
$\mathcal{F}$가 유한 표시라고 가정하자. $\mathcal{C} = \mathcal{C}'$라고
가정해도 좋고 그렇게 한다. 보조정리
\ref{lemma-deform-module-ringed-topoi}에 의해 다음 짧은 완전열이 있다.
$$
0 \to \mathcal{I} \otimes_{\mathcal{O}_X} \mathcal{F} \to
\mathcal{F}' \to \mathcal{F} \to 0
$$
$U$를 $\mathcal{F}|_U$가 다음 표시를 갖는 $\mathcal{C}$의 대상이라 하자.
$$
\mathcal{O}_U^{\oplus m} \to \mathcal{O}_U^{\oplus n} \to \mathcal{F}|_U \to 0
$$
덮개의 원소들로 $U$를 대체하여 준동형
$\mathcal{O}_U^{\oplus n} \to \mathcal{F}|_U$가 준동형
$(\mathcal{O}'_U)^{\oplus n} \to \mathcal{F}'|_U$로 올라간다고 가정할 수
있다. 유도된 준동형 $\mathcal{I}^{\oplus n} \to
\mathcal{I} \otimes \mathcal{F}$는 $\otimes$의 오른쪽 완전성에 의해
전사이다. 따라서 덮개의 원소들로 $U$를 다시 대체하면, 주어진 준동형
$\mathcal{O}_U^{\oplus m} \to \mathcal{O}_U^{\oplus n}$의 올림
$(\mathcal{O}'|_U)^{\oplus m} \to (\mathcal{O}'|_U)^{\oplus n}$
을 택하여
$$
(\mathcal{O}'_U)^{\oplus m} \to (\mathcal{O}'_U)^{\oplus n} \to
\mathcal{F}'|_U \to 0
$$
이 복합체가 되도록 할 수 있다. $\otimes$의 오른쪽 완전성을 다시
사용하면 이 복합체가 완전함을 알 수 있다.
\end{proof}

\begin{lemma}
\label{lemma-inf-map-rel-ringed-topoi}
$(f, f')$를 상황 \ref{situation-morphism-thickenings-ringed-topoi}에서와
같은 1차 두꺼워짐 사이의 사상이라 하자. $\mathcal{F}'$,
$\mathcal{G}'$를 $\mathcal{O}'$-가군이라 하고
$\mathcal{F} = i^*\mathcal{F}'$ 및 $\mathcal{G} = i^*\mathcal{G}'$라
놓자. $\varphi : \mathcal{F} \to \mathcal{G}$를 $\mathcal{O}$-선형
준동형이라 하자. $\mathcal{G}'$가 $\mathcal{O}_{\mathcal{B}'}$ 위에서
평탄하고 $(f, f')$가 두꺼워짐의 엄격 사상이라고 가정하자.
$\varphi$를 $\mathcal{O}'$-선형 준동형
$\varphi' : \mathcal{F}' \to \mathcal{G}'$로 올리는 방법들의 집합은
공집합이 아니면 다음 군 아래의 주동차 공간이다.
$$
\Hom_\mathcal{O}(\mathcal{F},
\mathcal{G} \otimes_\mathcal{O} f^*\mathcal{J})
$$
\end{lemma}

\begin{proof}
보조정리 \ref{lemma-inf-map-special-ringed-topoi}와
\ref{lemma-deform-module-ringed-topoi}을 결합하면 된다.
\end{proof}

\begin{lemma}
\label{lemma-inf-obs-map-special-ringed-topoi}
$i : (\Sh(\mathcal{C}), \mathcal{O}) \to (\Sh(\mathcal{D}), \mathcal{O}')$
를 환 달린 토포스의 1차 두꺼워짐이라 하자. $\mathcal{F}'$,
$\mathcal{G}'$를 $\mathcal{O}'$-가군이라 하고
$\mathcal{F} = i^*\mathcal{F}'$ 및 $\mathcal{G} = i^*\mathcal{G}'$라
놓자. $\varphi : \mathcal{F} \to \mathcal{G}$를 $\mathcal{O}$-선형
준동형이라 하자. 다음 원소가 존재한다.
$$
o(\varphi) \in
\Ext^1_\mathcal{O}(Li^*\mathcal{F}', \mathcal{I}\mathcal{G}')
$$
이 원소가 영인 것은 $\varphi$를 $\mathcal{O}'$-선형 준동형
$\varphi' : \mathcal{F}' \to \mathcal{G}'$로 올릴 수 있기 위한
필요충분조건이다.
\end{lemma}

\begin{proof}
보조정리 \ref{lemma-inf-map-special-ringed-topoi}의 증명으로부터, 다음
준동형에 따른 $\varphi$의 경계가 영인 것이 올림의 존재를 위한
필요충분조건임을 알 수 있다.
$$
\Hom_\mathcal{O}(\mathcal{F}, \mathcal{G}) =
\Hom_{\mathcal{O}'}(\mathcal{F}', \mathcal{G}) \longrightarrow
\Ext^1_{\mathcal{O}'}(\mathcal{F}', \mathcal{I}\mathcal{G}')
$$
한편
$$
\Ext^1_{\mathcal{O}'}(\mathcal{F}', \mathcal{I}\mathcal{G}') =
\Ext^1_\mathcal{O}(Li^*\mathcal{F}', \mathcal{I}\mathcal{G}')
$$
는 유도 범주에서의 $i_* = Ri_*$와 $Li^*$의 수반성에서 따른다
(사이트 위의 코호몰로지, 보조정리
\ref{sites-cohomology-lemma-adjoint}).
\end{proof}

\begin{lemma}
\label{lemma-inf-obs-map-rel-ringed-topoi}
$(f, f')$를 상황 \ref{situation-morphism-thickenings-ringed-topoi}에서와
같은 1차 두꺼워짐 사이의 사상이라 하자. $\mathcal{F}'$,
$\mathcal{G}'$를 $\mathcal{O}'$-가군이라 하고
$\mathcal{F} = i^*\mathcal{F}'$ 및 $\mathcal{G} = i^*\mathcal{G}'$라
놓자. $\varphi : \mathcal{F} \to \mathcal{G}$를 $\mathcal{O}$-선형
준동형이라 하자. $\mathcal{F}'$와 $\mathcal{G}'$가
$\mathcal{O}_{\mathcal{B}'}$ 위에서 평탄하고 $(f, f')$가 두꺼워짐의
엄격 사상이라고 가정하자. 다음 원소가 존재한다.
$$
o(\varphi) \in
\Ext^1_\mathcal{O}(\mathcal{F},
\mathcal{G} \otimes_\mathcal{O} f^*\mathcal{J})
$$
이 원소가 영인 것은 $\varphi$를 $\mathcal{O}'$-선형 준동형
$\varphi' : \mathcal{F}' \to \mathcal{G}'$로 올릴 수 있기 위한
필요충분조건이다.
\end{lemma}

\begin{proof}[첫째 증명]
보조정리의 가정 아래 다음 등식이 성립한다고 주장하면, 결과는 보조정리
\ref{lemma-inf-obs-map-special-ringed-topoi}에서 따른다.
$$
\Ext^1_\mathcal{O}(Li^*\mathcal{F}', \mathcal{I}\mathcal{G}') =
\Ext^1_\mathcal{O}(\mathcal{F},
\mathcal{G} \otimes_\mathcal{O} f^*\mathcal{J})
$$
실제로 보조정리 \ref{lemma-deform-module-ringed-topoi}에 의해
$\mathcal{I}\mathcal{G}' =
\mathcal{G} \otimes_\mathcal{O} f^*\mathcal{J}$
이다. 한편 다음에 유의하자.
$$
H^{-1}(Li^*\mathcal{F}') =
\text{Tor}_1^{\mathcal{O}'}(\mathcal{F}', \mathcal{O})
$$
(국소 계산은 생략한다). 다음 짧은 완전열을 사용하면
$$
0 \to \mathcal{I} \to \mathcal{O}' \to \mathcal{O} \to 0
$$
이 $\text{Tor}_1$은 준동형
$\mathcal{I} \otimes_\mathcal{O} \mathcal{F} \to \mathcal{I}\mathcal{F}'$
의 핵으로 계산됨을 알 수 있다. 이 핵은 보조정리
\ref{lemma-deform-module-ringed-topoi}의 마지막 명제에 의해 영이다.
따라서 $\tau_{\geq -1}Li^*\mathcal{F}' = \mathcal{F}$이다. 한편
$$
\Ext^1_\mathcal{O}(Li^*\mathcal{F}',
\mathcal{I}\mathcal{G}') =
\Ext^1_\mathcal{O}(\tau_{\geq -1}Li^*\mathcal{F}',
\mathcal{I}\mathcal{G}')
$$
는 유도 범주, 보조정리
\ref{derived-lemma-negative-vanishing}의 쌍대형에서 따른다.
\end{proof}

\begin{proof}[둘째 증명]
보조정리 \ref{lemma-inf-obs-map-ringed-topoi}를 다음과 같이 적용할 수 있다.
보조정리 \ref{lemma-deform-module-ringed-topoi}에 의해
$\mathcal{K} = \mathcal{I} \otimes_\mathcal{O} \mathcal{F}$이고
$\mathcal{L} = \mathcal{I} \otimes_\mathcal{O} \mathcal{G}$
이고 $c_{\mathcal{F}'} = 1 \otimes 1$이고
$c_{\mathcal{G}'} = 1 \otimes 1$이다. $\psi = 1 \otimes \varphi$로
택하면 보조정리의 도표가 가환한다. 따라서
$o(\varphi) = o(\varphi, 1 \otimes \varphi)$로 택하면 된다.
\end{proof}

\begin{lemma}
\label{lemma-inf-ext-rel-ringed-topoi}
$(f, f')$를 상황 \ref{situation-morphism-thickenings-ringed-topoi}에서와
같은 1차 두꺼워짐 사이의 사상이라 하자. $\mathcal{F}$를
$\mathcal{O}$-가군이라 하자. $(f, f')$가 두꺼워짐의 엄격 사상이고
$\mathcal{F}$가 $\mathcal{O}_\mathcal{B}$ 위에서 평탄하다고 가정하자.
$\mathcal{O}_{\mathcal{B}'}$ 위에서 평탄한 $\mathcal{O}'$-가군
$\mathcal{F}'$와 동형사상 $\alpha : i^*\mathcal{F}' \to \mathcal{F}$로
이루어진 쌍 $(\mathcal{F}', \alpha)$가 하나라도 존재하면, 그러한 쌍들의
동형류 집합은 다음 군 아래의 주동차 공간이다.
$\Ext^1_\mathcal{O}(
\mathcal{F}, \mathcal{I} \otimes_\mathcal{O} \mathcal{F})$.
\end{lemma}

\begin{proof}
그러한 가군이 하나 존재한다고 가정하면, 보조정리
\ref{lemma-deform-module-ringed-topoi}에 의해 표준 준동형
$$
f^*\mathcal{J} \otimes_\mathcal{O} \mathcal{F} \to
\mathcal{I} \otimes_\mathcal{O} \mathcal{F}
$$
은 동형사상이다. $\mathcal{K} = 
\mathcal{I} \otimes_\mathcal{O} \mathcal{F}$ 및 $c = 1$로 놓고 보조정리
\ref{lemma-inf-ext-ringed-topoi}를 적용하자. 보조정리
\ref{lemma-deform-module-ringed-topoi}에 의해 대응하는 확대
$\mathcal{F}'$는 모두 $\mathcal{O}_{\mathcal{B}'}$ 위에서 평탄하다.
\end{proof}

\begin{lemma}
\label{lemma-inf-obs-ext-rel-ringed-topoi}
$(f, f')$를 상황 \ref{situation-morphism-thickenings-ringed-topoi}에서와
같은 1차 두꺼워짐 사이의 사상이라 하자. $\mathcal{F}$를
$\mathcal{O}$-가군이라 하자. $(f, f')$가 두꺼워짐의 엄격 사상이고
$\mathcal{F}$가 $\mathcal{O}_\mathcal{B}$ 위에서 평탄하다고 가정하자.
$\mathcal{O}_{\mathcal{B}'}$ 위에서 평탄하며
$i^*\mathcal{F}' \cong \mathcal{F}$를 만족하는 $\mathcal{O}'$-가군
$\mathcal{F}'$가 존재하기 위한 필요충분조건은 다음과 같다.
\begin{enumerate}
\item 표준 준동형
$f^*\mathcal{J} \otimes_\mathcal{O} \mathcal{F} \to
\mathcal{I} \otimes_\mathcal{O} \mathcal{F}$
은 동형사상이다.
\item 보조정리 \ref{lemma-inf-obs-ext-ringed-topoi}의 동치류
$o(\mathcal{F}, \mathcal{I} \otimes_\mathcal{O} \mathcal{F}, 1)
\in \Ext^2_\mathcal{O}(
\mathcal{F}, \mathcal{I} \otimes_\mathcal{O} \mathcal{F})$
는 영이다.
\end{enumerate}
\end{lemma}

\begin{proof}
이는 보조정리 \ref{lemma-deform-module-ringed-topoi}에 주어진
$\mathcal{O}_{\mathcal{B}'}$ 위에서 평탄한 $\mathcal{O}'$-가군의
특징짓기와 보조정리 \ref{lemma-inf-obs-ext-ringed-topoi}에서 바로 따른다.
\end{proof}






\section{환 달린 토포스의 평탄한 두꺼워짐 위 평탄한 가군에 대한 응용}
\label{section-flat-ringed-topoi}

\noindent
다음 환 달린 토포스의 가환 도표를 생각하자.
$$
\xymatrix{
(\Sh(\mathcal{C}), \mathcal{O}) \ar[r]_i \ar[d]_f &
(\Sh(\mathcal{D}), \mathcal{O}') \ar[d]^{f'} \\
(\Sh(\mathcal{B}), \mathcal{O}_\mathcal{B}) \ar[r]^t &
(\Sh(\mathcal{B}'), \mathcal{O}_{\mathcal{B}'})
}
$$
가로 화살표들은 상황
\ref{situation-morphism-thickenings-ringed-topoi}에서와 같은 1차
두꺼워짐이다.
$\mathcal{I} = \Ker(i^\sharp) \subset \mathcal{O}'$ 및
$\mathcal{J} = \Ker(t^\sharp) \subset \mathcal{O}_{\mathcal{B}'}$라 놓자.
$\mathcal{F}$를 $\mathcal{O}$-가군이라 하고 다음을 가정하자.
\begin{enumerate}
\item $(f, f')$는 두꺼워짐의 엄격 사상이다.
\item $f'$는 평탄하다.
\item $\mathcal{F}$는 $\mathcal{O}_\mathcal{B}$ 위에서 평탄하다.
\end{enumerate}
(1) $+$ (2)로부터 $\mathcal{I} = f^*\mathcal{J}$가 따름에 유의하자
(보조정리 \ref{lemma-deform-module-ringed-topoi}를 $\mathcal{O}'$에
적용한다). 이 가정들 아래에서는 앞 절의 이론이 특히 간결해진다.
지금까지 얻은 결과를 다음 보조정리로 요약한다.

\begin{lemma}
\label{lemma-flat-ringed-topoi}
위 상황에서 다음이 성립한다.
\begin{enumerate}
\item $\mathcal{O}_{\mathcal{B}'}$ 위에서 평탄하고
$i^*\mathcal{F}' \cong \mathcal{F}$를 만족하는 $\mathcal{O}'$-가군
$\mathcal{F}'$가 존재하기 위한 필요충분조건은 보조정리
\ref{lemma-inf-obs-ext-ringed-topoi}의 동치류
$o(\mathcal{F}, f^*\mathcal{J} \otimes_\mathcal{O} \mathcal{F}, 1)
\in \Ext^2_\mathcal{O}(
\mathcal{F}, f^*\mathcal{J} \otimes_\mathcal{O} \mathcal{F})$
가 영인 것이다.
\item 그러한 가군이 존재하면 올림들의 동형류 집합은 다음 군 아래의
주동차 공간이다.
$\Ext^1_\mathcal{O}(
\mathcal{F}, f^*\mathcal{J} \otimes_\mathcal{O} \mathcal{F})$.
\item 올림 $\mathcal{F}'$가 주어졌을 때, 당기면
$\text{id}_\mathcal{F}$가 되는 $\mathcal{F}'$의 자기동형사상들의 집합은
다음 군과 표준적으로 동형이다.
$\Ext^0_\mathcal{O}(
\mathcal{F}, f^*\mathcal{J} \otimes_\mathcal{O} \mathcal{F})$.
\end{enumerate}
\end{lemma}

\begin{proof}
위에서 $\mathcal{I} = f^*\mathcal{J}$임을 보았으므로 (1)은 보조정리
\ref{lemma-inf-obs-ext-rel-ringed-topoi}에서 따른다. (2)는 보조정리
\ref{lemma-inf-ext-rel-ringed-topoi}에서, (3)은 보조정리
\ref{lemma-inf-map-rel-ringed-topoi}에서 따른다.
\end{proof}

\begin{situation}
\label{situation-morphism-flat-thickenings-ringed-topoi}
$f : (\Sh(\mathcal{C}), \mathcal{O}) \to
(\Sh(\mathcal{B}), \mathcal{O}_\mathcal{B})$를 환 달린 토포스의 사상이라
하자. 다음 가환 도표를 생각하자.
$$
\xymatrix{
(\Sh(\mathcal{C}'_1), \mathcal{O}'_1) \ar[r]_h \ar[d]_{f'_1} &
(\Sh(\mathcal{C}'_2), \mathcal{O}'_2) \ar[d]_{f'_2} \\
(\Sh(\mathcal{B}'_1), \mathcal{O}_{\mathcal{B}'_1}) \ar[r] &
(\Sh(\mathcal{B}'_2), \mathcal{O}_{\mathcal{B}'_2})
}
$$
여기서 $h$는 $(\Sh(\mathcal{C}), \mathcal{O})$의 1차 두꺼워짐 사이의
사상이고, 아래 가로 화살표는
$(\Sh(\mathcal{B}), \mathcal{O}_\mathcal{B})$의 1차 두꺼워짐 사이의
사상이며, 각 $f'_i$의 제한은 $f$이고, 두 쌍 $(f, f_i')$는 모두
두꺼워짐의 엄격 사상이며, 두 $f'_i$는 모두 평탄하다. 끝으로
$\mathcal{F}$를 $\mathcal{O}_\mathcal{B}$ 위에서 평탄한
$\mathcal{O}$-가군이라 하자.
\end{situation}

\begin{lemma}
\label{lemma-functorial-ringed-topoi}
상황 \ref{situation-morphism-flat-thickenings-ringed-topoi}에서 장애 동치류
$o(\mathcal{F}, f^*\mathcal{J}_2 \otimes_\mathcal{O} \mathcal{F}, 1)$
는 장애 동치류
$o(\mathcal{F}, f^*\mathcal{J}_1 \otimes_\mathcal{O} \mathcal{F}, 1)$
로 가며, 이때 표준 준동형은 다음과 같다.
$$
\Ext^2_\mathcal{O}(
\mathcal{F}, f^*\mathcal{J}_2 \otimes_\mathcal{O} \mathcal{F})
\to \Ext^2_\mathcal{O}(
\mathcal{F}, f^*\mathcal{J}_1 \otimes_\mathcal{O} \mathcal{F})
$$
\end{lemma}

\begin{proof}
주 \ref{remark-obstruction-extension-functorial-ringed-topoi}에서 따른다.
\end{proof}

\begin{situation}
\label{situation-ses-flat-thickenings-ringed-topoi}
$f : (\Sh(\mathcal{C}), \mathcal{O}) \to
(\Sh(\mathcal{B}), \mathcal{O}_\mathcal{B})$를 환 달린 토포스의 사상이라
하자. 다음 가환 도표를 생각하자.
$$
\xymatrix{
(\Sh(\mathcal{C}'_1), \mathcal{O}'_1) \ar[r]_h \ar[d]_{f'_1} &
(\Sh(\mathcal{C}'_2), \mathcal{O}'_2) \ar[r] \ar[d]_{f'_2} &
(\Sh(\mathcal{C}'_3), \mathcal{O}'_3) \ar[d]_{f'_3} \\
(\Sh(\mathcal{B}'_1), \mathcal{O}_{\mathcal{B}'_1}) \ar[r] &
(\Sh(\mathcal{B}'_2), \mathcal{O}_{\mathcal{B}'_2}) \ar[r] &
(\Sh(\mathcal{B}'_3), \mathcal{O}_{\mathcal{B}'_3})
}
$$
여기서 (a) 윗행은 $(\Sh(\mathcal{C}), \mathcal{O})$의 1차 두꺼워짐들의
짧은 완전열이고, (b) 아랫행은
$(\Sh(\mathcal{B}), \mathcal{O}_\mathcal{B})$의 1차 두꺼워짐들의 짧은
완전열이며, (c) 각 $f'_i$의 제한은 $f$이고, (d) 각 쌍 $(f, f_i')$는
두꺼워짐의 엄격 사상이며, (e) 각 $f'_i$는 평탄하다. 끝으로
$\mathcal{F}'_2$를 $\mathcal{O}_{\mathcal{B}'_2}$ 위에서 평탄한
$\mathcal{O}'_2$-가군이라 하고
$\mathcal{F} = \mathcal{F}'_2 \otimes \mathcal{O}$라 놓자.
$\pi : (\Sh(\mathcal{C}'_1), \mathcal{O}'_1) \to
(\Sh(\mathcal{C}), \mathcal{O})$를 표준 분할이라 하자
(주 \ref{remark-short-exact-sequence-thickenings-ringed-topoi}).
\end{situation}

\begin{lemma}
\label{lemma-verify-iv-ringed-topoi}
상황 \ref{situation-ses-flat-thickenings-ringed-topoi}에서 가군
$\pi^*\mathcal{F}$와 $h^*\mathcal{F}'_2$는
$\mathcal{O}_{\mathcal{B}'_1}$ 위에서 평탄하고
$(\Sh(\mathcal{C}), \mathcal{O})$ 위로 제한하면 $\mathcal{F}$가 되는
$\mathcal{O}'_1$-가군이다. 이 둘의 차이(보조정리
\ref{lemma-flat-ringed-topoi})는
$\Ext^1_\mathcal{O}(\mathcal{F},
f^*\mathcal{J}_1 \otimes_\mathcal{O} \mathcal{F})$
의 원소 $\theta$이고, 그 경계는
$\Ext^2_\mathcal{O}(\mathcal{F},
f^*\mathcal{J}_3 \otimes_\mathcal{O} \mathcal{F})$
에서 $\mathcal{F}$를 $\mathcal{O}_{\mathcal{B}'_3}$ 위에서 평탄한
$\mathcal{O}'_3$-가군으로 올리는 데 대한 장애 원소(보조정리
\ref{lemma-flat-ringed-topoi})와 같다.
\end{lemma}

\begin{proof}
$\pi^*\mathcal{F}$와 $h^*\mathcal{F}'_2$는 모두
$(\Sh(\mathcal{C}), \mathcal{O})$ 위로 제한하면 $\mathcal{F}$가 되고,
$\pi^*\mathcal{F} \to \mathcal{F}$와
$h^*\mathcal{F}'_2 \to \mathcal{F}$의 핵은
$f^*\mathcal{J}_1 \otimes_\mathcal{O} \mathcal{F}$로 주어진다. 따라서
보조정리 \ref{lemma-deform-module-ringed-topoi}에 의해 두 가군은
평탄하다. 다음 가군의 열이 짧은 완전열이므로 경계를 취할 수 있다.
$$
0 \to f^*\mathcal{J}_3 \otimes_\mathcal{O} \mathcal{F} \to
f^*\mathcal{J}_2 \otimes_\mathcal{O} \mathcal{F} \to
f^*\mathcal{J}_1 \otimes_\mathcal{O} \mathcal{F} \to 0
$$
이는 상황 \ref{situation-ses-flat-thickenings-ringed-topoi}의 가정들과
$\mathcal{F}$가 $\mathcal{O}_\mathcal{B}$ 위에서 평탄하다는 사실에서
따른다. 장애 동치류에 대한 명제는 주
\ref{remark-complex-thickenings-and-ses-modules-ringed-topoi}의 결과를 이
특수한 상황에 그대로 옮긴 것이다.
\end{proof}







\section{환 달린 토포스의 변형과 소박한 코탄젠트 복합체}
\label{section-deformations-ringed-topoi}

\noindent
이 절에서는 소박한 코탄젠트 복합체를 사용하여 변형 이론의 일부를
다룬다. 환 달린 토포스의 1차 두꺼워짐
$t : (\Sh(\mathcal{B}), \mathcal{O}_\mathcal{B}) \to
(\Sh(\mathcal{B}'), \mathcal{O}_{\mathcal{B}'})$에서 시작하자.
$\mathcal{J} = \Ker(t^\sharp)$라 쓰고 $\mathcal{B}$와 $\mathcal{B}'$의
바탕 토포스를 동일시한다. 또한 환 달린 토포스의 사상
$f : (\Sh(\mathcal{C}), \mathcal{O}) \to
(\Sh(\mathcal{B}), \mathcal{O}_\mathcal{B})$, $\mathcal{O}$-가군
$\mathcal{G}$, 그리고 $f^{-1}\mathcal{O}_\mathcal{B}$-가군층의 준동형
$f^{-1}\mathcal{J} \to \mathcal{G}$가 주어졌다고 가정하자. 이 절에서는
다음 도표의 물음표를 채울 수 있는지를 묻는다.
\begin{equation}
\label{equation-to-solve-ringed-topoi}
\vcenter{
\xymatrix{
0 \ar[r] & \mathcal{G} \ar[r] & {?} \ar[r] & \mathcal{O} \ar[r] & 0 \\
0 \ar[r] & f^{-1}\mathcal{J} \ar[u]^c \ar[r] &
f^{-1}\mathcal{O}_{\mathcal{B}'} \ar[u] \ar[r] &
f^{-1}\mathcal{O}_\mathcal{B} \ar[u] \ar[r] & 0
}
}
\end{equation}
또한 해가 존재하면 얼마나 유일한지도 묻는다. 더 정확히는, 1차 두꺼워짐
$i : (\Sh(\mathcal{C}), \mathcal{O}) \to (\Sh(\mathcal{C}'), \mathcal{O}')$
과 (\ref{equation-morphism-thickenings-ringed-topoi})에서와 같은
두꺼워짐의 사상 $(f, f')$를 찾는다. 여기서 $\Ker(i^\sharp)$를
$\mathcal{G}$와 동일시하며 $(f')^\sharp$가 주어진 사상 $c$를 유도해야
한다. 이때 $(\Sh(\mathcal{C}'), \mathcal{O}')$을
(\ref{equation-to-solve-ringed-topoi})의 {\it 해}라고 부른다.


\begin{lemma}
\label{lemma-huge-diagram-ringed-topoi}
다음 환 달린 토포스의 사상들의 가환 도표가 주어졌다고 가정하자.
\begin{equation}
\label{equation-huge-1-ringed-topoi}
\vcenter{
\xymatrix{
& (\Sh(\mathcal{C}_2), \mathcal{O}_2) \ar[r]_{i_2} \ar[d]_{f_2} \ar[ddl]_g &
(\Sh(\mathcal{C}'_2), \mathcal{O}'_2) \ar[d]^{f'_2} \\
&
(\Sh(\mathcal{B}_2), \mathcal{O}_{\mathcal{B}_2}) \ar[r]^{t_2} \ar[ddl]|\hole &
(\Sh(\mathcal{B}'_2), \mathcal{O}_{\mathcal{B}'_2}) \ar[ddl] \\
(\Sh(\mathcal{C}_1), \mathcal{O}_1) \ar[r]_{i_1} \ar[d]_{f_1} &
(\Sh(\mathcal{C}'_1), \mathcal{O}'_1) \ar[d]^{f'_1} \\
(\Sh(\mathcal{B}_1), \mathcal{O}_{\mathcal{B}_1}) \ar[r]^{t_1} &
(\Sh(\mathcal{B}'_1), \mathcal{O}_{\mathcal{B}'_1})
}
}
\end{equation}
가로 화살표들은 1차 두꺼워짐이다. $\mathcal{G}_j = \Ker(i_j^\sharp)$라
놓고, 다음 가환 도표를 유도하는 $g^{-1}\mathcal{O}_1$-가군의 준동형
$\nu : g^{-1}\mathcal{G}_1 \to \mathcal{G}_2$
가 주어졌다고 가정하자.
\begin{equation}
\label{equation-huge-2-ringed-topoi}
\vcenter{
\xymatrix{
& 0 \ar[r] & \mathcal{G}_2 \ar[r] &
\mathcal{O}'_2 \ar[r] &
\mathcal{O}_2 \ar[r] & 0 \\
& 0 \ar[r]|\hole &
f_2^{-1}\mathcal{J}_2 \ar[u]_{c_2} \ar[r] &
f_2^{-1}\mathcal{O}_{\mathcal{B}'_2} \ar[u] \ar[r]|\hole &
f_2^{-1}\mathcal{O}_{\mathcal{B}_2} \ar[u] \ar[r] & 0 \\
0 \ar[r] &
\mathcal{G}_1 \ar[ruu] \ar[r] &
\mathcal{O}'_1 \ar[r] &
\mathcal{O}_1 \ar[ruu] \ar[r] & 0 \\
0 \ar[r] &
f_1^{-1}\mathcal{J}_1 \ar[ruu]|\hole \ar[u]^{c_1} \ar[r] &
f_1^{-1}\mathcal{O}_{\mathcal{B}'_1} \ar[ruu]|\hole \ar[u] \ar[r] &
f_1^{-1}\mathcal{O}_{\mathcal{B}_1} \ar[ruu]|\hole \ar[u] \ar[r] & 0
}
}
\end{equation}
이 도표의 앞면과 뒷면은 (\ref{equation-to-solve-ringed-topoi})의 해이다.
(북북서 방향의 화살표들은 정의역에 $g^{-1}$을 적용한 뒤
$\mathcal{C}_2$ 위에서 얻는 사상들이다.)
\begin{enumerate}
\item 다음 군에는 표준 원소가 존재한다.
$\Ext^1_{\mathcal{O}_2}(
Lg^*\NL_{\mathcal{O}_1/\mathcal{O}_{\mathcal{B}_1}}, \mathcal{G}_2)$
이 원소가 영인 것은 $\nu$와 양립하며
(\ref{equation-huge-1-ringed-topoi})에 들어맞는 환 달린 토포스의 사상
$(\Sh(\mathcal{C}'_2), \mathcal{O}'_2) \to
(\Sh(\mathcal{C}'_1), \mathcal{O}'_1)$이 존재하기 위한 필요충분조건이다.
\item $\nu$와 양립하며 (\ref{equation-huge-1-ringed-topoi})에 들어맞는 사상
$(\Sh(\mathcal{C}'_2), \mathcal{O}'_2) \to
(\Sh(\mathcal{C}'_1), \mathcal{O}'_1)$
이 존재하면, 그러한 모든 사상의 집합은 다음 군 아래의 주동차 공간이다.
$$
\Hom_{\mathcal{O}_1}(
\Omega_{\mathcal{O}_1/\mathcal{O}_{\mathcal{B}_1}}, g_*\mathcal{G}_2) =
\Hom_{\mathcal{O}_2}(
g^*\Omega_{\mathcal{O}_1/\mathcal{O}_{\mathcal{B}_1}}, \mathcal{G}_2) =
\Ext^0_{\mathcal{O}_2}(
Lg^*\NL_{\mathcal{O}_1/\mathcal{O}_{\mathcal{B}_1}}, \mathcal{G}_2).
$$
\end{enumerate}
\end{lemma}

\begin{proof}
이 보조정리의 증명은 보조정리
\ref{lemma-huge-diagram-ringed-spaces}의 증명과 같다. 독자에게 이 증명
대신 그 증명을 읽기를 권한다. 등장하는 각 두꺼워짐의 바탕 토포스를
동일시하겠다(이미 명제에서 이 관례를 사용했다). 보조정리의 마지막
명제에 나오는 등식들은 정의에서 바로 따른다. 따라서 증명의 나머지에서는
다음 군들을 사용한다.
$\Ext^k_{\mathcal{O}_2}(
Lg^*\NL_{\mathcal{O}_1/\mathcal{O}_{\mathcal{B}_1}}, \mathcal{G}_2)$,
$k = 0, 1$. 먼저 보조정리에 나오는 모든 환 달린 토포스의 바탕
토포스가 같은 경우로 환원할 수 있음을 보이자.

\medskip\noindent
이를 위해 $g^{-1}\NL_{\mathcal{O}_1/\mathcal{O}_{\mathcal{B}_1}}$이
환층의 준동형
$g^{-1}f_1^{-1}\mathcal{O}_{\mathcal{B}_1} \to g^{-1}\mathcal{O}_1$의
소박한 코탄젠트 복합체와 같음에 유의하자. 사이트 위의 가군, 보조정리
\ref{sites-modules-lemma-pullback-differentials}를 보라. 더욱이
$\NL_{\mathcal{O}_1/\mathcal{O}_{\mathcal{B}_1}}$의 차수 $0$인 항은
평탄한 $\mathcal{O}_1$-가군이므로 표준 준동형
$$
Lg^*\NL_{\mathcal{O}_1/\mathcal{O}_{\mathcal{B}_1}}
\longrightarrow
g^{-1}\NL_{\mathcal{O}_1/\mathcal{O}_{\mathcal{B}_1}}
\otimes_{g^{-1}\mathcal{O}_1} \mathcal{O}_2
$$
은 차수 $0$과 $-1$의 코호몰로지층에서 동형사상을 유도한다. 따라서
보조정리의 Ext 군들을 다음 군들로 바꾸어도 된다.
$$
\Ext^k_{g^{-1}\mathcal{O}_1}(
g^{-1}\NL_{\mathcal{O}_1/\mathcal{O}_{\mathcal{B}_1}}, \mathcal{G}_2) =
\Ext^k_{g^{-1}\mathcal{O}_1}(
\NL_{g^{-1}\mathcal{O}_1/g^{-1}f_1^{-1}\mathcal{O}_{\mathcal{B}_1}},
\mathcal{G}_2)
$$
(\ref{equation-huge-1-ringed-topoi})에 들어맞고 $\nu$와 양립하는
환 달린 토포스의 사상
$(\Sh(\mathcal{C}'_2), \mathcal{O}'_2) \to
(\Sh(\mathcal{C}'_1), \mathcal{O}'_1)$들의 집합은 $f^\sharp$ 및 $\nu$와
양립하는
$g^{-1}f_1^{-1}\mathcal{O}_{\mathcal{B}'_1}$-대수 준동형
$g^{-1}\mathcal{O}'_1 \to \mathcal{O}'_2$들의 집합과 일대일 대응한다.
따라서 사이트 $\mathcal{C}$ 위의 층들로 이루어진 도표
(\ref{equation-huge-2-ringed-topoi})가 주어졌고(바탕 토포스에서
$f_1 = f_2 = \text{id}$), 그 도표에 들어맞는 환층의 준동형
$\mathcal{O}'_1 \to \mathcal{O}'_2$를 찾는 경우를 가정해도 된다.

\medskip\noindent
보조정리의 증명의 나머지에서는 모든 바탕 위상공간이 같다고 가정한다.
즉 사이트 $\mathcal{C}$ 위의 층들로 이루어진 도표
(\ref{equation-huge-2-ringed-topoi})가 주어졌고(바탕 토포스에서
$f_1 = f_2 = \text{id}$), 그 도표에 들어맞는 환층의 준동형
$\mathcal{O}'_1 \to \mathcal{O}'_2$를 찾는다. Ext 군으로는
$\Ext^k_{\mathcal{O}_1}(
\NL_{\mathcal{O}_1/\mathcal{O}_{\mathcal{B}_1}}, \mathcal{G}_2)$, $k = 0, 1$.

\medskip\noindent
1단계. 장애 동치류의 구성. 다음 집합층을 생각하자.
$$
\mathcal{E} = \mathcal{O}'_1 \times_{\mathcal{O}_2} \mathcal{O}'_2
$$
여기에는 전사 사상 $\alpha : \mathcal{E} \to \mathcal{O}_1$이 따라오므로
$\NL_{\mathcal{O}_1/\mathcal{O}_{\mathcal{B}_1}}$ 대신
$\NL(\alpha)$를 사용할 수 있다. 사이트 위의 가군, 보조정리
\ref{sites-modules-lemma-NL-up-to-qis}를 보라. 다음과 같이 놓자.
$$
\mathcal{I}' =
\Ker(\mathcal{O}_{\mathcal{B}'_1}[\mathcal{E}] \to \mathcal{O}_1)
\quad\text{그리고}\quad
\mathcal{I} =
\Ker(\mathcal{O}_{\mathcal{B}_1}[\mathcal{E}] \to \mathcal{O}_1)
$$
핵이 $\mathcal{J}_1\mathcal{O}_{\mathcal{B}'_1}[\mathcal{E}]$인 전사
준동형 $\mathcal{I}' \to \mathcal{I}$이 있다. 두 개의
$\mathcal{O}_{\mathcal{B}'_2}$-대수 준동형
$$
a : \mathcal{O}_{\mathcal{B}'_1}[\mathcal{E}] \to \mathcal{O}'_1
\quad\text{그리고}\quad
b : \mathcal{O}_{\mathcal{B}'_1}[\mathcal{E}] \to \mathcal{O}'_2
$$
를 얻는다. 이들은 각각 준동형
$a|_{\mathcal{I}'} : \mathcal{I}' \to \mathcal{G}_1$과
$b|_{\mathcal{I}'} : \mathcal{I}' \to \mathcal{G}_2$를 유도한다.
$a$와 $b$는 모두 $(\mathcal{I}')^2$을 영으로 보낸다. 더욱이
(\ref{equation-huge-2-ringed-topoi})의 왼쪽 정사각형이 가환하므로
$\mathcal{G}_2$로 가는 사상으로서 $a$와 $b$는
$\mathcal{J}_1\mathcal{O}_{\mathcal{B}'_1}[\mathcal{E}]$ 위에서 일치한다.
따라서 차
$b|_{\mathcal{I}'} - \nu \circ a|_{\mathcal{I}'}$
는 잘 정의된 $\mathcal{O}_1$-선형 준동형
$$
\xi : \mathcal{I}/\mathcal{I}^2 \longrightarrow \mathcal{G}_2
$$
를 유도한다. 이 준동형은 $\mathcal{I}$의 국소절단 $f$의 동치류를
$a(f') - \nu(b(f'))$로 보내는데, 여기서 $f'$는 $f$를
$\mathcal{I}'$의 국소절단으로 올린 것이다. 그 상을
$[\xi] \in \Ext^1_{\mathcal{O}_1}(\NL(\alpha), \mathcal{G}_2)$
라 하자(아래를 보라).

\medskip\noindent
2단계. $[\xi]$의 소멸은 필요하다. 다음과 같이 쓰자. $\Omega =
\Omega_{\mathcal{O}_{\mathcal{B}_1}[\mathcal{E}]/\mathcal{O}_{\mathcal{B}_1}}
\otimes_{\mathcal{O}_{\mathcal{B}_1}[\mathcal{E}]} \mathcal{O}_1$.
$\NL(\alpha) = (\mathcal{I}/\mathcal{I}^2 \to \Omega)$가 다음 완전
삼각형에 들어맞음에 유의하자.
$$
\Omega[0] \to
\NL(\alpha) \to
\mathcal{I}/\mathcal{I}^2[1] \to
\Omega[1]
$$
따라서 어떤 준동형 $\Omega \to \mathcal{G}_2$에 대하여 $\xi$가 합성
$\mathcal{I}/\mathcal{I}^2 \to \Omega \to \mathcal{G}_2$일 때, 그리고
그때에만 $[\xi]$가 영이다. (\ref{equation-huge-2-ringed-topoi})에
들어맞는 환층의 준동형
$\varphi : \mathcal{O}'_1 \to \mathcal{O}'_2$가 존재한다고 하자. 이 경우 준동형
$\mathcal{O}'_1[\mathcal{E}] \to \mathcal{G}_2$,
$f' \mapsto b(f') - \varphi(a(f'))$를 생각하자. 계산하면 이 준동형은
$\mathcal{J}_1\mathcal{O}_{\mathcal{B}'_1}[\mathcal{E}]$을 영으로 보내며
미분 $\mathcal{O}_{\mathcal{B}_1}[\mathcal{E}] \to \mathcal{G}_2$를
유도함을 알 수 있다. 그 결과 얻는 선형 준동형
$\Omega \to \mathcal{G}_2$는 이 경우 $[\xi] = 0$임을 보인다.

\medskip\noindent
3단계. $[\xi]$의 소멸은 충분하다.
$\theta : \Omega \to \mathcal{G}_2$를 $\mathcal{O}_1$-선형 준동형이라
하자. 이때 $\xi$는
$\theta \circ (\mathcal{I}/\mathcal{I}^2 \to \Omega)$와 같다고
가정한다. 계산하면
$$
b + \theta \circ d :
\mathcal{O}_{\mathcal{B}'_1}[\mathcal{E}]
\longrightarrow
\mathcal{O}'_2
$$
은 $\mathcal{I}'$을 영으로 보내므로, (\ref{equation-huge-2-ringed-topoi})에
들어맞는 준동형 $\mathcal{O}'_1 \to \mathcal{O}'_2$를 정한다.

\medskip\noindent
위의 특수한 경우에서 (2)의 증명. 생략한다. 도움말: 이는 보조정리
\ref{lemma-huge-diagram}의 (2)의 증명과 완전히 같다.
\end{proof}

\begin{lemma}
\label{lemma-NL-represent-ext-class-ringed-topoi}
$\mathcal{C}$를 사이트라 하자. $\mathcal{A} \to \mathcal{B}$를
$\mathcal{C}$ 위 환층의 준동형이라 하고 $\mathcal{G}$를
$\mathcal{B}$-가군이라 하자.
$\xi \in \Ext^1_\mathcal{B}(\NL_{\mathcal{B}/\mathcal{A}}, \mathcal{G})$라
하자. 집합층의 사상 $\alpha : \mathcal{E} \to \mathcal{B}$를 적절히
택하면, $\xi \in \Ext^1_\mathcal{B}(\NL(\alpha), \mathcal{G})$는 준동형
$\mathcal{I}/\mathcal{I}^2 \to \mathcal{G}$의 동치류로 표현된다
(표기는 증명을 보라).
\end{lemma}

\begin{proof}
사상 $\alpha : \mathcal{E} \to \mathcal{B}$가 주어지고
$\mathcal{A}[\mathcal{E}] \to \mathcal{B}$가 핵이 $\mathcal{I}$인
전사 준동형이면, 복합체
$\NL(\alpha) = (\mathcal{I}/\mathcal{I}^2 \to 
\Omega_{\mathcal{A}[\mathcal{E}]/\mathcal{A}}
\otimes_{\mathcal{A}[\mathcal{E}]} \mathcal{B})$는
$\NL_{\mathcal{B}/\mathcal{A}}$와 표준적으로 동형임을 상기하자.
사이트 위의 가군, 보조정리 \ref{sites-modules-lemma-NL-up-to-qis}를
보라. 더욱이
$\Omega = \Omega_{\mathcal{A}[\mathcal{E}]/\mathcal{A}}
\otimes_{\mathcal{A}[\mathcal{E}]} \mathcal{B}$는 다음 준층에 연관된
층이다.
$U \mapsto \bigoplus_{e \in \mathcal{E}(U)} \mathcal{B}(U)$.
다시 말해 $\Omega$는 집합층 $\mathcal{E}$ 위의 자유
$\mathcal{B}$-가군이고, 특히 표준 사상 $\mathcal{E} \to \Omega$가 있다.

\medskip\noindent
이제 어떤 $\mathcal{E}$를 택하자(예를 들어 소박한 코탄젠트 복합체의
정의에서처럼 $\mathcal{E} = \mathcal{B}$로 택한다). $\xi$를 준동형
$\mathcal{I}/\mathcal{I}^2 \to \mathcal{G}$의 동치류로 표현하는 데 대한
장애 원소는 $\Ext^1_\mathcal{B}(\Omega, \mathcal{G})$에 속한다.
이 원소가 $\mathcal{B}$-가군의 확대
$0 \to \mathcal{G} \to \mathcal{H} \to \Omega \to 0$로 표현된다고
하자. 유도된 사상 $\alpha' : \mathcal{E}' \to \mathcal{B}$를 갖는
집합층
$\mathcal{E}' = \mathcal{E} \times_\Omega \mathcal{H}$
를 생각하자. $\mathcal{I}' =
\Ker(\mathcal{A}[\mathcal{E}'] \to \mathcal{B})$ 및
$\Omega' = \Omega_{\mathcal{A}[\mathcal{E}']/\mathcal{A}}
\otimes_{\mathcal{A}[\mathcal{E}']} \mathcal{B}$라 놓자.
준동형 $\Omega' \to \Omega$에 따라 확대 $\mathcal{H}$를 당기면
분할된다. 실제로 $\Omega'$은 집합층 $\mathcal{E}'$ 위의 자유
$\mathcal{B}$-가군이고, 구성상 다음 가환 도표가 있기 때문이다.
$$
\xymatrix{
\mathcal{E}' \ar[r] \ar[d] & \mathcal{E} \ar[d] \\
\mathcal{H} \ar[r] & \Omega
}
$$
따라서 준동형 $\NL(\alpha') \to \NL(\alpha)$에 따른 $\xi$의 당김은
$\Ext^1_\mathcal{B}(\Omega', \mathcal{G})$에서 영으로 간다. 이로써
증명이 끝난다.
\end{proof}

\begin{lemma}
\label{lemma-choices-ringed-topoi}
(\ref{equation-to-solve-ringed-topoi})의 해가 존재하면, 해들의 동형류
집합은\par\noindent
$\Ext^1_\mathcal{O}(
\NL_{\mathcal{O}/\mathcal{O}_\mathcal{B}}, \mathcal{G})$ 아래의 주동차
공간이다.
\end{lemma}

\begin{proof}
(\ref{equation-to-solve-ringed-topoi})의 두 해 $\mathcal{O}'_1$과
$\mathcal{O}'_2$가 주어지면, 보조정리
\ref{lemma-huge-diagram-ringed-topoi}에 의해 준동형
$\mathcal{O}'_1 \to \mathcal{O}'_2$의 존재에 대한 장애 원소
$o(\mathcal{O}'_1, \mathcal{O}'_2) \in \Ext^1_\mathcal{O}(
\NL_{\mathcal{O}/\mathcal{O}_\mathcal{B}}, \mathcal{G})$
를 곧바로 얻는다. 이 원소는 동형사상의 존재에 대한 장애 원소이므로
동형류들을 구별한다. 따라서 증명을 끝내려면 해 $\mathcal{O}'$과 원소
$\xi \in \Ext^1_\mathcal{O}(
\NL_{\mathcal{O}/\mathcal{O}_\mathcal{B}}, \mathcal{G})$
가 주어졌을 때 $o(\mathcal{O}', \mathcal{O}'_\xi) = \xi$를 만족하는
두 번째 해 $\mathcal{O}'_\xi$를 찾을 수 있음을 보이면 충분하다.

\medskip\noindent
동치류 $\xi$에 대하여 보조정리
\ref{lemma-NL-represent-ext-class-ringed-topoi}에서와 같이
$\alpha : \mathcal{E} \to \mathcal{O}$를 택하자. 핵이
$\mathcal{I}$인 전사 준동형
$f^{-1}\mathcal{O}_\mathcal{B}[\mathcal{E}] \to \mathcal{O}$과 그에
대응하는 소박한 코탄젠트 복합체
$\NL(\alpha) = (\mathcal{I}/\mathcal{I}^2 \to
\Omega_{f^{-1}\mathcal{O}_\mathcal{B}[\mathcal{E}]/
f^{-1}\mathcal{O}_\mathcal{B}}
\otimes_{f^{-1}\mathcal{O}_\mathcal{B}[\mathcal{E}]} \mathcal{O})$.
이 보조정리에 의해 $\xi$는 준동형
$\delta : \mathcal{I}/\mathcal{I}^2 \to \mathcal{G}$.
의 동치류이다. $\mathcal{E}$를
$\mathcal{E} \times_\mathcal{O} \mathcal{O}'$로 바꾸면, $\alpha$가
사상 $\alpha' : \mathcal{E} \to \mathcal{O}'$를 거쳐 분해된다고
가정할 수도 있다.

\medskip\noindent
이 선택들은 $f^{-1}\mathcal{O}_{\mathcal{B}'}$-대수 준동형
$\varphi : \mathcal{O}_{\mathcal{B}'}[\mathcal{E}] \to \mathcal{O}'$.
를 정한다. $\mathcal{I}' = \Ker(\varphi)$라 놓자. $\varphi$는 준동형
$\varphi|_{\mathcal{I}'} : \mathcal{I}' \to \mathcal{G}$
를 유도하고 $\mathcal{O}'$은 다음 도표에서와 같은 밀어내기이다.
$$
\xymatrix{
0 \ar[r] & \mathcal{G} \ar[r] & \mathcal{O}' \ar[r] &
\mathcal{O} \ar[r] & 0 \\
0 \ar[r] & \mathcal{I}' \ar[u]^{\varphi|_{\mathcal{I}'}} \ar[r] &
f^{-1}\mathcal{O}_{\mathcal{B}'}[\mathcal{E}] \ar[u] \ar[r] &
\mathcal{O} \ar[u]_{=} \ar[r] & 0
}
$$
$\psi : \mathcal{I}' \to \mathcal{G}$를
$\varphi|_{\mathcal{I}'}$와 다음 합성의 합이라 하자.
$$
\mathcal{I}' \to \mathcal{I}'/(\mathcal{I}')^2 \to
\mathcal{I}/\mathcal{I}^2 \xrightarrow{\delta} \mathcal{G}.
$$
그러면 $\psi$에 따른 밀어내기는 위와 같은 도표에 들어맞는 또 다른
환의 확대 $\mathcal{O}'_\xi$이다. 계산하면(생략한다)
$o(\mathcal{O}', \mathcal{O}'_\xi) = \xi$임을 알 수 있다.
\end{proof}

\begin{lemma}
\label{lemma-extensions-of-relative-ringed-topoi}
$f : (\Sh(\mathcal{C}), \mathcal{O}) \to
(\Sh(\mathcal{B}), \mathcal{O}_\mathcal{B})$를 환 달린 토포스의
사상이라 하고 $\mathcal{G}$를 $\mathcal{O}$-가군이라 하자. 다음
$f^{-1}\mathcal{O}_\mathcal{B}$-대수의 확대들의 동형류 집합을 생각하자.
$$
0 \to \mathcal{G} \to \mathcal{O}' \to \mathcal{O} \to 0
$$
여기서 $\mathcal{G}$는 제곱이 영인 아이디얼이다.\footnote{다시 말해,
$\mathcal{O}$-가군의 동형사상
$\mathcal{G} \to \Ker(i^\sharp)$가 주어진
$(\Sh(\mathcal{B}), \mathcal{O}_\mathcal{B})$ 위의 1차 두꺼워짐
$i : (\Sh(\mathcal{C}), \mathcal{O}) \to (\Sh(\mathcal{C}), \mathcal{O}')$
들의 동형류 집합이다.}
이 집합은 다음 군과 표준적으로 전단사이다.
$\Ext^1_\mathcal{O}(\NL_{\mathcal{O}/\mathcal{O}_\mathcal{B}}, \mathcal{G})$.
\end{lemma}

\begin{proof}
이를 증명하기 위해, (\ref{equation-to-solve-ringed-topoi})이 다음
도표로 주어지는 경우에 앞의 결과들을 적용하자.
$$
\xymatrix{
0 \ar[r] &
\mathcal{G} \ar[r] &
{?} \ar[r] &
\mathcal{O} \ar[r] & 0 \\
0 \ar[r] &
0 \ar[u] \ar[r] &
f^{-1}\mathcal{O}_\mathcal{B} \ar[u] \ar[r]^{\text{id}} &
f^{-1}\mathcal{O}_\mathcal{B} \ar[u] \ar[r] & 0
}
$$
해 $\mathcal{G} \oplus \mathcal{O}$가 존재하므로, 이 보조정리는
보조정리 \ref{lemma-choices-ringed-topoi}에서 따른다. 전단사의 직접
구성은 아래의 주를 보라.
\end{proof}

\begin{remark}
\label{remark-extensions-of-relative-ringed-topoi}
$f : (\Sh(\mathcal{C}), \mathcal{O}) \to
(\mathcal{B}, \mathcal{O}_\mathcal{B})$와 $\mathcal{G}$를 보조정리
\ref{lemma-extensions-of-relative-ringed-topoi}에서와 같이 택하자.
보조정리에서와 같은 확대
$0 \to \mathcal{G} \to \mathcal{O}' \to \mathcal{O} \to 0$
를 생각하자. 집합층 $\mathcal{E}$와 다음 가환 도표를 택할 수 있다.
$$
\xymatrix{
\mathcal{E} \ar[d]_{\alpha'} \ar[rd]^\alpha \\
\mathcal{O}' \ar[r] & \mathcal{O}
}
$$
여기서 $f^{-1}\mathcal{O}_\mathcal{B}[\mathcal{E}] \to \mathcal{O}$은
핵이 $\mathcal{J}$인 전사 준동형이다. 예를 들어 $\mathcal{O}'$로
전사하는 임의의 집합층을 택할 수 있다. 그러면
$$
\NL_{\mathcal{O}/\mathcal{O}_\mathcal{B}} \cong \NL(\alpha) = 
\left(
\mathcal{J}/\mathcal{J}^2
\longrightarrow
\Omega_{f^{-1}\mathcal{O}_\mathcal{B}[\mathcal{E}]/
f^{-1}\mathcal{O}_\mathcal{B}}
\otimes_{f^{-1}\mathcal{O}_\mathcal{B}[\mathcal{E}]} \mathcal{O}\right)
$$
이다. 사이트 위의 가군, 절 \ref{sites-modules-section-netherlander},
특히 보조정리 \ref{sites-modules-lemma-NL-up-to-qis}를 보라. 물론
$\alpha'$는 준동형
$f^{-1}\mathcal{O}_\mathcal{B}[\mathcal{E}] \to \mathcal{O}'$
를 정하고, 이는 다시 준동형
$$
\mathcal{J}/\mathcal{J}^2 \longrightarrow \mathcal{G}
$$
를 정하며, 이는 다시 다음 군의 원소를 정한다.\par\noindent
$\Ext^1_\mathcal{O}(\NL(\alpha), \mathcal{G}) =
\Ext^1_\mathcal{O}(\NL_{\mathcal{O}/\mathcal{O}_\mathcal{B}}, \mathcal{G})$
보조정리의 전단사에서 이 원소는 $\mathcal{O}'$에 대응한다.
\end{remark}

\begin{lemma}
\label{lemma-extensions-of-relative-ringed-topoi-functorial}
환 달린 토포스의 사상
$f : (\Sh(\mathcal{C}), \mathcal{O}_\mathcal{C}) \to
(\Sh(\mathcal{B}), \mathcal{O}_\mathcal{B})$와
$g : (\Sh(\mathcal{D}), \mathcal{O}_\mathcal{D}) \to
(\Sh(\mathcal{C}), \mathcal{O}_\mathcal{C})$를 생각하자. $\mathcal{F}$를
$\mathcal{O}_\mathcal{C}$-가군, $\mathcal{G}$를
$\mathcal{O}_\mathcal{D}$-가군이라 하고
$c : g^*\mathcal{F} \to \mathcal{G}$를 $\mathcal{O}_\mathcal{D}$-선형
준동형이라 하자. 끝으로 다음을 생각하자.
\begin{enumerate}
\item[(a)]
$0 \to \mathcal{F} \to \mathcal{O}_{\mathcal{C}'} \to
\mathcal{O}_\mathcal{C} \to 0$
은
$\xi \in \Ext^1_{\mathcal{O}_\mathcal{C}}(
\NL_{\mathcal{O}_\mathcal{C}/\mathcal{O}_\mathcal{B}}, \mathcal{F})$에
대응하는 $f^{-1}\mathcal{O}_\mathcal{B}$-대수의 확대이다.
\item[(b)]
$0 \to \mathcal{G} \to \mathcal{O}_{\mathcal{D}'} \to
\mathcal{O}_\mathcal{D} \to 0$
은
$\zeta \in \Ext^1_{\mathcal{O}_\mathcal{D}}(
\NL_{\mathcal{O}_\mathcal{D}/\mathcal{O}_\mathcal{B}}, \mathcal{G})$에
대응하는 $g^{-1}f^{-1}\mathcal{O}_\mathcal{B}$-대수의 확대이다.
\end{enumerate}
보조정리 \ref{lemma-extensions-of-relative-ringed-topoi}를 보라. 다음
환 달린 토포스의 사상
$$
g' :
(\Sh(\mathcal{D}), \mathcal{O}_{\mathcal{D}'})
\longrightarrow
(\Sh(\mathcal{C}), \mathcal{O}_{\mathcal{C}'})
$$
이 $(\Sh(\mathcal{B}), \mathcal{O}_\mathcal{B})$ 위에서 존재하고
$g$ 및 $c$와 양립할 필요충분조건은 $\xi$와 $\zeta$의 상이 다음
군에서 같은 원소인 것이다.
$\Ext^1_{\mathcal{O}_\mathcal{D}}(
Lg^*\NL_{\mathcal{O}_\mathcal{C}/\mathcal{O}_\mathcal{B}}, \mathcal{G})$.
\end{lemma}

\begin{proof}
다음 준동형들이 있으므로 명제는 의미가 있다.
$$
\Ext^1_{\mathcal{O}_\mathcal{C}}(
\NL_{\mathcal{O}_\mathcal{C}/\mathcal{O}_\mathcal{B}}, \mathcal{F}) \to
\Ext^1_{\mathcal{O}_\mathcal{D}}(
Lg^*\NL_{\mathcal{O}_\mathcal{C}/\mathcal{O}_\mathcal{B}}, Lg^*\mathcal{F}) \to
\Ext^1_{\mathcal{O}_\mathcal{D}}
(Lg^*\NL_{\mathcal{O}_\mathcal{C}/\mathcal{O}_\mathcal{B}}, \mathcal{G})
$$
첫 준동형들은 사상
$Lg^*\mathcal{F} \to g^*\mathcal{F} \xrightarrow{c} \mathcal{G}$를
사용하고, 다음 준동형은
$$
\Ext^1_{\mathcal{O}_Y}(
\NL_{\mathcal{O}_\mathcal{D}/\mathcal{O}_\mathcal{B}}, \mathcal{G}) \to
\Ext^1_{\mathcal{O}_Y}(
Lg^*\NL_{\mathcal{O}_\mathcal{C}/\mathcal{O}_\mathcal{B}}, \mathcal{G})
$$
사상 $Lg^*\NL_{\mathcal{O}_\mathcal{C}/\mathcal{O}_\mathcal{B}} \to
\NL_{\mathcal{O}_\mathcal{D}/\mathcal{O}_\mathcal{B}}$를 사용한다.
보조정리 \ref{lemma-huge-diagram-ringed-topoi}을 다음 도표에 적용하면
보조정리의 명제를 얻을 수 있다.
$$
\xymatrix{
& 0 \ar[r] &
\mathcal{G} \ar[r] &
\mathcal{O}_{\mathcal{D}'} \ar[r] &
\mathcal{O}_\mathcal{D} \ar[r] & 0 \\
& 0 \ar[r]|\hole & 0 \ar[u] \ar[r] &
g^{-1}f^{-1}\mathcal{O}_\mathcal{B} \ar[u] \ar[r]|\hole &
g^{-1}f^{-1}\mathcal{O}_\mathcal{B} \ar[u] \ar[r] & 0 \\
0 \ar[r] &
\mathcal{F} \ar[ruu] \ar[r] &
\mathcal{O}_{\mathcal{C}'} \ar[r] &
\mathcal{O}_\mathcal{C} \ar[ruu] \ar[r] & 0 \\
0 \ar[r] & 0 \ar[ruu]|\hole \ar[u] \ar[r] &
f^{-1}\mathcal{O}_\mathcal{B} \ar[ruu]|\hole \ar[u] \ar[r] &
f^{-1}\mathcal{O}_\mathcal{B} \ar[ruu]|\hole \ar[u] \ar[r] & 0
}
$$
여기에 보조정리 \ref{lemma-extensions-of-relative-ringed-topoi}와
\ref{lemma-huge-diagram-ringed-topoi}의 증명에 나오는 구성들 사이의
양립성을 함께 사용한다. 그 양립성의 명제와 증명은 생략한다. 직접
논증은 아래의 주를 보라.
\end{proof}

\begin{remark}
\label{remark-extensions-of-relative-ringed-topoi-functorial}
$f : (\Sh(\mathcal{C}), \mathcal{O}_\mathcal{C}) \to
(\Sh(\mathcal{B}), \mathcal{O}_\mathcal{B})$,
$g : (\Sh(\mathcal{D}), \mathcal{O}_\mathcal{D}) \to
(\Sh(\mathcal{C}), \mathcal{O}_\mathcal{C})$,
$\mathcal{F}$,
$\mathcal{G}$,
$c : g^*\mathcal{F} \to \mathcal{G}$,
$0 \to \mathcal{F} \to \mathcal{O}_{\mathcal{C}'} \to
\mathcal{O}_\mathcal{C} \to 0$,
$\xi \in \Ext^1_{\mathcal{O}_\mathcal{C}}(
\NL_{\mathcal{O}_\mathcal{C}/\mathcal{O}_\mathcal{B}}, \mathcal{F})$,
$0 \to \mathcal{G} \to \mathcal{O}_{\mathcal{D}'} \to
\mathcal{O}_\mathcal{D} \to 0$, 그리고
$\zeta \in \Ext^1_{\mathcal{O}_\mathcal{D}}(
\NL_{\mathcal{O}_\mathcal{D}/\mathcal{O}_\mathcal{B}}, \mathcal{G})$
를 보조정리 \ref{lemma-extensions-of-relative-ringed-topoi-functorial}에서와
같이 택하자. $c : g^{-1}\mathcal{F} \to \mathcal{G}$에 따른 밀어내기를
사용하여 다음 확대를 구성할 수 있다.
$$
\xymatrix{
0 \ar[r] &
\mathcal{G} \ar[r] &
\mathcal{O}'_1 \ar[r] &
g^{-1}\mathcal{O}_\mathcal{C} \ar[r] & 0 \\
0 \ar[r] &
g^{-1}\mathcal{F} \ar[u]^c \ar[r] &
g^{-1}\mathcal{O}_{\mathcal{C}'} \ar[u] \ar[r] &
g^{-1}\mathcal{O}_\mathcal{C} \ar@{=}[u] \ar[r] & 0
}
$$
$g^\sharp : g^{-1}\mathcal{O}_\mathcal{C} \to \mathcal{O}_\mathcal{D}$에
따른 당김을 사용하여 다음 확대를 구성할 수 있다.
$$
\xymatrix{
0 \ar[r] &
\mathcal{G} \ar[r] &
\mathcal{O}_{\mathcal{D}'} \ar[r] &
\mathcal{O}_\mathcal{D} \ar[r] & 0 \\
0 \ar[r] &
\mathcal{G} \ar@{=}[u] \ar[r] &
\mathcal{O}'_2 \ar[u] \ar[r] &
g^{-1}\mathcal{O}_\mathcal{C} \ar[u] \ar[r] & 0
}
$$
도표 추적에 따르면, 환 달린 토포스의 사상
$g' : (\Sh(\mathcal{D}), \mathcal{O}_{\mathcal{D}'}) \to
(\Sh(\mathcal{C}), \mathcal{O}_{\mathcal{C}'})$
이 $(\Sh(\mathcal{B}), \mathcal{O}_\mathcal{B})$ 위에서 존재하고
$g$ 및 $c$와 양립할 필요충분조건은
$g^{-1}\mathcal{O}_\mathcal{C}$의 $\mathcal{G}$에 의한
$g^{-1}f^{-1}\mathcal{O}_\mathcal{B}$-대수 확대들로서
$\mathcal{O}'_1$과 $\mathcal{O}'_2$가 동형인 것이다. 보조정리
\ref{lemma-extensions-of-relative-ringed-topoi}에 의해 이 확대들은
다음 등식의 왼쪽 항으로 분류된다.
$$
\Ext^1_{g^{-1}\mathcal{O}_\mathcal{C}}(
\NL_{g^{-1}\mathcal{O}_\mathcal{C}/g^{-1}f^{-1}\mathcal{O}_\mathcal{B}},
\mathcal{G}) =
\Ext^1_{\mathcal{O}_\mathcal{D}}(
Lg^*\NL_{\mathcal{O}_\mathcal{C}/\mathcal{O}_\mathcal{B}}, \mathcal{G})
$$
여기서 등식은 텐서-Hom 수반성과 다음 등식들에서 따른다.
$$
\vcenter{\halign{\hfil\ensuremath{\displaystyle #}\hfil\cr
\NL_{g^{-1}\mathcal{O}_\mathcal{C}/g^{-1}f^{-1}\mathcal{O}_\mathcal{B}} =
g^{-1}\NL_{\mathcal{O}_\mathcal{C}/\mathcal{O}_\mathcal{B}}\cr
\text{그리고}\cr
Lg^*\NL_{\mathcal{O}_\mathcal{C}/\mathcal{O}_\mathcal{B}} =
g^{-1}\NL_{\mathcal{O}_\mathcal{C}/\mathcal{O}_\mathcal{B}}
\otimes_{g^{-1}\mathcal{O}_X}^\mathbf{L} \mathcal{O}_Y\cr}}
$$
첫째 등식은 사이트 위의 가군, 보조정리
\ref{sites-modules-lemma-pullback-NL}를 보라. 둘째 등식은 유도 당김의
정의에서 따른다. 따라서 보조정리
\ref{lemma-extensions-of-relative-ringed-topoi-functorial}을 보이려면
$\mathcal{O}'_1$이 $\xi$의 상에 대응하고 $\mathcal{O}'_2$가
$\zeta$의 상에 대응함을 보이면 충분하다. $\xi$와
$\mathcal{O}'_1$ 사이의 대응은 주
\ref{remark-extensions-of-relative-ringed-topoi}에서 동치류 $\xi$를
구성한 방법에서 바로 알 수 있다. $\zeta$와 $\mathcal{O}'_2$ 사이의
대응을 위해 먼저 다음 가환 도표를 택한다.
$$
\xymatrix{
\mathcal{E} \ar[d]_{\beta'} \ar[rd]^\beta \\
\mathcal{O}_{\mathcal{D}'} \ar[r] & \mathcal{O}_\mathcal{D}
}
$$
여기서 $g^{-1}f^{-1}\mathcal{O}_\mathcal{B}[\mathcal{E}] \to
\mathcal{O}_\mathcal{D}$는 핵이 $\mathcal{K}$인 전사 준동형이다.
다음으로 가환 도표
$$
\xymatrix{
\mathcal{E} \ar[d]_{\beta'} &
\mathcal{E}' \ar[l]^\varphi \ar[d]_{\alpha'} \ar[rd]^\alpha \\
\mathcal{O}_{\mathcal{D}'} &
\mathcal{O}'_2 \ar[l] \ar[r] &
g^{-1}\mathcal{O}_\mathcal{C}
}
$$
를 택한다. 여기서
$g^{-1}f^{-1}\mathcal{O}_\mathcal{B}[\mathcal{E}'] \to
g^{-1}\mathcal{O}_\mathcal{C}$는 핵이 $\mathcal{J}$인 전사
준동형이다. 예를 들어 집합층으로서
$\mathcal{E}' = \mathcal{E} \amalg \mathcal{O}'_2$로 택하면 된다.
사상 $\varphi$는 복합체의 사상 $\NL(\alpha) \to \NL(\beta)$를
유도하고(표기는 가군, 절 \ref{modules-section-netherlander}에서와
같다), 특히
$\bar\varphi : \mathcal{J}/\mathcal{J}^2 \to \mathcal{K}/\mathcal{K}^2$.
를 유도한다. 이때
$\NL(\alpha) \cong \NL_{\mathcal{O}_\mathcal{D}/\mathcal{O}_\mathcal{B}}$
이고
$\NL(\beta) \cong
\NL_{g^{-1}\mathcal{O}_\mathcal{C}/g^{-1}f^{-1}\mathcal{O}_\mathcal{B}}$
이다. 복합체의 사상 $\NL(\alpha) \to \NL(\beta)$는
사상
$Lg^*\NL_{\mathcal{O}_\mathcal{C}/\mathcal{O}_\mathcal{B}} \to
\NL_{\mathcal{O}_\mathcal{D}/\mathcal{O}_\mathcal{B}}$
를 나타낸다. 이 사상은 보조정리
\ref{lemma-extensions-of-relative-ringed-topoi-functorial}의 명제에서
사용한 것이다(그 증명의 첫 부분을 보라). 이제 $\zeta$는 $\beta'$가
유도하는 준동형 $\mathcal{K}/\mathcal{K}^2 \to \mathcal{G}$의
동치류에 대응한다. 주
\ref{remark-extensions-of-relative-ringed-topoi}를 보라. 마찬가지로
확대 $\mathcal{O}'_2$는 $\alpha'$가 유도하는 준동형
$\mathcal{J}/\mathcal{J}^2 \to \mathcal{G}$에 대응한다. 위의 가환
도표에 의해 이 준동형은 $\beta'$가 유도하는 준동형
$\mathcal{K}/\mathcal{K}^2 \to \mathcal{G}$와 다음 준동형의 합성이다.
$\bar\varphi : \mathcal{J}/\mathcal{J}^2 \to \mathcal{K}/\mathcal{K}^2$.
이것이 필요한 양립성을 증명한다.
\end{remark}

\begin{lemma}
\label{lemma-parametrize-solutions-ringed-topoi}
$t : (\Sh(\mathcal{B}), \mathcal{O}_\mathcal{B}) \to
(\Sh(\mathcal{B}'), \mathcal{O}_{\mathcal{B}'})$,
$\mathcal{J} = \Ker(t^\sharp)$,
$f : (\Sh(\mathcal{C}), \mathcal{O}) \to
(\Sh(\mathcal{B}), \mathcal{O}_\mathcal{B})$, $\mathcal{G}$, 그리고
$c : \mathcal{J} \to \mathcal{G}$를
(\ref{equation-to-solve-ringed-topoi})에서와 같이 택하자. 보조정리
\ref{lemma-extensions-of-relative-ringed-topoi}에 의해
$\mathcal{O}_\mathcal{B}$의 $\mathcal{J}$에 의한 확대
$\mathcal{O}_{\mathcal{B}'}$에 대응하는 원소를
$\xi \in \Ext^1_{\mathcal{O}_\mathcal{B}}(
\NL_{\mathcal{O}_\mathcal{B}/\mathcal{O}_{\mathcal{B}'}}, \mathcal{J})$
라 쓰자. 해들의 동형류 집합은 다음 준동형에서 $\xi$의 상 위의 올과
표준적으로 전단사이다.
$$
\Ext^1_\mathcal{O}(\NL_{\mathcal{O}/\mathcal{O}_{\mathcal{B}'}},
\mathcal{G})\to
\Ext^1_\mathcal{O}(
Lf^*\NL_{\mathcal{O}_\mathcal{B}/\mathcal{O}_{\mathcal{B}'}}, \mathcal{G})
$$
\end{lemma}

\begin{proof}
보조정리 \ref{lemma-extensions-of-relative-ringed-topoi}를
$t \circ f : (\Sh(\mathcal{C}), \mathcal{O}) \to
(\Sh(\mathcal{B}'), \mathcal{O}_{\mathcal{B}'})$
와 $\mathcal{O}$-가군 $\mathcal{G}$에 적용하면, 다음 군의 원소
$\zeta$들이
$\Ext^1_\mathcal{O}(\NL_{\mathcal{O}/\mathcal{O}_{\mathcal{B}'}},
\mathcal{G})$
다음 $f^{-1}\mathcal{O}_{\mathcal{B}'}$-대수의 확대들을 매개화함을
알 수 있다.
$0 \to \mathcal{G} \to \mathcal{O}' \to \mathcal{O} \to 0$
보조정리
\ref{lemma-extensions-of-relative-ringed-topoi-functorial}을 다음에
적용하자.
$$
(\Sh(\mathcal{C}), \mathcal{O}) \xrightarrow{f}
(\Sh(\mathcal{B}), \mathcal{O}_\mathcal{B}) \xrightarrow{t}
(\Sh(\mathcal{B}'), \mathcal{O}_{\mathcal{B}'})
$$
그리고 $c : \mathcal{J} \to \mathcal{G}$를 함께 사용한다. 그러면
다음 사상
$$
f' : 
(\Sh(\mathcal{C}), \mathcal{O}')
\longrightarrow
(\Sh(\mathcal{B}'), \mathcal{O}_{\mathcal{B}'})
$$
이 $(\Sh(\mathcal{B}'), \mathcal{O}_{\mathcal{B}'})$ 위에서 존재하고
$c$ 및 $f$와 양립할 필요충분조건은 $\zeta$가 $\xi$로 가는 것이다.
이는 정확히 $\mathcal{O}'$이
(\ref{equation-to-solve-ringed-topoi})의 해라는 뜻이다.
\end{proof}






















\section{대수공간의 변형}
\label{section-deformations-spaces}

\noindent
이 절에서는 절 \ref{section-deformations-ringed-topoi}의 결과들이
대수공간의 변형에 대하여 무엇을 뜻하는지 자세히 설명한다.

\begin{lemma}
\label{lemma-match-thickenings}
$S$를 스킴이라 하자. $i : Z \to Z'$를 $S$ 위 대수공간의 사상이라
하자. 다음은 서로 동치이다.
\begin{enumerate}
\item $i$는 공간의 사상의 심화, 절
\ref{spaces-more-morphisms-section-thickenings}에서 정의한 대수공간의
두꺼워짐이다.
\item 연관된 환 달린 토포스의 사상
$i_{small} : (\Sh(Z_\etale), \mathcal{O}_Z) \to
(\Sh(Z'_\etale), \mathcal{O}_{Z'})$
은(공간의 성질, 보조정리
\ref{spaces-properties-lemma-morphism-ringed-topoi})
절 \ref{section-thickenings-ringed-topoi}의 의미에서 두꺼워짐이다.
\end{enumerate}
\end{lemma}

\begin{proof}
이것이 자명한 사실이 아님을 강조한다.

\medskip\noindent
(1)을 가정하자. 공간의 사상의 심화, 보조정리
\ref{spaces-more-morphisms-lemma-thickening-equivalence}
에 의해 사상 $i$는 작은 에탈 사이트들의 동치, 따라서 토포스들의
동치를 유도한다. 물론 두꺼워짐의 정의에 의해 $i^\sharp$는 국소
멱영인 핵을 갖는 전사 준동형이다.

\medskip\noindent
(2)를 가정하자. (이 방향은 덜 중요하고 호기심에 가깝다.) 임의의 에탈
사상 $Y' \to Z'$에 대하여 $Y = Z \times_{Z'} Y'$와 $Y'$은 같은 에탈
토포스를 갖는다. 특히 준콤팩트성은 토포스 이론적 개념이므로
(사이트, 보조정리 \ref{sites-lemma-quasi-compact}) $Y'$이 준콤팩트일
필요충분조건은 $Y$가 준콤팩트인 것이다. 더 나아가 $Y'$이 준콤팩트이고
준분리일 필요충분조건은 $Y$가 준콤팩트이고 준분리인 것이다. 실제로
$Y'$ 위 에탈인 모든 준콤팩트 대수공간 $Y'_1, Y'_2$에 대하여
$Y'_1 \times_{Y'} Y'_2$가 준콤팩트이라는 조건으로 $Y'$의 준분리성을
특징지을 수 있다. $Y'$이 아핀이라고 하자. 그러면 대수공간 $Y$는
준콤팩트이고 준분리이다. 임의의 준연접 $\mathcal{O}_Y$-가군
$\mathcal{F}$에 대하여
$H^q(Y, \mathcal{F}) = H^q(Y', (Y \to Y')_*\mathcal{F})$
이다. 이는 에탈 토포스들이 같기 때문이다. 한편 직접상은 준연접이고
(공간의 사상, 보조정리 \ref{spaces-morphisms-lemma-pushforward})
$Y$는 아핀이므로 $H^q(Y', (Y \to Y')_*\mathcal{F}) = 0$이다.
공간의 코호몰로지, 명제
\ref{spaces-cohomology-proposition-vanishing-affine}
에 의해 $Y'$은 아핀이다(이 보조정리의 이 방향에는 이 인용을 피하는
증명이 틀림없이 있을 것이다). 따라서 $i$는 아핀 사상이다. 아핀인
경우에는 절 \ref{section-thickenings-ringed-topoi}의 조건들로부터
$i$가 대수공간의 두꺼워짐임이 쉽게 따른다.
\end{proof}

\begin{lemma}
\label{lemma-deform-spaces}
$S$를 스킴이라 하자. $Y \subset Y'$을 $S$ 위 대수공간의 1차
두꺼워짐이라 하자. $f : X \to Y$를 $S$ 위 대수공간의 평탄한 사상이라
하자. $S$ 위 대수공간의 평탄한 사상 $f' : X' \to Y'$과 $Y$ 위의
동형사상 $a : X \to X' \times_{Y'} Y$가 존재하면 다음이 성립한다.
\begin{enumerate}
\item 쌍 $(f' : X' \to Y', a)$의 동형류 집합은
$\Ext^1_{\mathcal{O}_X}(\NL_{X/Y}, f^*\mathcal{C}_{Y/Y'})$ 아래의 주동차
공간이다.
\item $X' \times_{Y'} Y$ 위에서 항등사상으로 환원되는 $Y'$ 위의
자기동형사상 $\varphi : X' \to X'$들의 집합은
$\Ext^0_{\mathcal{O}_X}(\NL_{X/Y}, f^*\mathcal{C}_{Y/Y'})$이다.
\end{enumerate}
\end{lemma}

\begin{proof}
환 달린 토포스의 변형에 관한 결과를 보조정리의 대수공간들의 작은 에탈
토포스에 적용하겠다. $X$를 $X'$의 닫힌 부분공간으로 생각할 수 있으며,
그러면 $(f, f') : (X \subset X') \to (Y \subset Y')$는 1차
두꺼워짐들의 사상이다. 보조정리 \ref{lemma-match-thickenings}에 의해
이는 환 달린 토포스의 두꺼워짐들의 사상으로 옮겨진다. 공간의 사상의
심화, 보조정리
\ref{spaces-more-morphisms-lemma-deform}
(또는 더 일반적인 보조정리 \ref{lemma-deform-module-ringed-topoi})에
의해 $X'$에서 $X$의 아이디얼층은 $f^*\mathcal{C}_{Y'/Y}$와 같으며,
이는 실제로 $Y'$ 위에서 $X'$의 평탄성과 동치이다. 따라서 다음 가환
도표가 있다.
$$
\xymatrix{
0 \ar[r] & f^*\mathcal{C}_{Y/Y'} \ar[r] &
\mathcal{O}_{X'} \ar[r] &
\mathcal{O}_X \ar[r] & 0 \\
0 \ar[r] &
f_{small}^{-1}\mathcal{C}_{Y/Y'} \ar[u] \ar[r] &
f_{small}^{-1}\mathcal{O}_{Y'} \ar[u] \ar[r] &
f_{small}^{-1}\mathcal{O}_Y \ar[u] \ar[r] & 0
}
$$
(\ref{equation-to-solve-ringed-topoi})와 비교하라. (2)에서와 같은
자기동형사상 $\varphi$는 위의 도표에 들어맞는 자기동형사상
$\varphi^\sharp : \mathcal{O}_{X'} \to \mathcal{O}_{X'}$를 준다.
거꾸로, 위의 층들의 도표에 들어맞는 자기동형사상
$\alpha : \mathcal{O}_{X'} \to \mathcal{O}_{X'}$
는 공간의 사상의 심화, 보조정리
\ref{spaces-more-morphisms-lemma-first-order-thickening-maps}
에 의해 (2)에서와 같은 어떤 자기동형사상 $\varphi$에 대한
$\varphi^\sharp$와 같다. 끝으로 공간의 사상의 심화, 보조정리
\ref{spaces-more-morphisms-lemma-first-order-thickening}
에 의해, 다음 도표에 들어맞는 $X_\etale$ 위의 또 다른 환층
$\mathcal{A}$를 찾으면
$$
\xymatrix{
0 \ar[r] & f^*\mathcal{C}_{Y/Y'} \ar[r] &
\mathcal{A} \ar[r] &
\mathcal{O}_X \ar[r] & 0 \\
0 \ar[r] &
f_{small}^{-1}\mathcal{C}_{Y/Y'} \ar[u] \ar[r] &
f_{small}^{-1}\mathcal{O}_{Y'} \ar[u] \ar[r] &
f_{small}^{-1}\mathcal{O}_Y \ar[u] \ar[r] & 0
}
$$
1차 두꺼워짐 $X \subset X''$이 존재하고
$\mathcal{O}_{X''} = \mathcal{A}$이다. 공간의 사상의 심화, 보조정리
\ref{spaces-more-morphisms-lemma-first-order-thickening-maps}
를 한 번 더 적용하면, 원하는 성질을 모두 갖는 사상
$(f, f'') : (X \subset X'') \to (Y \subset Y')$를 얻는다. 따라서
(1)은 보조정리 \ref{lemma-choices-ringed-topoi}에서, (2)는 보조정리
\ref{lemma-huge-diagram-ringed-topoi}의 (2)에서 따른다. 공간의 사상의
심화, 절
\ref{spaces-more-morphisms-section-netherlander}
에서 대수공간의 사상에 대하여 정의한 $\NL_{X/Y}$는
절 \ref{section-deformations-ringed-topoi}에서 사용한 $\NL_{X/Y}$와
일치함에 유의하라.
\end{proof}

\noindent
$S$를 스킴이라 하자. $f : X \to B$를 $S$ 위 대수공간의 사상이라 하자.
$\mathcal{F} \to \mathcal{G}$를 $\mathcal{O}_X$-가군의 준동형이라 하자
(준연접일 필요는 없다). 다음 함자를 생각하자.
$$
F :
\left\{
\begin{matrix}
\text{다음 }f^{-1}\mathcal{O}_B\text{-대수의 확대}\\
0 \to \mathcal{F} \to \mathcal{O}' \to \mathcal{O}_X \to 0\\
\mathcal{F}\text{는 제곱이 영인 아이디얼}
\end{matrix}
\right\}
\longrightarrow
\left\{
\begin{matrix}
\text{다음 }f^{-1}\mathcal{O}_B\text{-대수의 확대}\\
0 \to \mathcal{G} \to \mathcal{O}' \to \mathcal{O}_X \to 0\\
\mathcal{G}\text{는 제곱이 영인 아이디얼}
\end{matrix}
\right\}
$$
이 함자는 밀어내기로 주어진다.

\begin{lemma}
\label{lemma-thickening-space-quasi-coherent}
위의 상황에서 $X$가 준콤팩트이고 준분리이며
$DQ_X(\mathcal{F}) \to DQ_X(\mathcal{G})$가 동형사상이라고 가정하자
(공간의 유도 범주, 절
\ref{spaces-perfect-section-better-coherator}). 그러면 함자 $F$는
범주들의 동치이다.
\end{lemma}

\begin{proof}
$\NL_{X/B}$가 $D_\QCoh(\mathcal{O}_X)$의 대상임을 상기하자. 공간의
사상의 심화, 보조정리
\ref{spaces-more-morphisms-lemma-netherlander-quasi-coherent}를 보라.
따라서 우리의 가정으로부터 모든 $i$에 대하여 다음 준동형들이
동형사상임이 따른다.
$$
\Ext^i_X(\NL_{X/B}, \mathcal{F}) \longrightarrow
\Ext^i_X(\NL_{X/B}, \mathcal{G})
$$
보조정리 \ref{lemma-huge-diagram-ringed-topoi}에 의해 이는 함자가
충실충만임을 뜻한다. 한편 두 범주에 각각 해
$\mathcal{O}_X \oplus \mathcal{F}$와
$\mathcal{O}_X \oplus \mathcal{G}$가 있으므로, 보조정리
\ref{lemma-choices-ringed-topoi}에 의해 이 함자는 본질적으로
전사적이다.
\end{proof}

\noindent
$S$를 스킴이라 하자. $B \subset B'$을 $S$ 위 대수공간의 1차
두꺼워짐이라 하고 그 아이디얼층을 $\mathcal{J}$라 하자. 이를 준연접
$\mathcal{O}_B$-가군 또는 $B'$ 위의 준연접 아이디얼층으로 본다.
공간의 사상의 심화, 절
\ref{spaces-more-morphisms-section-thickenings}를 보라.
$f : X \to B$를 $S$ 위 대수공간의 사상이라 하자.
$\mathcal{F} \to \mathcal{G}$를 $\mathcal{O}_X$-가군의 준동형이라
하자(준연접일 필요는 없다).
$c : f^{-1}\mathcal{J} \to \mathcal{F}$를
$f^{-1}\mathcal{O}_B$-가군의 준동형이라 하고, 그 합성을
$c' : f^{-1}\mathcal{J} \to \mathcal{G}$라 쓰자. 다음 함자를 생각하자.
$$
FT :
\{(\ref{equation-to-solve-ringed-topoi})\text{의 }
\mathcal{F}\text{와 }c\text{에 대한 해}\}
\longrightarrow
\{(\ref{equation-to-solve-ringed-topoi})\text{의 }
\mathcal{G}\text{와 }c'\text{에 대한 해}\}
$$
이 함자는 밀어내기로 주어진다.

\begin{lemma}
\label{lemma-thickening-over-thickening-space-quasi-coherent}
위의 상황에서 $X$가 준콤팩트이고 준분리이며
$DQ_X(\mathcal{F}) \to DQ_X(\mathcal{G})$가 동형사상이라고 가정하자
(공간의 유도 범주, 절
\ref{spaces-perfect-section-better-coherator}). 그러면 함자 $FT$는
범주들의 동치이다.
\end{lemma}

\begin{proof}
(\ref{equation-to-solve-ringed-topoi})의 $\mathcal{F}$에 대한 해는 특히
다음 $f^{-1}\mathcal{O}_{B'}$-대수의 확대를 준다.
$$
0 \to \mathcal{F} \to \mathcal{O}' \to \mathcal{O}_X \to 0
$$
여기서 $\mathcal{F}$는 제곱이 영인 아이디얼이다. $\mathcal{G}$에
대해서도 마찬가지이다. 더욱이 그러한 확대가 주어지면 준동형
$c_{\mathcal{O}'} : f^{-1}\mathcal{J} \to \mathcal{F}$를 얻는다.
따라서 $c = c_{\mathcal{O}'}$를 만족하는 그러한
$f^{-1}\mathcal{O}_{B'}$-대수 확대들의 충만 부분범주를 살펴보는
것이다. 이제 $\mathcal{O}'' = F(\mathcal{O}')$이고, 여기서 $F$가
보조정리 \ref{lemma-thickening-space-quasi-coherent}의 동치
(이번에는 $X \to B'$에 적용한다)라면, $c_{\mathcal{O}''}$은
$c_{\mathcal{O}'}$와 준동형 $\mathcal{F} \to \mathcal{G}$의 합성임이
분명하다. 이로써 보조정리가 증명되었다.
\end{proof}





\section{복합체의 변형}
\label{section-deformations-complexes}

\noindent
이 절은 다음 절을 위한 준비이다. 가환대수학 장들의 내용을 가능한 한
많이 사용하겠다.

\begin{lemma}
\label{lemma-canonical-class-algebra}
$R' \to R$을 핵이 제곱이 영인 아이디얼 $I$인 환의 전사 준동형이라
하자. 모든 $K \in D^-(R)$에 대하여 다음 표준 사상이 존재한다.
$$
\omega(K) : K \longrightarrow K \otimes_R^\mathbf{L} I[2]
$$
이는 $D(R)$에서 정의되며 다음 성질들을 갖는다.
\begin{enumerate}
\item $K' \otimes_{R'}^\mathbf{L} R = K$를 만족하는 $K' \in D(R')$가
존재할 필요충분조건은 $\omega(K) = 0$인 것이다.
\item $D^-(R)$에서 $K \to L$이 주어지면 다음 도표는 가환한다.
$$
\xymatrix{
K \ar[d] \ar[rr]_-{\omega(K)} & &
K \otimes^\mathbf{L}_R I[2] \ar[d] \\
L \ar[rr]^-{\omega(L)} & &
L \otimes^\mathbf{L}_R I[2]
}
$$
\item $\omega(K)$의 구성은 환 준동형 $R' \to S'$와 양립한다
(정확한 명제는 증명을 보라).
\end{enumerate}
\end{lemma}

\begin{proof}
$K$를 나타내는 자유 $R$-가군의 위로 유계인 복합체 $K^\bullet$을
택하자. 그러면 $K^n$을 올리는 자유 $R'$-가군 $(K')^n$을 택할 수 있다.
또한 $K^\bullet$의 미분 $d^n_K : K^n \to K^{n + 1}$을 올리는
$R'$-가군 준동형 $(d')^n_K : (K')^n \to (K')^{n + 1}$을 택할 수 있다.
합성
$$
(d')^{n + 1}_K \circ (d')^n_K : (K')^n \to (K')^{n + 2}
$$
은 영이 아닐 수도 있지만 다음과 같이 분해된다.
$$
(K')^n \to K^n \xrightarrow{\omega^n_K}
K^{n + 2} \otimes_R I = I(K')^{n + 2} \to (K')^{n + 2}
$$
이는 $d^{n + 1} \circ d^n = 0$이기 때문이다. 계산하면
$\omega^n_K$가 복합체의 사상을 정함을 알 수 있다. 이 복합체의 사상이
$\omega(K)$를 정한다.

\medskip\noindent
이 구성이 위로 유계인 자유 $R$-가군 복합체의 사상
$\alpha^\bullet : K^\bullet \to L^\bullet$ 및 주어진 올림
$(K')^n, (L')^n, (d')^n_K, (d')^n_L$과 양립함을 증명하자. 구체적으로
성분 $\alpha^n : K^n \to L^n$을 올리는
$(\alpha')^n : (K')^n \to (L')^n$을 택한다. 앞에서와 같이
다음 합성의 분해를 얻는다.
$$
(K')^n \to K^n \xrightarrow{h^n}
L^{n + 1} \otimes_R I = I(L')^{n + 1} \to (L')^{n + 2}
$$
이는 $(d')^n_L \circ (\alpha')^n -
(\alpha')^{n + 1} \circ (d')_K^n$의 분해이다. 간단한 계산으로
다음을 보일 수 있다.
$$
\omega^n_L \circ \alpha^n =
(d_L^{n + 1} \otimes \text{id}_I) \circ h^n + h^{n + 1} \circ d_K^n +
(\alpha^{n + 2} \otimes \text{id}_I) \circ \omega^n_K
$$
이는 이 특수한 경우에 보조정리 (2)의 도표가 가환함을 증명한다.
$K$를 나타내는 위로 유계인 자유 복합체를 서로 다르게 두 번 택하고
이 결과를 적용하면 $\omega(K)$가 잘 정의됨을 알 수 있다. 물론
(2)도 일반적으로 성립한다.

\medskip\noindent
$K$가 $D^-(R')$의 $K'$로 올라가면, $K'$를 자유 $R'$-가군의 위로
유계인 복합체로 나타낼 수 있으므로 $\omega(K) = 0$임을 곧바로 알 수
있다. 거꾸로 우리의 선택 $K^\bullet$, $(K')^n$, $(d')^n_K$로
돌아가자. $\omega(K) = 0$이면 다음을 만족하는 준동형
$g^n : K^n \to K^{n + 1} \otimes_R I$를 찾을 수 있다.
$$
\omega^n = (d_K^{n + 1} \otimes \text{id}_I) \circ g^n +
g^{n + 1} \circ d_K^n
$$
이는 미분을
$(d')^n_K - g^n : (K')^n \to (K')^{n + 1}$
로 택하면 $K^\bullet$을 올리는 자유 $R'$-가군의 복합체를 얻는다는
뜻이다. 이로써 (1)을 증명했다.

\medskip\noindent
끝으로 (3)은 다음을 뜻한다. $R' \to S'$을 환 준동형이라 하자.
$S = S' \otimes_{R'} R$라 놓고, $S' \to S$의 제곱이 영인 핵
$J = IS' \subset S'$라 쓰자. 그러면 $K \in D^-(R)$가 주어졌을 때
다음 가환 도표를 얻는다.
$$
\xymatrix{
K \otimes_R^\mathbf{L} S \ar[d] \ar[rr]_-{\omega(K) \otimes \text{id}} & &
(K \otimes^\mathbf{L}_R I[2]) \otimes_R^\mathbf{L} S \ar[d] \\
K \otimes_R^\mathbf{L} S \ar[rr]^-{\omega(K \otimes_R^\mathbf{L} S)} & &
(K \otimes_R^\mathbf{L} S) \otimes^\mathbf{L}_S J[2]
}
$$
여기서 오른쪽 세로 화살표는 다음 사상에서 온다.
$$
(K \otimes^\mathbf{L}_R I[2]) \otimes_R^\mathbf{L} S =
(K \otimes_R^\mathbf{L} S) \otimes_S^\mathbf{L}
(I \otimes_R^\mathbf{L} S)[2] \longrightarrow
(K \otimes_R^\mathbf{L} S) \otimes_S^\mathbf{L} J[2]
$$
위에서와 같이 $K^\bullet$, $(K')^n$, $(d')^n_K$를 택하자. 그러면
$\omega(K \otimes_R^\mathbf{L} S)$의 구성에
$K^\bullet \otimes_R S$, $(K')^n \otimes_{R'} S'$,
$(d')^n_K \otimes \text{id}_{S'}$를 사용할 수 있다. 이 선택들에
대해서는 복합체의 사상 수준에서 가환성을 곧바로 확인할 수 있다.
\end{proof}







\section{환 달린 토포스 위 복합체의 변형}
\label{section-thickenings-complexes}

\noindent
이 절의 내용은 \cite{lieblich-complexes}에서 가져왔다.

\medskip\noindent
이 절의 내용은 절 \ref{section-thickenings-ringed-topoi}에서 정의한 환
달린 토포스의 1차 두꺼워짐이라는 설정에서 성립한다. 다만 표기를
단순하게 하기 위해 바탕 사이트 $\mathcal{C}$와 $\mathcal{D}$가 같다고
가정하겠다. 또한 문헌의 통상적인 표기에 맞추어 환층의 전사 준동형
$\mathcal{O}' \to \mathcal{O}$를
$\mathcal{O} \to \mathcal{O}_0$로 나타내겠다.

\begin{lemma}
\label{lemma-lift-complex}
$\mathcal{C}$를 사이트라 하자. $\mathcal{O} \to \mathcal{O}_0$를
환층의 전사 준동형이라 하자. 다음 자료가 주어졌다고 가정하자.
\begin{enumerate}
\item 평탄한 $\mathcal{O}$-가군 $\mathcal{G}^n$,
\item $\mathcal{O}$-가군의 준동형 $\mathcal{G}^n \to \mathcal{G}^{n + 1}$,
\item $\mathcal{O}_0$-가군의 복합체 $\mathcal{K}_0^\bullet$,
\item $\mathcal{O}$-가군의 준동형 $\mathcal{G}^n \to \mathcal{K}_0^n$.
\end{enumerate}
이 자료는 다음을 만족한다.
\begin{enumerate}
\item[(a)] $n \gg 0$일 때 $H^n(\mathcal{K}_0^\bullet) = 0$이다.
\item[(b)] $n \gg 0$일 때 $\mathcal{G}^n = 0$이다.
\item[(c)] 다음과 같이 놓자.
$\mathcal{G}^n_0 = \mathcal{G}^n \otimes_\mathcal{O} \mathcal{O}_0$
그러면 유도된 준동형들은 복합체 $\mathcal{G}_0^\bullet$과 복합체의
사상 $\mathcal{G}_0^\bullet \to \mathcal{K}_0^\bullet$를 정한다.
\end{enumerate}
그러면 다음이 존재한다.
\begin{enumerate}
\item[(\romannumeral1)]
평탄한 $\mathcal{O}$-가군 $\mathcal{F}^n$,
\item[(\romannumeral2)]
$\mathcal{O}$-가군의 준동형 $\mathcal{F}^n \to \mathcal{F}^{n + 1}$,
\item[(\romannumeral3)]
$\mathcal{O}$-가군의 준동형 $\mathcal{F}^n \to \mathcal{K}_0^n$,
\item[(\romannumeral4)]
$\mathcal{O}$-가군의 준동형 $\mathcal{G}^n \to \mathcal{F}^n$.
\end{enumerate}
이들은 $n \gg 0$일 때 $\mathcal{F}^n = 0$이고, 모든 $n$에 대하여
다음 도표가 가환하며,
$$
\xymatrix{
\mathcal{G}^n \ar[r] \ar[d] & \mathcal{G}^{n + 1} \ar[d] \\
\mathcal{F}^n \ar[r] & \mathcal{F}^{n + 1}
}
$$
합성 $\mathcal{G}^n \to \mathcal{F}^n \to \mathcal{K}_0^n$은 주어진
준동형 $\mathcal{G}^n \to \mathcal{K}_0^n$이다. 또한
$\mathcal{F}^n_0 = \mathcal{F}^n \otimes_\mathcal{O} \mathcal{O}_0$
라 놓으면 복합체 $\mathcal{F}_0^\bullet$과 준동형사상인 복합체의 사상
$\mathcal{F}_0^\bullet \to \mathcal{K}_0^\bullet$을 얻는다.
\end{lemma}

\begin{proof}
$e$에 관한 내림 귀납법으로, $n \geq e$에 대하여 다음 가환 도표에
들어맞는 $\mathcal{F}^n$, $\mathcal{G}^n \to \mathcal{F}^n$,
$\mathcal{F}^n \to \mathcal{F}^{n + 1}$을 찾을 수 있음을 증명하겠다.
$$
\xymatrix{
\ldots \ar[r] &
\mathcal{G}^{e - 1} \ar[r] \ar@/_2pc/[dd] &
\mathcal{G}^e \ar[d] \ar[r] \ar@/_2pc/[dd] &
\mathcal{G}^{e + 1} \ar[d] \ar[r] \ar@/_2pc/[dd]|\hole &
\ldots \\
& &
\mathcal{F}^e \ar[d] \ar[r] &
\mathcal{F}^{e + 1} \ar[d] \ar[r] & \ldots \\
\ldots  \ar[r] &
\mathcal{K}_0^{e - 1} \ar[r] &
\mathcal{K}_0^e \ar[r] &
\mathcal{K}_0^{e + 1} \ar[r] & \ldots
}
$$
여기서 $\mathcal{F}_0^\bullet$은 복합체이고, 유도된 준동형
$\mathcal{F}_0^\bullet \to \mathcal{K}_0^\bullet$은 $n > e$일 때
$H^n$에서 동형사상을, $n = e$일 때 전사 준동형을 유도한다.
$e \gg 0$이면 가정 (a), (b)에 의해 $n \geq e$에 대하여
$\mathcal{F}^n = 0$으로 택할 수 있으므로 이는 참이다.

\medskip\noindent
귀납 단계. $\mathcal{F}^{e - 1}$과 준동형들
$\mathcal{G}^{e - 1} \to \mathcal{F}^{e - 1}$,
$\mathcal{F}^{e - 1} \to \mathcal{F}^e$,
$\mathcal{F}^{e - 1} \to \mathcal{K}_0^{e - 1}$을 구성해야 한다.
$\mathcal{F}^{e - 1} = A \oplus B \oplus C$를 세 부분의 직합으로
택하겠다.

\medskip\noindent
첫째 부분으로 $A = \mathcal{G}^{e - 1}$를 택하고,
$\mathcal{G}^{e - 1} \to \mathcal{F}^{e - 1}$은 첫 번째 직합 성분의
포함사상으로 택한다. 준동형 $A \to \mathcal{K}^{e - 1}_0$과
$A \to \mathcal{F}^e$은 명백한 것들로 택한다.

\medskip\noindent
$B$를 택하기 위해 귀납 가정으로부터 얻는 다음 전사 준동형을 생각하자.
$$
\gamma :
\Ker(\mathcal{F}^e_0 \to \mathcal{F}^{e + 1}_0)
\longrightarrow
\Ker(\mathcal{K}^e_0 \to \mathcal{K}^{e + 1}_0)/
\Im(\mathcal{K}^{e - 1}_0 \to \mathcal{K}^e_0)
$$
집합 $I$, 각 $i \in I$에 대하여 $\mathcal{C}$의 대상 $U_i$, 그리고
다음 절단들을 택할 수 있다.
$s_i \in \mathcal{F}^e(U_i)$, $t_i \in \mathcal{K}^{e - 1}_0(U_i)$
이들은 다음을 만족한다.
\begin{enumerate}
\item $s_i$는 $\Ker(\gamma) \subset
\Ker(\mathcal{F}^e_0 \to \mathcal{F}^{e + 1}_0)$의 절단으로 간다.
\item $s_i$와 $t_i$는 $\mathcal{K}^e_0$의 같은 절단으로 간다.
\item 절단들 $s_i$는 $\mathcal{O}_0$-가군으로서 $\Ker(\gamma)$를
생성한다.
\end{enumerate}
이에 대한 자세한 설명은 생략한다. 여기서는
$\mathcal{F}^e \to \mathcal{F}^e_0$가 집합층의 전사 사상이라는 사실을
사용한다. 이제 다음과 같이 놓는다.
$$
B = \bigoplus\nolimits_{i \in I} j_{U_i!}\mathcal{O}_{U_i}
$$
그리고 $s_i$와 $t_i$를 사용하여 직합 성분
$j_{U_i!}\mathcal{O}_{U_i}$의 상을 정함으로써 준동형
$B \to \mathcal{F}^e$과 $B \to \mathcal{K}_0^{e - 1}$을 정의한다.

\medskip\noindent
$\mathcal{F}^{e - 1} = A \oplus B$와 위의 준동형들을 사용하면,
$\mathcal{F}_0^\bullet \to \mathcal{K}_0^\bullet$이 $n \geq e$일 때
$H^n$에서 동형사상을 유도하는 $e - 1$에 대한 위와 같은 도표를 얻는다.
$H^{e - 1}$에서 이 준동형이 전사이도록 하기 위해 직합 성분 $C$를
다음과 같이 택한다. 집합 $J$, 각 $j \in J$에 대한 $\mathcal{C}$의
대상 $U_j$, 그리고 $U_j$ 위의
$\Ker(\mathcal{K}^{e - 1}_0 \to \mathcal{K}^e_0)$의 절단 $t_j$를
택하되, 이 절단들이 $\mathcal{O}_0$ 위에서 이 핵을 생성하도록 한다.
이제 다음과 같이 놓는다.
$$
C = \bigoplus\nolimits_{j \in J} j_{U_j!}\mathcal{O}_{U_j}
$$
그리고 $C \to \mathcal{F}^e$는 영 준동형으로, 준동형
$C \to \mathcal{K}_0^{e - 1}$은 $s_j$를 사용하여 직합 성분
$j_{U_j!}\mathcal{O}_{U_j}$의 상을 정함으로써 정의한다.
$\mathcal{F}^{e - 1} = A \oplus B \oplus C$와 위에서 지정한 준동형들을
택하면 귀납 단계가 끝난다.
\end{proof}

\begin{lemma}
\label{lemma-canonical-class}
$\mathcal{C}$를 사이트라 하자. $\mathcal{O} \to \mathcal{O}_0$를 핵이
제곱이 영인 아이디얼층 $\mathcal{I}$인 환층의 전사 준동형이라 하자.
$D^-(\mathcal{O}_0)$의 모든 대상 $K_0$에 대하여 다음 표준 사상이
존재한다.
$$
\omega(K_0) :
K_0 \longrightarrow
K_0 \otimes_{\mathcal{O}_0}^\mathbf{L} \mathcal{I}[2]
$$
이는 $D(\mathcal{O}_0)$에서 정의되며,
$D^-(\mathcal{O}_0)$의 임의의 사상 $K_0 \to L_0$에 대하여 다음 도표가
가환한다.
$$
\xymatrix{
K_0 \ar[d] \ar[rr]_-{\omega(K_0)} & &
(K_0 \otimes^\mathbf{L}_{\mathcal{O}_0} \mathcal{I})[2] \ar[d] \\
L_0 \ar[rr]^-{\omega(L_0)} & &
(L_0 \otimes^\mathbf{L}_{\mathcal{O}_0} \mathcal{I})[2]
}
$$
\end{lemma}

\begin{proof}
$K_0$를 $\mathcal{O}_0$-가군의 임의의 복합체
$\mathcal{K}_0^\bullet$로 나타내자. 모든 $n$에 대하여
$\mathcal{G}^n = 0$으로 놓고 보조정리 \ref{lemma-lift-complex}를
적용한다. 보조정리가 만드는 준동형을
$d : \mathcal{F}^n \to \mathcal{F}^{n + 1}$라 쓰자. 그러면
$d \circ d : \mathcal{F}^n \to \mathcal{F}^{n + 2}$
는 $\mathcal{I}$을 법으로 영이다. $\mathcal{F}^n$이 평탄하므로
$\mathcal{I}\mathcal{F}^n =
\mathcal{F}^n \otimes_{\mathcal{O}} \mathcal{I} =
\mathcal{F}^n_0 \otimes_{\mathcal{O}_0} \mathcal{I}$이다. 따라서
복합체의 표준 사상
$$
d \circ d : \mathcal{F}_0^\bullet
\longrightarrow
(\mathcal{F}_0^\bullet \otimes_{\mathcal{O}_0} \mathcal{I})[2]
$$
을 얻는다. $\mathcal{F}_0^\bullet$은 평탄한 $\mathcal{O}_0$-가군의
위로 유계인 복합체이므로 K-평탄하고, 유도 텐서곱을 계산하는 데 사용할
수 있다. 또한 구성상 복합체의 사상
$\mathcal{F}_0^\bullet \to \mathcal{K}_0^\bullet$은 준동형사상이다.
따라서 방금 구성한 사상의 정의역과 공역은 각각 $K_0$와
$K_0 \otimes_{\mathcal{O}_0}^\mathbf{L} \mathcal{I}[2]$를 나타내며,
원하는 사상 $\omega(K_0)$를 얻는다.

\medskip\noindent
이 절차가 복합체의 사상과 양립함을 보이자.
$\mathcal{L}_0^\bullet$이 $D^-(\mathcal{O}_0)$의 다른 대상을 나타내고,
다음이 복합체의 사상이라고 가정하자.
$$
\mathcal{K}_0^\bullet \longrightarrow \mathcal{L}_0^\bullet
$$
복합체 $\mathcal{L}_0^\bullet$, 평탄한 가군 $\mathcal{F}^n$, 준동형
$\mathcal{F}^n \to \mathcal{F}^{n + 1}$, 그리고 합성
$\mathcal{F}^n \to \mathcal{K}_0^n \to \mathcal{L}_0^n$
에 보조정리 \ref{lemma-lift-complex}를 적용한다(사용한 글자의 순서가
뒤집힌 점을 양해해 주기 바란다). 그러면 보조정리의 모든 성질을 갖는
평탄한 가군 $\mathcal{G}^n$, 준동형
$\mathcal{F}^n \to \mathcal{G}^n$,
$\mathcal{G}^n \to \mathcal{G}^{n + 1}$,
$\mathcal{G}^n \to \mathcal{L}_0^n$을 얻는다. 그러면 다음이 복합체의
가환 도표임은 분명하다.
$$
\xymatrix{
\mathcal{F}_0^\bullet \ar[d] \ar[r] &
(\mathcal{F}_0^\bullet \otimes_{\mathcal{O}_0} \mathcal{I})[2] \ar[d] \\
\mathcal{G}_0^\bullet \ar[r] &
(\mathcal{G}_0^\bullet \otimes_{\mathcal{O}_0} \mathcal{I})[2]
}
$$

\medskip\noindent
이제 $\omega(K_0)$가 잘 정의됨을 보이기 위해, $K_0$를 나타내는 두
$\mathcal{O}_0$-가군 복합체
$\mathcal{K}_0^\bullet$과 $(\mathcal{K}'_0)^\bullet$
와 다음 두 체계가 주어졌다고 가정하자.
$(\mathcal{F}^n, d : \mathcal{F}^n \to \mathcal{F}^{n + 1},
\mathcal{F}^n \to \mathcal{K}_0^n)$
그리고
$((\mathcal{F}')^n, d : (\mathcal{F}')^n \to (\mathcal{F}')^{n + 1},
(\mathcal{F}')^n \to \mathcal{K}_0^n)$
이다. 그러면 복합체 $(\mathcal{K}''_0)^\bullet$과 준동형사상들
$\mathcal{K}_0^\bullet  \to (\mathcal{K}''_0)^\bullet$
및
$(\mathcal{K}'_0)^\bullet  \to (\mathcal{K}''_0)^\bullet$
을 택할 수 있다. 이들은 두 복합체가 유도 범주에서 모두 $K_0$를
나타낸다는 사실을 실현한다. 다음으로 앞 문단의 결과를
$$
(\mathcal{K}_0)^\bullet \oplus (\mathcal{K}'_0)^\bullet
\longrightarrow
(\mathcal{K}''_0)^\bullet
$$
에 적용한다. 그러면 다음 가환 도표를 얻는다.
$$
\xymatrix{
\mathcal{F}_0^\bullet \oplus (\mathcal{F}'_0)^\bullet
\ar[d] \ar[r] &
(\mathcal{F}_0^\bullet \otimes_{\mathcal{O}_0} \mathcal{I})[2] \oplus
((\mathcal{F}'_0)^\bullet \otimes_{\mathcal{O}_0} \mathcal{I})[2] \ar[d] \\
\mathcal{G}_0^\bullet \ar[r] &
(\mathcal{G}_0^\bullet \otimes_{\mathcal{O}_0} \mathcal{I})[2]
}
$$
세로 화살표들은 각 직합 성분에서 준동형사상을 주므로
$D(\mathcal{O}_0)$에서 원하는 가환성을 얻는다.

\medskip\noindent
잘 정의됨을 보였으므로 사상과의 양립성에 관한 명제는 두 번째 문단의
결과에서 따른다.
\end{proof}

\begin{lemma}
\label{lemma-induced-map}
$(\mathcal{C}, \mathcal{O})$를 환 달린 사이트라 하자.
$\alpha : K \to L$을 $D^-(\mathcal{O})$의 사상이라 하자.
$\mathcal{F}$를 $\mathcal{O}$-가군층이라 하고
$n \in \mathbf{Z}$라 하자.
\begin{enumerate}
\item $i \geq n$일 때 $H^i(\alpha)$가 동형사상이면, $i \geq n$일 때
$H^i(\alpha \otimes_\mathcal{O}^\mathbf{L}
\text{id}_\mathcal{F})$도 동형사상이다.
\item $i > n$일 때 $H^i(\alpha)$가 동형사상이고 $i = n$일 때
전사 준동형이면, $i > n$일 때
$H^i(\alpha \otimes_\mathcal{O}^\mathbf{L}
\text{id}_\mathcal{F})$가 동형사상이고 $i = n$일 때 전사 준동형이다.
\end{enumerate}
\end{lemma}

\begin{proof}
완전 삼각형
$$
K \to L \to C \to K[1]
$$
을 택하자. (2)의 경우 $i \geq n$일 때 $H^i(C) = 0$이다. 따라서
유도 범주, 보조정리 \ref{derived-lemma-negative-vanishing}의 쌍대에
의해 $i \geq n$일 때
$H^i(C \otimes_\mathcal{O}^\mathbf{L} \mathcal{F}) = 0$이다. 이는 다시
$i > n$일 때
$H^i(\alpha \otimes_\mathcal{O}^\mathbf{L} \text{id}_\mathcal{F})$
가 동형사상이고 $i = n$일 때 전사 준동형임을 보인다. (1)의 경우에는
더 나아가 $H^{n - 1}(L) \to H^{n - 1}(C)$가 전사 준동형이다.
다음 도표를 생각하면
$$
\xymatrix{
H^{n - 1}(L) \otimes_\mathcal{O} \mathcal{F} \ar[r] \ar[d] &
H^{n - 1}(C) \otimes_\mathcal{O} \mathcal{F} \ar@{=}[d] \\
H^{n - 1}(L \otimes_\mathcal{O}^\mathbf{L} \mathcal{F}) \ar[r] &
H^{n - 1}(C \otimes_\mathcal{O}^\mathbf{L} \mathcal{F})
}
$$
아래쪽 가로 화살표가 전사 준동형임을 알 수 있다. 이를 앞에서 말한 것과
합치면
$H^n(\alpha \otimes_\mathcal{O}^\mathbf{L} \text{id}_\mathcal{F})$
가 동형사상임을 얻는다.
\end{proof}

\begin{lemma}
\label{lemma-canonical-class-obstruction}
$\mathcal{C}$를 사이트라 하자. $\mathcal{O} \to \mathcal{O}_0$를 핵이
제곱이 영인 아이디얼층 $\mathcal{I}$인 환층의 전사 준동형이라 하자.
$D^-(\mathcal{O}_0)$의 모든 대상 $K_0$에 대하여 다음은 서로 동치이다.
\begin{enumerate}
\item 보조정리 \ref{lemma-canonical-class}에서 구성한 동치류
$\omega(K_0) \in
\Ext^2_{\mathcal{O}_0}(K_0, K_0 \otimes_{\mathcal{O}_0} \mathcal{I})$
는 영이다.
\item $D(\mathcal{O}_0)$에서
$K \otimes_\mathcal{O}^\mathbf{L} \mathcal{O}_0 = K_0$
를 만족하는 $K \in D^-(\mathcal{O})$가 존재한다.
\end{enumerate}
\end{lemma}

\begin{proof}
$K$를 (2)에서와 같이 택하자. 그러면 $K$를 평탄한
$\mathcal{O}$-가군의 위로 유계인 복합체 $\mathcal{F}^\bullet$로
나타낼 수 있다. 이때
$\mathcal{F}_0^\bullet =
\mathcal{F}^\bullet \otimes_{\mathcal{O}} \mathcal{O}_0$
은 $D(\mathcal{O}_0)$에서 $K_0$를 나타낸다.
$\mathcal{F}^\bullet$은 복합체이므로
$d_{\mathcal{F}^\bullet} \circ d_{\mathcal{F}^\bullet} = 0$이다.
따라서 $\omega(K_0)$의 구성 자체에서 이것이 영임을 알 수 있다.

\medskip\noindent
(1)을 가정하자. $\mathcal{F}^n$과
$d : \mathcal{F}^n \to \mathcal{F}^{n + 1}$을 $\omega(K_0)$의
구성에서와 같이 택하자. $\omega(K_0)$가 영이라는 사실은 사상
$$
\omega = d \circ d : \mathcal{F}_0^\bullet
\longrightarrow
(\mathcal{F}_0^\bullet \otimes_{\mathcal{O}_0} \mathcal{I})[2]
$$
이 $D(\mathcal{O}_0)$에서 영임을 뜻한다. 유도 범주를
$\mathcal{O}_0$-가군 복합체의 호모토피 범주의 국소화로 정의하는
것에 의해, 준동형사상
$\alpha : \mathcal{G}_0^\bullet \to \mathcal{F}_0^\bullet$
과 다음을 만족하는 $\mathcal{O}_0$-가군 준동형
$h^n : \mathcal{G}_0^n \to
\mathcal{F}_0^{n + 1} \otimes_\mathcal{O} \mathcal{I}$
가 존재한다.
$$
\omega \circ \alpha =
d_{\mathcal{F}_0^\bullet \otimes \mathcal{I}} \circ h +
h \circ d_{\mathcal{G}_0^\bullet}
$$
다음과 같이 놓는다.
$$
\mathcal{H}^n = \mathcal{F}^n \times_{\mathcal{F}^n_0} \mathcal{G}_0^n
$$
그리고 다음을 정의한다.
$$
d' : \mathcal{H}^n \longrightarrow \mathcal{H}^{n + 1},\quad
(f^n, g_0^n) \longmapsto (d(f^n) - h^n(g_0^n), d(g_0^n))
$$
여기서는 다음 등식을 사용한 명백한 표기를 썼다.
$\mathcal{F}_0^{n + 1} \otimes_{\mathcal{O}_0} \mathcal{I} =
\mathcal{F}^{n + 1} \otimes_\mathcal{O} \mathcal{I} =
\mathcal{I}\mathcal{F}^{n + 1} \subset \mathcal{F}^{n + 1}$.
$h^n$의 선택과 $\omega$의 정의를 사용하면
$d' \circ d' = 0$임을 확인할 수 있다. 따라서
$\mathcal{H}^\bullet$은 $D(\mathcal{O})$의 대상을 정한다. 한편
$\mathcal{O}$-가군 복합체의 다음 짧은 완전열이 있다.
$$
0 \to \mathcal{F}_0^\bullet \otimes_{\mathcal{O}_0} \mathcal{I} \to
\mathcal{H}^\bullet \to \mathcal{G}_0^\bullet \to 0
$$
이제
$\mathcal{H}^\bullet \otimes_\mathcal{O}^\mathbf{L} \mathcal{O}_0$
가 $K_0$와 동형임을 보이면 된다. 준동형사상
$\mathcal{E}^\bullet \to \mathcal{H}^\bullet$
을 택하되, $\mathcal{E}^\bullet$은 평탄한 $\mathcal{O}$-가군의 위로
유계인 복합체로 택한다. 그러면 다음 가환 도표를 얻는다.
$$
\xymatrix{
0 \ar[r] &
\mathcal{E}^\bullet \otimes_\mathcal{O} \mathcal{I} \ar[d]^\beta \ar[r] &
\mathcal{E}^\bullet \ar[d]^\gamma \ar[r] &
\mathcal{E}_0^\bullet \ar[d]^\delta \ar[r] &
0 \\
0 \ar[r] &
\mathcal{F}_0^\bullet \otimes_{\mathcal{O}_0} \mathcal{I} \ar[r] &
\mathcal{H}^\bullet \ar[r] &
\mathcal{G}_0^\bullet \ar[r] &
0
}
$$
$\delta$가 준동형사상이라고 주장한다. $i \gg 0$일 때
$H^i(\delta)$가 동형사상이므로, $i \geq n$일 때 $H^i(\delta)$가
동형사상이라는 $n$에 관한 내림 귀납법을 사용할 수 있다. 다음에
유의하자.
$\mathcal{E}^\bullet \otimes_\mathcal{O} \mathcal{I}$
은
$\mathcal{E}_0^\bullet \otimes_{\mathcal{O}_0}^\mathbf{L} \mathcal{I}$,
를 나타내고,
$\mathcal{F}_0^\bullet \otimes_{\mathcal{O}_0} \mathcal{I}$
은
$\mathcal{G}_0^\bullet \otimes_{\mathcal{O}_0}^\mathbf{L} \mathcal{I}$,
를 나타내며,
$\beta = \delta \otimes_{\mathcal{O}_0}^\mathbf{L} \text{id}_\mathcal{I}$
는 $D(\mathcal{O}_0)$에서의 사상으로서 성립한다. 마지막 주장은
다음 등식에서 따른다.
$\beta =
(\alpha \otimes \text{id}_\mathcal{I})
\circ
(\delta \otimes \text{id}_\mathcal{I})$.
$H^i(\delta)$가 차수 $\geq n$에서 동형사상이라고 가정하자. 방금 말한
것과 보조정리 \ref{lemma-induced-map}에 의해 $\beta$에 대해서도 같은
것이 참이다. 이제 다음 도표를 보자.
$$
\xymatrix@C=1.2em{
H^{n - 1}(\mathcal{E}^\bullet \otimes_\mathcal{O} \mathcal{I})
\ar[r] \ar[d]^{H^{n - 1}(\beta)} &
H^{n - 1}(\mathcal{E}^\bullet) \ar[r] \ar[d] &
H^{n - 1}(\mathcal{E}_0^\bullet) \ar[r] \ar[d]^{H^{n - 1}(\delta)} &
H^n(\mathcal{E}^\bullet \otimes_\mathcal{O} \mathcal{I})
\ar[r] \ar[d]^{H^n(\beta)} &
H^n(\mathcal{E}^\bullet) \ar[d] \\
H^{n - 1}(\mathcal{F}_0^\bullet \otimes_\mathcal{O} \mathcal{I}) \ar[r] &
H^{n - 1}(\mathcal{H}^\bullet) \ar[r] &
H^{n - 1}(\mathcal{G}_0^\bullet) \ar[r] &
H^n(\mathcal{F}_0^\bullet \otimes_\mathcal{O} \mathcal{I}) \ar[r] &
H^n(\mathcal{H}^\bullet)
}
$$
호몰로지, 보조정리 \ref{homology-lemma-four-lemma}를 사용하면
$H^{n - 1}(\delta)$가 전사 준동형임을 알 수 있다. 이는 다시 보조정리
\ref{lemma-induced-map}에 의해 $H^{n - 1}(\beta)$가 전사 준동형임을
뜻한다. 호몰로지, 보조정리 \ref{homology-lemma-four-lemma}를 다시
사용하면 $H^{n - 1}(\delta)$가 동형사상임을 알 수 있다. 따라서
귀납법으로 주장이 성립하며, $\delta$는 준동형사상이다. 이것이
보이고자 한 것이다.
\end{proof}

\begin{lemma}
\label{lemma-lift-map-complexes}
$\mathcal{C}$를 사이트라 하자. $\mathcal{O} \to \mathcal{O}_0$를
환층의 전사 준동형이라 하자. 다음 자료가 주어졌다고 가정하자.
\begin{enumerate}
\item $\mathcal{O}$-가군의 복합체 $\mathcal{F}^\bullet$,
\item $\mathcal{O}_0$-가군의 복합체 $\mathcal{K}_0^\bullet$,
\item 준동형사상 $\mathcal{K}_0^\bullet \to
\mathcal{F}^\bullet \otimes_\mathcal{O} \mathcal{O}_0$.
\end{enumerate}
그러면 준동형사상 $\mathcal{G}^\bullet \to \mathcal{F}^\bullet$이
존재하여, 복합체의 사상
$\mathcal{G}^\bullet  \otimes_\mathcal{O} \mathcal{O}_0 \to
\mathcal{F}^\bullet \otimes_\mathcal{O} \mathcal{O}_0$
은 $\mathcal{O}_0$-가군 복합체의 호모토피 범주에서
$\mathcal{K}_0^\bullet$을 거쳐 분해된다.
\end{lemma}

\begin{proof}
다음과 같이 놓자.
$\mathcal{F}_0^\bullet =
\mathcal{F}^\bullet \otimes_\mathcal{O} \mathcal{O}_0$.
유도 범주, 보조정리 \ref{derived-lemma-make-surjective}에 의해 주어진
사상의 다음 분해가 존재한다.
$$
\mathcal{K}_0^\bullet \to \mathcal{L}_0^\bullet \to \mathcal{F}_0^\bullet
$$
여기서 첫 번째 화살표는 호모토피까지 역을 가지며 두 번째 화살표는
각 항에서 분할 전사 준동형이다. 따라서
$\mathcal{K}_0^\bullet \to \mathcal{F}_0^\bullet$이 각 항에서
전사라고 가정해도 된다. 이 경우 다음과 같이 택한다.
$$
\mathcal{G}^n = \mathcal{F}^n \times_{\mathcal{F}^n_0} \mathcal{K}_0^n
$$
그러면 모든 것이 분명하다.
\end{proof}

\begin{lemma}
\label{lemma-inf-obs-map-defo-complex}
$\mathcal{C}$를 사이트라 하자. $\mathcal{O} \to \mathcal{O}_0$를
커널이 제곱 영 아이디얼 층 $\mathcal{I}$인 환층의 전사 준동형이라 하자.
$K, L \in D^-(\mathcal{O})$라 하자. $D^-(\mathcal{O}_0)$에서
$K_0 = K \otimes_\mathcal{O}^\mathbf{L} \mathcal{O}_0$ 및
$L_0 = L \otimes_\mathcal{O}^\mathbf{L} \mathcal{O}_0$로 놓자.
$D(\mathcal{O}_0)$의 사상 $\alpha_0 : K_0 \to L_0$가 주어지면
정준 원소
$$
o(\alpha_0) \in \Ext^1_{\mathcal{O}_0}(K_0,
L_0 \otimes_{\mathcal{O}_0}^\mathbf{L} \mathcal{I})
$$
가 존재한다. 이 원소가 소멸하는 것은
$\alpha_0 = \alpha \otimes_\mathcal{O}^\mathbf{L} \text{id}$를 만족하는
$D(\mathcal{O})$의 사상 $\alpha : K \to L$가 존재하기 위한
필요충분조건이다.
\end{lemma}

\begin{proof}
$\alpha_0$를 올리는 $\alpha : K \to L$를 찾는 것은
합성 $K \xrightarrow{\alpha} L \to L_0$가
합성 $K \to K_0 \xrightarrow{\alpha_0} L_0$와 같도록 하는
$\alpha : K \to L$를 찾는 것과 같다. 짧은 완전열
$0 \to \mathcal{I} \to \mathcal{O} \to \mathcal{O}_0 \to 0$
은 $D(\mathcal{O})$에서 정준 완전 삼각형
$$
L \otimes_\mathcal{O}^\mathbf{L} \mathcal{I} \to
L \to
L_0 \to
(L \otimes_\mathcal{O}^\mathbf{L} \mathcal{I})[1]
$$
을 준다. 유도 범주, 보조정리
\ref{derived-lemma-representable-homological}에 의해 합성
$$
K \to K_0 \xrightarrow{\alpha_0} L_0 \to
(L \otimes_\mathcal{O}^\mathbf{L} \mathcal{I})[1]
$$
이 영인 것은 $\alpha_0$를 올리는 $\alpha : K \to L$를 찾을 수
있기 위한 필요충분조건이다. 이 합성은 수반성에 의해
$$
\Hom_{D(\mathcal{O})}(K, (L \otimes_\mathcal{O}^\mathbf{L} \mathcal{I})[1]) =
\Hom_{D(\mathcal{O}_0)}(K_0,
(L \otimes_\mathcal{O}^\mathbf{L} \mathcal{I})[1]) =
\Ext^1_{\mathcal{O}_0}(K_0,
L_0 \otimes_{\mathcal{O}_0}^\mathbf{L} \mathcal{I})
$$
의 원소이다.
\end{proof}

\begin{lemma}
\label{lemma-first-order-defos-complex}
$\mathcal{C}$를 사이트라 하자. $\mathcal{O} \to \mathcal{O}_0$를
커널이 제곱 영 아이디얼 층 $\mathcal{I}$인 환층의 전사 준동형이라 하자.
$K_0 \in D^-(\mathcal{O})$라 하자. $K_0$의 올림이란
$D^-(\mathcal{O})$의 대상 $K$와 동형사상
$\alpha_0 : K \otimes_\mathcal{O}^\mathbf{L} \mathcal{O}_0 \to K_0$
$D(\mathcal{O}_0)$로 이루어진 쌍 $(K, \alpha_0)$을 말한다.
\begin{enumerate}
\item 올림 $(K, \alpha)$가 주어지면 이 쌍의 자기동형군은 정준적으로
다음 사상의 여핵이다.
$$
\Ext^{-1}_{\mathcal{O}_0}(K_0, K_0)
\longrightarrow
\Hom_{\mathcal{O}_0}(K_0, K_0 \otimes_{\mathcal{O}_0}^\mathbf{L} \mathcal{I})
$$
\item 올림이 존재하면, 올림의 동형류 집합은
$\Ext^1_{\mathcal{O}_0}(K_0,
K_0 \otimes_{\mathcal{O}_0}^\mathbf{L} \mathcal{I})$의 작용에 대한
주동차 공간을 이룬다.
\end{enumerate}
\end{lemma}

\begin{proof}
$(K, \alpha)$의 자기동형은
$\varphi \otimes_\mathcal{O} \text{id}_{\mathcal{O}_0} = \text{id}$.
를 만족하는 $D(\mathcal{O})$의 사상 $\varphi : K \to K$이다.
이는
$$
K \xrightarrow{\varphi - \text{id}} K \to
K \otimes_\mathcal{O}^\mathbf{L} \mathcal{O}_0
$$
가 영이라는 것과 같다. 따라서 자기동형군은 다음 사상의 여핵이다.
$$
\Hom_\mathcal{O}(K, K_0[-1])
\longrightarrow
\Hom_\mathcal{O}(K, K_0 \otimes_{\mathcal{O}_0}^\mathbf{L} \mathcal{I})
$$
실제로 이는 $D(\mathcal{O})$의 완전 삼각형
$$
K \otimes_\mathcal{O}^\mathbf{L} \mathcal{I} \to
K \to
K \otimes_\mathcal{O}^\mathbf{L} \mathcal{O}_0 \to
(K \otimes_\mathcal{O}^\mathbf{L} \mathcal{I})[1]
$$
과 유도 범주, 보조정리 \ref{derived-lemma-representable-homological}에서
따른다. 이를 보조정리에 나타난 군으로 옮기려면 제한 함자
$D(\mathcal{O}_0) \to D(\mathcal{O})$와
$- \otimes_\mathcal{O} \mathcal{O}_0 : D(\mathcal{O}) \to D(\mathcal{O}_0)$
사이의 수반성을 사용한다. 이로써 (1)이 증명되었다.

\medskip\noindent
(2)의 증명. $D(\mathcal{O})$에서
$K_0 = K \otimes_\mathcal{O}^\mathbf{L} \mathcal{O}_0$라고 가정하자.
보조정리 \ref{lemma-inf-obs-map-defo-complex}에 의해 올림
$(K', \alpha_0)$을 $\alpha_0$를 올리는 데 대한 장애 $o(\alpha_0)$로
보내는 사상은 쌍들의 동형류 집합에서
$\Ext^1_{\mathcal{O}_0}(K_0,
K_0 \otimes_{\mathcal{O}_0}^\mathbf{L} \mathcal{I})$로 가는
정준 단사 사상을 정의한다. 증명을 마치기 위해 이 사상이 전사임을
보이자. 보조정리에 나오는 $\Ext^1$의 원소
$\xi : K_0 \to (K_0 \otimes_{\mathcal{O}_0}^\mathbf{L} \mathcal{I})[1]$를
택하자. $K$를 나타내는 평탄 $\mathcal{O}$-가군의 위로 유계인 복합체
$\mathcal{F}^\bullet$을 택하자. 사상 $\xi$는
$t \circ s^{-1}$로 나타낼 수 있다. 여기서
$s : \mathcal{K}_0^\bullet \to \mathcal{F}_0^\bullet$은
준동형사상이고
$t : \mathcal{K}_0^\bullet \to
\mathcal{F}_0^\bullet \otimes_{\mathcal{O}_0} \mathcal{I}[1]$
는 복합체의 사상이다. 보조정리 \ref{lemma-lift-map-complexes}에 의해
$\mathcal{O}$-가군 복합체의 준동형사상
$\mathcal{G}^\bullet \to \mathcal{F}^\bullet$이 존재하고,
$\mathcal{G}_0^\bullet \to \mathcal{F}_0^\bullet$이 호모토피까지
$s$를 거쳐 분해된다고 가정할 수 있다. 또한 평탄 $\mathcal{O}$-가군의
위로 유계인 복합체에서 $\mathcal{G}^\bullet$로 가는 준동형사상을
택하여 대체함으로써, $\mathcal{G}^\bullet$이 그런 복합체라고
가정해도 된다. 그러면 $\xi$는 복합체의 사상
$t : \mathcal{G}_0^\bullet \to
\mathcal{F}_0^\bullet \otimes_{\mathcal{O}_0} \mathcal{I}[1]$
과 준동형사상 $\mathcal{G}_0^\bullet \to \mathcal{F}_0^\bullet$로
나타남을 알 수 있다. 다음과 같이 놓자.
$$
\mathcal{H}^n = \mathcal{F}^n \times_{\mathcal{F}_0^n} \mathcal{G}_0^n
$$
그 미분은 다음과 같다.
$$
\mathcal{H}^n \to \mathcal{H}^{n + 1},\quad
(f^n, g_0^n) \mapsto (d(f^n) + t(g_0^n), d(g_0^n))
$$
이는
$\mathcal{F}_0^{n + 1} \otimes_{\mathcal{O}_0} \mathcal{I} =
\mathcal{F}^{n + 1} \otimes_\mathcal{O} \mathcal{I} =
\mathcal{I}\mathcal{F}^{n + 1} \subset \mathcal{F}^{n + 1}$이므로
잘 정의된다. $\mathcal{H}^\bullet$이 $\mathcal{O}$-가군의 복합체임을
보이는 계산은 생략한다. 구성에 의해 $\mathcal{O}$-가군 복합체의
짧은 완전열
$$
0 \to \mathcal{F}_0^\bullet \otimes_{\mathcal{O}_0} \mathcal{I} \to
\mathcal{H}^\bullet \to \mathcal{G}_0^\bullet \to 0
$$
이 존재한다. 보조정리 \ref{lemma-canonical-class-obstruction}의 증명과
똑같이, 이 열이 동형사상
$\alpha_0 :
\mathcal{H}^\bullet \otimes_\mathcal{O}^\mathbf{L} \mathcal{O}_0 \to
\mathcal{G}_0^\bullet$를 $D(\mathcal{O}_0)$에서 유도함을 보일 수 있다.
다시 말해 쌍 $(\mathcal{H}^\bullet, \alpha_0)$을 구성했다.
$o(\alpha_0) = \xi$의 확인은 생략한다. 힌트는 다음과 같다. 이 경우
$\alpha_0$의 각 성분을 올리는 사상
$\mathcal{H}^n \to \mathcal{F}^n$이 있으므로(이들은 미분과 양립하지
않는다), $o(\alpha_0)$를 명시적으로 계산할 수 있다. 이로써 증명이
끝난다.
\end{proof}












\input{ko_chapters}

\bibliography{my}
\bibliographystyle{amsalpha}

\end{document}
